% Gestion de la santé et du bien-être social — Informatique 3ème — Cours signé OnBuch+
% Programme officiel MINESEC. Compiler avec : tectonic N12-gestion-sante-bien-etre-social.tex
\newcommand{\DOCMATIERE}{Informatique}
\newcommand{\DOCNIVEAU}{3ème}
\newcommand{\DOCMODULE}{Informatique – classe de 3ème}
\newcommand{\DOCLECON}{N12}
\newcommand{\DOCTITRE}{Gestion de la santé et du bien-être social}
% OnBuch+ — préambule « cours signés » ENRICHI (compilé avec tectonic / XeLaTeX)
% Système visuel « Pop » d'OnBuch+ : fond crème (jamais blanc), bordures encre
% épaisses, ombres « dures » décalées, titres Archivo Black en capitales.
% Version MATHÉMATIQUES : graphes (pgfplots/tikz), boîtes pédagogiques
% (définition, propriété, méthode, attention, le savais-tu, exercice résolu,
% exercice, TP, bilan), page de garde et signature OnBuch+ ancrée en bas.
%
% Macros attendues dans main.tex AVANT \input{preamble} :
%   \DOCMATIERE  \DOCNIVEAU  \DOCMODULE  \DOCLECON (numéro)  \DOCTITRE
\documentclass[11pt]{article}

\usepackage{fontspec}
\usepackage[french]{babel}
\usepackage[a4paper,margin=2cm,top=3.4cm,bottom=2.8cm,headheight=1.4cm,footskip=1.3cm]{geometry}
\usepackage{amsmath,amssymb,amsfonts}
\usepackage{mathtools} % \xmapsto, \coloneqq... souvent utilisés par les modèles
\usepackage{cancel} % simplifications barrées (bilans de forces)
% \abs{z} (module, valeur absolue) et \norm{u} (norme), utilisés par les modèles
\providecommand{\abs}{}\let\abs\relax\DeclarePairedDelimiter{\abs}{\lvert}{\rvert}
\providecommand{\norm}{}\let\norm\relax\DeclarePairedDelimiter{\norm}{\lVert}{\rVert}
\usepackage[table]{xcolor} % [table] : fournit aussi \rowcolors (alternance de lignes)
\usepackage{pagecolor}
\usepackage{tikz}
\usetikzlibrary{arrows.meta,positioning,calc,shapes.geometric,shapes.misc,shapes.arrows,shapes.symbols,decorations.pathmorphing,decorations.pathreplacing,decorations.markings,patterns,shadows,mindmap,trees,fit,backgrounds,matrix,intersections,angles,quotes}
\usepackage{pgfplots}
\usepackage[version=4]{mhchem} % équations nucléaires \ce{^{235}_{92}U}
\usepackage[european,siunitx]{circuitikz} % schémas électriques
\pgfplotsset{compat=1.18}
\usepgfplotslibrary{fillbetween,groupplots} % aires entre courbes (calcul intégral)
\usepackage{enumitem}
\usepackage{fancyhdr}
% [most] charge toutes les bibliothèques tcolorbox (listings, minted, raster,
% documentation...) et gonfle inutilement la mémoire TeX sur un document long
% (~300 boîtes) : on ne charge que ce qui est réellement utilisé.
\usepackage[skins,breakable]{tcolorbox}
\usepackage{graphicx}
\usepackage{colortbl}
\usepackage{booktabs}
\usepackage{tabularx}
\usepackage{diagbox} % case d'en-tête diagonale des tableaux à double entrée
% Colonnes extensibles centrée / gauche / droite (C, L, R), souvent utilisées par les modèles
\newcolumntype{C}{>{\centering\arraybackslash}X}
\newcolumntype{L}{>{\raggedright\arraybackslash}X}
\newcolumntype{R}{>{\raggedleft\arraybackslash}X}
\usepackage{array}
\usepackage{multirow}
\usepackage{multicol}
\usepackage{siunitx}
% parse-numbers=false : accepte aussi \SI{2,5 \times 10^{-3}}{...}
\sisetup{locale=FR,output-decimal-marker={,},group-minimum-digits=4,parse-numbers=false}
\DeclareSIUnit\ohm{\ensuremath{\symup{\Omega}}} % U+2126 absent de la police
\DeclareSIUnit\division{div} % réglages d'oscilloscope (ms/div, V/div)
\DeclareSIUnit\tr{tr} % tours (vitesse de rotation, ex. \SI{1500}{\tr\per\minute})
\DeclareSIUnit\bit{bit}
\DeclareSIUnit\byte{o} % octet
\DeclareSIUnit\inch{po} % pouce
\DeclareSIUnit\ppp{ppp} % points par pouce (résolution d'image)
\usepackage{qrcode}
% Blocs de code source (PHP, C, HTML, JavaScript, SQL) — spécifique à la
% matière Informatique : listings + habillage aux couleurs OnBuch+ (voir le
% style \lstset ci-dessous, à côté des autres boîtes pédagogiques).
\usepackage{listings}
% \degree : commande fréquemment utilisée par les modèles pour un angle ou une
% température en dehors de \SI{}{} (non fournie par les paquets chargés).
\providecommand{\degree}{^{\circ}}
% \qed (fin de démonstration), utilisé par les modèles sans amsthm
\providecommand{\qed}{\hfill$\square$}
% \barre : double barre des valeurs interdites dans un tableau de variations (tkz-tab)
\providecommand{\barre}{\ensuremath{\|}}
% environnement proof (amsthm non chargé) utilisé par les modèles
\ifdefined\proof\else\newenvironment{proof}[1][Démonstration]{\par\noindent\textit{#1.}\ }{\qed\par}\fi
\usepackage{lastpage}
\usepackage{hyperref}
\hypersetup{hidelinks}

% ============================================================
% Palette « Pop » OnBuch+ — mêmes hex que la classe P (plus_ui.dart)
% ============================================================
\definecolor{poporange}{HTML}{F59321}
\definecolor{popdark}{HTML}{E07A0C}
\definecolor{popcream}{HTML}{FFF6E8}
\definecolor{popink}{HTML}{1C1714}
\definecolor{poppurple}{HTML}{7A5AE0}
\definecolor{popgreen}{HTML}{2E9E6A}
\definecolor{poppink}{HTML}{E0457B}
\definecolor{popblue}{HTML}{2F7DE1}
\definecolor{popgold}{HTML}{C98A12}
\definecolor{popmuted}{HTML}{8A7F76}
\definecolor{popline}{HTML}{EADFD2}
\definecolor{popgray}{HTML}{8A7F76}
\colorlet{popgrey}{popgray}
% Alias tolérants (le modèle nomme parfois les couleurs différemment)
\colorlet{popwhite}{white}
\colorlet{popblack}{popink}
\colorlet{popred}{poppink}
\colorlet{popyellow}{popgold}
% Teintes claires pour fonds de tableaux / zones
\colorlet{popblueL}{popblue!10!white}
\colorlet{poporangeL}{poporange!14!white}
\colorlet{popgreenL}{popgreen!12!white}
\colorlet{poppurpleL}{poppurple!12!white}
\colorlet{poppinkL}{poppink!12!white}
\colorlet{popgoldL}{popgold!14!white}

\pagecolor{popcream}

% ============================================================
% Typographie
% ============================================================
\defaultfontfeatures{Ligatures=TeX}
\setmainfont{PlusJakartaSans-Medium.ttf}[Path=../../../fonts/,BoldFont=PlusJakartaSans-Bold.ttf]
% Maths en sans-serif (Fira Math) avec chiffres et lettres droites de Plus
% Jakarta, pour que les formules se fondent dans le texte.
\usepackage[math-style=ISO,bold-style=ISO]{unicode-math}
\setmathfont{FiraMath-Regular.otf}[Path=../../../fonts/]
\setmathfont{PlusJakartaSans-Medium.ttf}[Path=../../../fonts/,range={up/num,up/{Latin,latin},\mathup}]
\usepackage{icomma} % virgule décimale sans espace en mode math
% Caractères mathématiques Unicode absents des polices de texte (Plus Jakarta,
% Archivo Black) : rendus par leur équivalent mathématique, en texte comme en
% formule — sinon ils s'affichent en carré vide (ex. « Arithmétique dans ℤ »).
\usepackage{newunicodechar}
\newunicodechar{ℝ}{\ensuremath{\mathbb{R}}}
\newunicodechar{ℤ}{\ensuremath{\mathbb{Z}}}
\newunicodechar{ℂ}{\ensuremath{\mathbb{C}}}
\newunicodechar{ℕ}{\ensuremath{\mathbb{N}}}
\newunicodechar{ℚ}{\ensuremath{\mathbb{Q}}}
\newunicodechar{𝔻}{\ensuremath{\mathbb{D}}}
\newunicodechar{α}{\ensuremath{\alpha}}
\newunicodechar{β}{\ensuremath{\beta}}
\newunicodechar{γ}{\ensuremath{\gamma}}
\newunicodechar{δ}{\ensuremath{\delta}}
\newunicodechar{ε}{\ensuremath{\varepsilon}}
\newunicodechar{θ}{\ensuremath{\theta}}
\newunicodechar{λ}{\ensuremath{\lambda}}
\newunicodechar{μ}{\ensuremath{\mu}}
\newunicodechar{σ}{\ensuremath{\sigma}}
\newunicodechar{τ}{\ensuremath{\tau}}
\newunicodechar{φ}{\ensuremath{\varphi}}
\newunicodechar{ψ}{\ensuremath{\psi}}
\newunicodechar{ω}{\ensuremath{\omega}}
\newunicodechar{Σ}{\ensuremath{\Sigma}}
% majuscules produites par \MakeUppercase dans les titres (α-aminés -> Α)
\newunicodechar{Α}{\ensuremath{\alpha}}
\newunicodechar{Β}{\ensuremath{\beta}}
\newunicodechar{Γ}{\ensuremath{\Gamma}}
\newunicodechar{Δ}{\ensuremath{\Delta}}
\newunicodechar{Ε}{\ensuremath{\varepsilon}}
\newunicodechar{Θ}{\ensuremath{\Theta}}
\newunicodechar{Λ}{\ensuremath{\Lambda}}
\newunicodechar{Μ}{\ensuremath{\mu}}
\newunicodechar{Τ}{\ensuremath{\tau}}
\newunicodechar{Φ}{\ensuremath{\Phi}}
\newunicodechar{Ψ}{\ensuremath{\Psi}}
\newunicodechar{Ω}{\ensuremath{\Omega}}
\newunicodechar{↦}{\ensuremath{\mapsto}}
\newunicodechar{∈}{\ensuremath{\in}}
\newunicodechar{∉}{\ensuremath{\notin}}
\newunicodechar{∧}{\ensuremath{\wedge}}
\newunicodechar{⇔}{\ensuremath{\Leftrightarrow}}
\newunicodechar{⟺}{\ensuremath{\Longleftrightarrow}}
\newunicodechar{⇒}{\ensuremath{\Rightarrow}}
\newunicodechar{⟹}{\ensuremath{\Longrightarrow}}
\newunicodechar{∘}{\ensuremath{\circ}}
\newunicodechar{∪}{\ensuremath{\cup}}
\newunicodechar{∩}{\ensuremath{\cap}}
\newunicodechar{≡}{\ensuremath{\equiv}}
\newunicodechar{⊂}{\ensuremath{\subset}}
\newunicodechar{⊥}{\ensuremath{\perp}}
\newunicodechar{∃}{\ensuremath{\exists}}
\newunicodechar{∀}{\ensuremath{\forall}}
\newunicodechar{∛}{\ensuremath{\sqrt[3]{\,}}}
\newunicodechar{∜}{\ensuremath{\sqrt[4]{\,}}}
\newunicodechar{⃗}{\ensuremath{{}^{\to}}}
\newunicodechar{ⁿ}{\ifmmode^{n}\else\textsuperscript{\ensuremath{n}}\fi}
\newunicodechar{ˣ}{\ifmmode^{x}\else\textsuperscript{\ensuremath{x}}\fi}
\newunicodechar{ᵇ}{\ifmmode^{b}\else\textsuperscript{\ensuremath{b}}\fi}
\newunicodechar{ʳ}{\ifmmode^{r}\else\textsuperscript{\ensuremath{r}}\fi}
\newunicodechar{ᵘ}{\ifmmode^{u}\else\textsuperscript{\ensuremath{u}}\fi}
\newunicodechar{ʸ}{\ifmmode^{y}\else\textsuperscript{\ensuremath{y}}\fi}
\newunicodechar{ᵃ}{\ifmmode^{a}\else\textsuperscript{\ensuremath{a}}\fi}
\newunicodechar{ᵉ}{\ifmmode^{e}\else\textsuperscript{\ensuremath{e}}\fi}
\newunicodechar{ᵏ}{\ifmmode^{k}\else\textsuperscript{\ensuremath{k}}\fi}
\newunicodechar{ᵖ}{\ifmmode^{p}\else\textsuperscript{\ensuremath{p}}\fi}
\newunicodechar{ᴺ}{\ifmmode^{N}\else\textsuperscript{\ensuremath{N}}\fi}
\newunicodechar{ᶜ}{\ifmmode^{c}\else\textsuperscript{\ensuremath{c}}\fi}
\newunicodechar{ʰ}{\ifmmode^{h}\else\textsuperscript{\ensuremath{h}}\fi}
\newunicodechar{ᵠ}{\ifmmode^{q}\else\textsuperscript{\ensuremath{q}}\fi}
\newunicodechar{ₙ}{\ifmmode_{n}\else\textsubscript{\ensuremath{n}}\fi}
\newunicodechar{ₐ}{\ifmmode_{a}\else\textsubscript{\ensuremath{a}}\fi}
\newunicodechar{ᵢ}{\ifmmode_{i}\else\textsubscript{\ensuremath{i}}\fi}
\newunicodechar{ᵣ}{\ifmmode_{r}\else\textsubscript{\ensuremath{r}}\fi}
\newunicodechar{ₖ}{\ifmmode_{k}\else\textsubscript{\ensuremath{k}}\fi}
\newunicodechar{ᵦ}{\ifmmode_{\beta}\else\textsubscript{\ensuremath{\beta}}\fi}
\newunicodechar{ₓ}{\ifmmode_{x}\else\textsubscript{\ensuremath{x}}\fi}
\newfontfamily\popdisplay{ArchivoBlack-Regular.ttf}[Path=../../../fonts/]
\newfontfamily\poplogo{SpaceGrotesk-Bold.ttf}[Path=../../../fonts/]
\newfontfamily\popsemi{PlusJakartaSans-SemiBold.ttf}[Path=../../../fonts/]
\newfontfamily\popxbold{PlusJakartaSans-ExtraBold.ttf}[Path=../../../fonts/]
\newfontfamily\popxboldls{PlusJakartaSans-ExtraBold.ttf}[Path=../../../fonts/,LetterSpace=3.0]

\setlist[enumerate]{leftmargin=1.7em,itemsep=5pt,topsep=4pt}
\setlist[itemize]{leftmargin=1.7em,itemsep=4pt,topsep=4pt,label={\color{poporange}\textbullet}}
\setlength{\parindent}{0pt}
\setlength{\parskip}{5pt}
\renewcommand{\arraystretch}{1.25}

% pgfplots : style Pop par défaut
\pgfplotsset{
  every axis/.append style={axis line style={popink,line width=1pt},tick style={popink},
    grid style={popline,line width=0.5pt},label style={font=\small},tick label style={font=\footnotesize}},
  popcurve/.style={poporange,line width=1.8pt,no markers,smooth},
}
% Même style disponible dans un \draw TikZ simple (hors pgfplots)
\tikzset{popcurve/.style={poporange,line width=1.8pt}}
\tikzset{popgrid/.style={popline,very thin}}
\colorlet{popcurve}{poporange} % le nom du style est parfois utilisé comme couleur

% ============================================================
% Pill / squiggle
% ============================================================
\newcommand{\poppill}[3]{%
  \tikz[baseline=-0.55ex]\node[draw=popink,line width=1.2pt,rounded corners=2.6pt,%
    fill=#2,inner xsep=6.5pt,inner ysep=3pt]{%
    \popxboldls\fontsize{7}{7}\selectfont\color{#3}\MakeUppercase{#1}};%
}
\newcommand{\popsquiggle}[1]{%
  \tikz[baseline=-0.4ex]{
    \draw[#1,line width=2.6pt,line cap=round]
      plot [smooth] coordinates {(0,0.09) (0.34,0.01) (0.68,0.21) (1.02,0.01) (1.36,0.10)};
  }%
}
% Mot-clé surligné (marqueur orange sous le texte)
\newcommand{\cle}[1]{{\popxbold\color{popdark}#1}}

% ============================================================
% En-tête / pied
% ============================================================
\pagestyle{fancy}
\fancyhf{}
\renewcommand{\headrulewidth}{0pt}
\renewcommand{\footrulewidth}{0pt}
\lhead{\raisebox{-0.15ex}{\poplogo\fontsize{13}{13}\selectfont\color{popink}OnBuch\color{poporange}+}}
\rhead{\poppill{\DOCMATIERE\ \textbullet\ \DOCNIVEAU\ \textbullet\ Leçon \DOCLECON}{white}{popink}}
\lfoot{\popxbold\fontsize{7.6}{7.6}\selectfont\color{popmuted}\MakeUppercase{Cours signé OnBuch+}\ \textbullet\ {\popsemi\DOCTITRE}}
\rfoot{\tikz[baseline=-0.6ex]\node[rounded corners=8.5pt,fill=popink,minimum height=17pt,minimum width=17pt,inner xsep=5pt,inner sep=0]{\popxbold\fontsize{8}{8}\selectfont\color{popcream}\thepage\,/\,\pageref*{LastPage}};}

% ============================================================
% Titres
% ============================================================
\newcommand{\chaptitre}[2]{%
  \begin{tcolorbox}[enhanced,colback=poporange,colframe=popink,boxrule=2.6pt,arc=11pt,
      drop shadow=popink,left=15pt,right=15pt,top=12pt,bottom=11pt,before skip=3pt,after skip=18pt]
    \poppill{#2}{popink}{popcream}\par\vspace{9pt}
    {\raggedright\hyphenpenalty=10000\exhyphenpenalty=10000\popdisplay\fontsize{22}{24}\selectfont\color{popcream}\MakeUppercase{#1}\par}
  \end{tcolorbox}
}
% Section numérotée : pastille orange avec le numéro + titre Archivo
\newcounter{popsec}
\newcommand{\coursec}[1]{%
  \stepcounter{popsec}\setcounter{popsub}{0}%
  \par\vspace{16pt}\noindent
  \begin{minipage}{\linewidth}
  \tikz[baseline=-0.7ex]\node[draw=popink,line width=1.6pt,fill=poporange,rounded corners=4pt,minimum size=22pt,inner sep=0]{\popdisplay\fontsize{12}{12}\selectfont\color{popcream}\thepopsec};%
  \hspace{8pt}{\popdisplay\fontsize{15}{17}\selectfont\color{popink}\MakeUppercase{#1}}\par
  \vspace{4pt}\noindent{\color{poporange}\rule{36pt}{3.6pt}}
  \end{minipage}\par\nopagebreak\vspace{8pt}
}
\newcounter{popsub}
\newcommand{\courssub}[1]{%
  \stepcounter{popsub}%
  \par\vspace{9pt}\noindent{\popxbold\fontsize{11.5}{13}\selectfont\color{popink}{\color{poporange}\thepopsec.\thepopsub}\hspace{6pt}#1}\par\nopagebreak\vspace{4pt}
}

% ============================================================
% Boîtes pédagogiques — même recette : fond blanc, bordure épaisse colorée,
% ombre dure encre, badge plein. Titre optionnel [#1].
% ============================================================
\tcbset{popbase/.style={breakable,enhanced,colback=white,boxrule=2.3pt,arc=9pt,
  shadow={3pt}{-3pt}{0pt}{popink},left=14pt,right=14pt,top=11pt,bottom=11pt,before skip=11pt,after skip=15pt,
  fontupper=\popsemi\fontsize{10.6}{14.5}\selectfont\color{popink},
  fontlower=\popsemi\fontsize{10.6}{14.5}\selectfont\color{popink}}}
% Coche dessinée (le glyphe ✓ n'existe pas dans les polices du document)
\newcommand{\popcheck}[1]{\tikz[baseline=-0.5ex]\draw[#1,line width=1.6pt,line cap=round,line join=round](-0.13,0.0)--(-0.04,-0.09)--(0.13,0.1);}
\providecommand{\coche}{\popcheck{popgreen}}
\providecommand{\resultat}[1]{\par\smallskip\noindent\fcolorbox{popgreen}{popgreenL}{\parbox{\dimexpr\linewidth-2\fboxsep-2\fboxrule\relax}{\textbf{Résultat.} #1}}\par\smallskip}
\newcommand{\popboxhead}[3]{\poppill{#1}{#2}{white}\ifx\relax#3\relax\else\hspace{6pt}{\popxbold\fontsize{10.6}{12}\selectfont\color{popink}#3}\fi\par\vspace{7pt}}

\newtcolorbox{aretenir}[1][]{popbase,colframe=popgreen,before upper={\popboxhead{À retenir}{popgreen}{#1}}}
\newtcolorbox{exemplebox}[1][]{popbase,colframe=poporange,before upper={\popboxhead{Exemple}{poporange}{#1}}}
\newtcolorbox{definition}[1][]{popbase,colframe=popblue,before upper={\popboxhead{Définition}{popblue}{#1}}}
\newtcolorbox{propriete}[1][]{popbase,colframe=poppurple,before upper={\popboxhead{Propriété}{poppurple}{#1}}}
\newtcolorbox{methode}[1][]{popbase,colframe=popgold,before upper={\popboxhead{Méthode}{popgold}{#1}}}
\newtcolorbox{attention}[1][]{popbase,colframe=poppink,before upper={\popboxhead{Attention}{poppink}{#1}}}
\newtcolorbox{experience}[1][]{popbase,colframe=popblue,colback=popblueL,before upper={\popboxhead{Expérience}{popblue}{#1}}}
% « Le savais-tu ? » : carte violette pleine (contexte, culture, Cameroun)
\newtcolorbox{savaistu}[1][]{popbase,colback=poppurple,colframe=popink,
  fontupper=\popsemi\fontsize{10.4}{14.2}\selectfont\color{white},
  before upper={\poppill{Le savais-tu\,?}{white}{poppurple}\ifx\relax#1\relax\else\hspace{6pt}{\popxbold\color{white}#1}\fi\par\vspace{7pt}}}
% Exercice résolu : énoncé en haut, \tcblower puis la solution sur fond crème
\newcounter{popexo}\newcounter{popexr}
\newtcolorbox{exoresolu}[1][]{popbase,colframe=poporange,colbacklower=popcream,
  segmentation style={popink,line width=1.2pt,dashed},
  before upper={\stepcounter{popexr}\popboxhead{Exercice résolu \thepopexr}{poporange}{#1}},
  before lower={\poppill{Solution}{popink}{popcream}\par\vspace{6pt}}}
% Exercice (non résolu en place) — la correction est dans la partie Corrigés
\newtcolorbox{exercice}[1][]{popbase,colframe=popink,
  before upper={\stepcounter{popexo}\popboxhead{Exercice \thepopexo}{popink}{#1}}}
% Corrigé
\newtcolorbox{corrige}[1][]{popbase,colframe=popgreen,colback=popgreenL,
  before upper={\popboxhead{Corrigé}{popgreen}{#1}}}
% Objectifs / prérequis / bilan
\newtcolorbox{objectifs}{popbase,colframe=popink,colback=poporangeL,
  before upper={\popboxhead{Objectifs de la leçon}{popink}{}}}
\newtcolorbox{prerequis}{popbase,colframe=popink,colback=popblueL,
  before upper={\popboxhead{Prérequis}{popblue}{}}}
\newtcolorbox{bilan}{popbase,colframe=popink,colback=white,boxrule=2.8pt,
  before upper={\popboxhead{Fiche bilan}{poporange}{L'essentiel à connaître pour l'examen}}}
% Figure encadrée avec légende
\newtcolorbox{popfigure}[1][]{enhanced,breakable=false,colback=white,colframe=popline,boxrule=1.4pt,arc=8pt,
  left=10pt,right=10pt,top=10pt,bottom=8pt,before skip=10pt,after skip=12pt,halign=center,
  lower separated=false,halign lower=center,
  fontlower=\popsemi\fontsize{9}{11.5}\selectfont\color{popmuted},
  #1}
\newcommand{\legende}[1]{\par\vspace{6pt}{\popsemi\fontsize{9}{11.5}\selectfont\color{popmuted}#1}}

% ============================================================
% Bloc de code source — \begin{lstlisting}[language=Php]...\end{lstlisting}
% (ou language=C, HTML, JavaScript, SQL ; option omise = pseudo-code brut).
% Même recette visuelle que les figures (\popfigure) : cadre encre épais,
% fond clair (popcream), légende possible juste après via \legende{...}
% comme pour une figure (listings ne sait pas arrondir ses coins : on garde
% un cadre rectangulaire, cohérent avec les autres tableaux du document).
% CONTRAT : les modèles ne doivent utiliser QUE \begin{lstlisting}[...] tel
% quel (pas de \begin{tcblisting}, pas de \verb pour un bloc multi-lignes).
% ============================================================
\lstdefinestyle{poplst}{
  backgroundcolor=\color{popcream},
  basicstyle=\popsemi\fontsize{8.6}{11}\selectfont\color{popink},
  keywordstyle=\bfseries\color{popblue},
  commentstyle=\itshape\color{popmuted},
  stringstyle=\color{popgreen},
  numberstyle=\tiny\color{popmuted},
  identifierstyle=\color{popink},
  frame=l,
  rulecolor=\color{popink},
  framerule=1.6pt,
  framesep=7pt,
  framexleftmargin=7pt,
  framexrightmargin=7pt,
  xleftmargin=2pt,
  xrightmargin=2pt,
  breaklines=true,
  breakatwhitespace=false,
  showstringspaces=false,
  showspaces=false,
  showtabs=false,
  tabsize=2,
  columns=flexible,
  keepspaces=true,
  numbers=none,
  aboveskip=10pt,
  belowskip=10pt,
  captionpos=b,
}
\lstset{style=poplst, language=PHP}
% La distribution TeX embarquée ne fournit pas le fichier de langue
% JavaScript de listings (lstlang*.dat) : on la redéfinit nous-mêmes
% pour que \begin{lstlisting}[language=JavaScript] fonctionne.
\lstdefinelanguage{JavaScript}{
  keywords={break,case,catch,class,const,continue,debugger,default,delete,do,
    else,export,extends,finally,for,function,if,import,in,instanceof,let,new,
    return,static,super,switch,this,throw,try,typeof,var,void,while,with,yield,
    async,await,of,get,set},
  keywordstyle=\bfseries\color{popblue},
  ndkeywords={true,false,null,undefined,NaN,Infinity},
  ndkeywordstyle=\bfseries\color{poporange},
  identifierstyle=\color{popink},
  sensitive=true,
  comment=[l]{//},
  morecomment=[s]{/*}{*/},
  string=[b]",
  morestring=[b]',
  morestring=[b]`,
}
% Même limitation pour CSS et les fichiers de configuration (apacheconf,
% nginx, ini...) : ils ne sont pas fournis. CSS et un style générique de
% type "config" (commentaires # ou ;, pas de mots-clés) couvrent les besoins.
\lstdefinelanguage{CSS}{
  keywords={color,background,background-color,margin,padding,border,
    border-radius,font,font-size,font-weight,font-family,width,height,
    display,position,top,left,right,bottom,float,clear,text-align,
    line-height,list-style,cursor,overflow,box-shadow,transition,
    flex,grid,align-items,justify-content,z-index},
  keywordstyle=\bfseries\color{popblue},
  identifierstyle=\color{popink},
  sensitive=false,
  comment=[s]{/*}{*/},
  string=[b]",
  morestring=[b]',
}
\lstdefinelanguage{apacheconf}{
  sensitive=false,
  morecomment=[l]{\#},
}
\lstdefinelanguage{Apache}{
  sensitive=false,
  morecomment=[l]{\#},
}
\lstdefinelanguage{text}{sensitive=false}
\lstdefinelanguage{powershell}{
  keywords={New-Object,New-SmbShare,Get-Acl,Set-Acl,Get-ChildItem,Set-Content,
    Get-Content,Write-Host,ForEach-Object,Where-Object,Select-Object,
    Test-Path,Remove-Item,Copy-Item,Move-Item,Start-Process,Stop-Process,
    If,Else,ElseIf,ForEach,While,Function,Return,Try,Catch},
  keywordstyle=\bfseries\color{popblue},
  identifierstyle=\color{popink},
  sensitive=false,
  comment=[l]{\#},
  string=[b]",
  morestring=[b]',
}
\lstdefinelanguage{ini}{
  sensitive=false,
  morecomment=[l]{;},
  morecomment=[l]{\#},
}

% ============================================================
% Signature OnBuch+
% ============================================================
\newcommand{\poplogoinline}[1]{{\poplogo\fontsize{#1}{#1}\selectfont\color{popink}OnBuch\color{poporange}+}}
\newcommand{\signatureonbuch}{%
  \begin{tcolorbox}[enhanced,colback=popink,colframe=popink,boxrule=2.2pt,arc=10pt,
      drop shadow=poporange,left=14pt,right=14pt,top=10pt,bottom=10pt,before skip=0pt,after skip=0pt]
    \tikz[baseline=-0.7ex]\node[circle,fill=popgreen,draw=popcream,line width=1.3pt,minimum size=22pt,inner sep=0]{\popcheck{white}};%
    \hspace{9pt}\begin{minipage}[c]{0.62\linewidth}
      {\popxbold\fontsize{12}{13}\selectfont\color{popcream}Signé\ \textbullet\ {\poplogo OnBuch\color{poporange}+}}\par\vspace{2pt}
      {\raggedright\popsemi\fontsize{8}{10.5}\selectfont\color{popline}Ludovic A.\\\MakeUppercase{Conforme au programme officiel MINESEC}\par}
    \end{minipage}\hfill
    {\poplogo\fontsize{22}{22}\selectfont\color{popcream}OnBuch\color{poporange}+}
  \end{tcolorbox}%
}

% ============================================================
% Page de garde
% #1 titre · #2 matière · #3 niveau · #4 module · #5 n° leçon · #6 durée ·
% #7 description / objectifs (texte ou itemize)
% ============================================================
\newcommand{\pagegarde}[7]{%
  \thispagestyle{empty}
  \newgeometry{a4paper,margin=1.7cm,top=1.6cm,bottom=1.4cm}%
  \begin{tikzpicture}[remember picture,overlay]
    \fill[poporange!18] ([xshift=-3.2cm,yshift=-3.6cm]current page.north east) circle (4.6cm);
    \fill[poppurple!14] ([xshift=2.4cm,yshift=7.6cm]current page.south west) circle (3.8cm);
  \end{tikzpicture}%
  {\poplogo\fontsize{30}{30}\selectfont\color{popink}OnBuch\color{poporange}+}\hfill
  \raisebox{0.6ex}{\poppill{Cours signé \textbullet\ Premium}{popink}{popcream}}\par
  \vspace{1pt}\popsquiggle{poppurple}\par
  \vspace{8pt}
  \begin{tcolorbox}[enhanced,colback=poporange,colframe=popink,boxrule=2.8pt,arc=13pt,
      drop shadow=popink,left=18pt,right=18pt,top=12pt,bottom=13pt,after skip=10pt]
    \poppill{#2 \textbullet\ #3}{popink}{popcream}\hspace{5pt}\poppill{#4}{white}{popink}\par\vspace{8pt}
    {\popdisplay\fontsize{14}{14}\selectfont\color{popink}LEÇON #5}\par\vspace{4pt}
    {\raggedright\hyphenpenalty=10000\exhyphenpenalty=10000\popdisplay\fontsize{25}{27}\selectfont\color{popcream}\MakeUppercase{#1}\par}
    \vspace{8pt}
    {\popxbold\fontsize{9.5}{9.5}\selectfont\color{popink}\MakeUppercase{Durée conseillée : #6}}
  \end{tcolorbox}
  \begin{tcolorbox}[enhanced,colback=white,colframe=popink,boxrule=2.2pt,arc=10pt,
      drop shadow=popink,left=16pt,right=16pt,top=10pt,bottom=10pt,after skip=8pt]
    \poppill{Dans cette leçon}{poppurple}{white}\par\vspace{5pt}
    \popsemi\fontsize{9.3}{12.6}\selectfont\color{popink}#7
  \end{tcolorbox}
  \noindent\begin{minipage}[t]{0.487\linewidth}
    \begin{tcolorbox}[enhanced,colback=popgreen,colframe=popink,boxrule=2.2pt,arc=10pt,
        drop shadow=popink,left=11pt,right=11pt,top=10pt,bottom=10pt,height=3.1cm,valign=center]
      \tcbox[colback=white,colframe=popink,boxrule=1.3pt,arc=5pt,boxsep=0pt,nobeforeafter,
        left=3pt,right=3pt,top=3pt,bottom=3pt,valign=center]{\qrcode[height=1.9cm]{https://play.google.com/store/apps/details?id=com.luvvix.onbuch&hl=fr}}%
      \hspace{8pt}\begin{minipage}[c]{0.5\linewidth}
        {\popdisplay\fontsize{11}{12.5}\selectfont\color{white}\MakeUppercase{Télécharge OnBuch}}\par\vspace{4pt}
        {\raggedright\popsemi\fontsize{8.4}{11}\selectfont\color{white}Gratuit \textbullet\ Google Play\\Tuteur IA, annales, concours, résultats\par}
      \end{minipage}
    \end{tcolorbox}
  \end{minipage}\hfill
  \begin{minipage}[t]{0.487\linewidth}
    \begin{tcolorbox}[enhanced,colback=poppurple,colframe=popink,boxrule=2.2pt,arc=10pt,
        drop shadow=popink,left=11pt,right=11pt,top=10pt,bottom=10pt,height=3.1cm,valign=center]
      \tikz[baseline=-0.7ex]\node[circle,fill=white,draw=popink,line width=1.3pt,minimum size=1.6cm,inner sep=0]{\popdisplay\fontsize{24}{24}\selectfont\color{poppurple}+};%
      \hspace{8pt}\begin{minipage}[c]{0.6\linewidth}
        {\popdisplay\fontsize{11}{12.5}\selectfont\color{white}\MakeUppercase{Passe à OnBuch+}}\par\vspace{4pt}
        {\raggedright\popsemi\fontsize{8.4}{11}\selectfont\color{white}Tous les cours signés, Tuteur IA illimité, fiches PDF — onglet OnBuch+ de l'app\par}
      \end{minipage}
    \end{tcolorbox}
  \end{minipage}
  \vspace{8pt}
  \popcontenu
  \vfill
  \signatureonbuch
  \restoregeometry
  \newpage
}

% Rangée « Ce cours contient » (page de garde)
\newcommand{\poptile}[3]{% #1 couleur #2 gros texte #3 légende
  \begin{tcolorbox}[enhanced,colback=#1,colframe=popink,boxrule=2pt,arc=9pt,shadow={2.5pt}{-2.5pt}{0pt}{popink},
      left=6pt,right=6pt,top=9pt,bottom=9pt,halign=center,valign=center,height=1.75cm,width=\linewidth,nobeforeafter]
    {\popdisplay\fontsize{12.5}{13}\selectfont\color{white}\MakeUppercase{#2}}\par\vspace{3pt}
    {\centering\popsemi\fontsize{7.8}{9.5}\selectfont\color{white}#3\par}
  \end{tcolorbox}}
\newcommand{\popcontenu}{%
  \par\noindent{\popxboldls\fontsize{7.5}{7.5}\selectfont\color{popmuted}CE COURS SIGNÉ CONTIENT}\par\vspace{7pt}
  \noindent\begin{minipage}[t]{0.235\linewidth}\poptile{popblue}{Cours}{complet, illustré, pas à pas}\end{minipage}\hfill
  \begin{minipage}[t]{0.235\linewidth}\poptile{poporange}{Méthodes}{et exercices résolus}\end{minipage}\hfill
  \begin{minipage}[t]{0.235\linewidth}\poptile{poppink}{Exercices}{type examen + corrigés}\end{minipage}\hfill
  \begin{minipage}[t]{0.235\linewidth}\poptile{popgreen}{Fiche bilan}{l'essentiel à retenir}\end{minipage}\par}

% Sommaire
\newtcolorbox{sommaire}{enhanced,colback=white,colframe=popink,boxrule=2.2pt,arc=10pt,
  drop shadow=popink,left=18pt,right=18pt,top=13pt,bottom=13pt,before skip=6pt,after skip=20pt,
  before upper={\poppill{Sommaire}{poppurple}{white}\par\vspace{9pt}\popsemi\fontsize{10.6}{14.5}\selectfont\color{popink}}}

% Signature de fin de cours — poussée tout en bas de la dernière page
\newcommand{\findecours}{%
  \par\vfill\nopagebreak
  \begin{center}\popsquiggle{poporange}\end{center}\vspace{4pt}
  \signatureonbuch
}
\AtBeginDocument{\def\llbracket{\mathopen{[\mkern-3.5mu[}}\def\rrbracket{\mathclose{]\mkern-3.5mu]}}} % intervalles d'entiers (glyphe ⟦ absent de la police)
% Glyphes absents de Fira Math : redéfinis à partir de glyphes disponibles
\AtBeginDocument{%
  \def\setminus{\mathbin{\backslash}}\let\smallsetminus\setminus
  \def\vdots{\mathord{\vbox{\baselineskip3.2pt\lineskiplimit0pt\kern4pt\hbox{.}\hbox{.}\hbox{.}}}}%
  \def\ddots{\mathinner{\mkern1mu\raise6.5pt\hbox{.}\mkern2mu\raise3.6pt\hbox{.}\mkern2mu\raise0.7pt\hbox{.}\mkern1mu}}%
  \def\triangle{\mathord{\scalebox{0.9}{$\symup{\Delta}$}}}%
  \def\nabla{\mathord{\raisebox{\depth}{\scalebox{1}[-1]{$\symup{\Delta}$}}}}%
  \def\checkmark{\popcheck{popdark}}%
  \def\star{\mathbin{\ast}}\def\bigstar{\mathord{\ast}}%
  \def\diamond{\mathbin{\scalebox{0.6}{$\Diamond$}}}%
  \def\bigcup{\mathop{\mathchoice{\raisebox{-0.3em}{\scalebox{1.7}{$\cup$}}}{\scalebox{1.25}{$\cup$}}{\cup}{\cup}}\displaylimits}%
  \def\bigcap{\mathop{\mathchoice{\raisebox{-0.3em}{\scalebox{1.7}{$\cap$}}}{\scalebox{1.25}{$\cap$}}{\cap}{\cap}}\displaylimits}%
  \def\lesseqqgtr{\mathrel{\lessgtr}}\def\bigoplus{\mathop{\oplus}\displaylimits}%
}
\providecommand{\ssi}{si et seulement si} % abréviation fréquente
\AtBeginDocument{\providecommand{\wideparen}{\overparen}} % arc de cercle

%% FIGFIX-BEGIN v5 — correctifs globaux des figures (docs/onbuch-generation/tools/figfix)
\usepackage{adjustbox}\usepackage{etoolbox}\tcbuselibrary{raster}
% toute figure TikZ/pgfplots plus large que la ligne est réduite à la largeur disponible
\BeforeBeginEnvironment{tikzpicture}{\begin{adjustbox}{max width=\linewidth}}
\AfterEndEnvironment{tikzpicture}{\end{adjustbox}}
% légendes pgfplots lisibles par-dessus les courbes
\pgfplotsset{every axis/.append style={legend style={fill=white,fill opacity=0.92,text opacity=1,draw=black!15,inner sep=3pt}}}
% étiquettes d'arêtes (syntaxe edge["..."]) avec fond blanc
\tikzset{every edge quotes/.append style={fill=white,inner sep=1.2pt}}
%% FIGFIX-END
\begin{document}

\pagegarde{Gestion de la santé et du bien-être social}{Informatique}{3ème}{Informatique – classe de 3ème}{N12}{1 séance (fiche de progression harmonisée)}{Cette leçon sensibilise l'élève aux bonnes pratiques d'utilisation des outils numériques afin de préserver sa santé physique, son équilibre psychologique et son bien-être social. Elle fournit des règles concrètes d'hygiène du numérique applicables dans la vie courante, à l'école comme à la maison.
\vspace{4pt}
\begin{multicols}{2}\small
\begin{itemize}
\item À la fin de cette leçon, je sais définir l'hygiène du numérique et citer ses enjeux pour la santé et le bien-être social.
\item À la fin de cette leçon, je sais identifier les risques physiques liés à un usage excessif ou mal maîtrisé des écrans (fatigue visuelle, troubles du sommeil, mauvaise posture).
\item À la fin de cette leçon, je sais identifier les risques psychologiques et sociaux du numérique (addiction, cyberharcèlement, isolement, comparaison sociale).
\item À la fin de cette leçon, je sais appliquer les règles d'ergonomie et de bon usage des écrans (règle des 20-20-20, posture, luminosité, temps d'écran).
\item À la fin de cette leçon, je sais adopter des comportements responsables sur les réseaux sociaux (respect de la vie privée, e-réputation, netiquette).
\item À la fin de cette leçon, je sais analyser une situation concrète de la vie courante et proposer des solutions d'hygiène numérique adaptées.
\end{itemize}\end{multicols}}

\begin{sommaire}
\begin{multicols}{2}
\begin{enumerate}
\item Notion d'hygiène du numérique
\item Risques physiques liés aux écrans
\item Risques psychologiques et sociaux
\item Les bonnes pratiques d'hygiène numérique
\item Usage responsable des réseaux sociaux et bien-être social
\item Élaborer et justifier une démarche d'hygiène numérique
\item Activité d'intégration
\item Méthodes et exercices résolus
\item Exercices
\item Corrigés des exercices
\item Fiche bilan
\end{enumerate}
\end{multicols}
\end{sommaire}

\begin{objectifs}
\begin{itemize}
\item À la fin de cette leçon, je sais définir l'hygiène du numérique et citer ses enjeux pour la santé et le bien-être social.
\item À la fin de cette leçon, je sais identifier les risques physiques liés à un usage excessif ou mal maîtrisé des écrans (fatigue visuelle, troubles du sommeil, mauvaise posture).
\item À la fin de cette leçon, je sais identifier les risques psychologiques et sociaux du numérique (addiction, cyberharcèlement, isolement, comparaison sociale).
\item À la fin de cette leçon, je sais appliquer les règles d'ergonomie et de bon usage des écrans (règle des 20-20-20, posture, luminosité, temps d'écran).
\item À la fin de cette leçon, je sais adopter des comportements responsables sur les réseaux sociaux (respect de la vie privée, e-réputation, netiquette).
\item À la fin de cette leçon, je sais analyser une situation concrète de la vie courante et proposer des solutions d'hygiène numérique adaptées.
\item À la fin de cette leçon, je sais rédiger et justifier une démarche, une réponse ou une conclusion concernant un problème de bien-être numérique.
\item À la fin de cette leçon, je sais élaborer un plan personnel d'hygiène du numérique et le justifier.
\end{itemize}
\end{objectifs}

\begin{prerequis}
\begin{itemize}
\item Notions de base sur l'ordinateur, le smartphone et Internet (vues en 4ème et en début de 3ème)
\item Utilisation courante des réseaux sociaux et des messageries (WhatsApp, Facebook, TikTok)
\item Notion de sécurité sur Internet et de mot de passe (leçons précédentes de 3ème)
\item Notions élémentaires de santé et d'hygiène de vie (SVT / éducation à la santé)
\item Capacité à rédiger un court texte argumenté
\end{itemize}
\end{prerequis}

\newpage

\coursec{Notion d'hygiène du numérique}

\begin{exemplebox}[Situation de vie : le soir d'Aïcha]
Aïcha est élève de 3ème dans un collège de Yaoundé. Chaque soir, après ses devoirs (parfois avant...), elle passe plus de 4 heures sur son smartphone : réseaux sociaux, jeux en ligne et vidéos. Depuis quelques semaines, elle a mal aux yeux, dort mal, ses notes baissent et elle s'est disputée avec une amie à cause d'une photo publiée sans autorisation. Son père, inquiet, se demande comment l'aider à mieux utiliser ses outils numériques.
\end{exemplebox}

\textbf{Problématique :} comment Aïcha peut-elle adopter une bonne \cle{hygiène du numérique} pour protéger sa santé, son bien-être et ses relations sociales ?

Pour répondre à cette question, commençons par le commencement : que signifie exactement « hygiène du numérique » ?

\courssub{Qu'est-ce que l'hygiène du numérique ?}

\textbf{Intuition : à quel problème répond cette notion ?}

Quand on parle d'« hygiène », tu penses sans doute à l'hygiène du corps : se laver les mains, se brosser les dents, manger équilibré. Ces habitudes ne servent pas à empêcher de vivre, bien au contraire : elles permettent de profiter de la vie \emph{sans tomber malade}. L'hygiène du numérique fonctionne exactement pareil. Le smartphone, l'ordinateur ou la télévision sont des outils formidables ; le problème n'est pas de les utiliser, mais de les utiliser \emph{de la bonne manière, au bon moment, pendant une durée raisonnable}, pour qu'ils restent des alliés et ne deviennent pas des ennemis de la santé, du sommeil et des études.

\begin{definition}[Hygiène du numérique]
L'\cle{hygiène du numérique} est l'ensemble des règles, bonnes pratiques et habitudes d'utilisation des outils numériques (ordinateurs, smartphones, tablettes, consoles de jeux, Internet) visant à préserver la \textbf{santé physique}, l'\textbf{équilibre mental} et le \textbf{bien-être social} de l'utilisateur.
\end{definition}

Autrement dit : utiliser le numérique \emph{sans qu'il nuise} ni à ton corps, ni à ton moral, ni à tes relations avec les autres.

\begin{attention}[Erreur fréquente]
Beaucoup d'élèves croient que l'hygiène du numérique consiste seulement à « ne pas trop jouer ». C'est faux ! Elle concerne aussi la \textbf{posture} devant l'écran (dos, yeux, poignets), le \textbf{sommeil} (écrans tard le soir), les \textbf{relations} (respect des autres en ligne, photos publiées) et la \textbf{vie privée} (informations personnelles partagées). C'est un ensemble complet de bonnes habitudes, pas une seule interdiction.
\end{attention}

\courssub{Les outils numériques de notre quotidien}

Au Cameroun comme ailleurs, les outils numériques sont partout. Les principaux concernés par l'hygiène du numérique sont :

\begin{itemize}
\item le \cle{téléphone mobile} (smartphone) : appels, messages, réseaux sociaux, jeux, vidéos — c'est l'outil le plus utilisé par les jeunes ;
\item l'\cle{ordinateur} : à la maison, au collège (salle informatique) ou au \cle{cybercafé} pour les recherches, les exposés et la bureautique ;
\item la \cle{télévision} : programmes, films, séries ;
\item la \cle{console de jeux} : jeux vidéo seul ou en réseau ;
\item les \cle{réseaux sociaux} (WhatsApp, Facebook, TikTok, etc.) : communication, partage de photos et de vidéos.
\end{itemize}

Chacun de ces outils est utile. Mais chacun peut devenir source de fatigue, de conflit ou de dépendance s'il est mal utilisé.

\begin{savaistu}[Le numérique des jeunes Camerounais]
Au Cameroun, la grande majorité des jeunes accède à Internet principalement via le \textbf{téléphone mobile}, grâce aux forfaits Internet des opérateurs. Beaucoup se connectent tard le soir, allongés dans leur lit, écran à quelques centimètres des yeux. Or la lumière des écrans et l'excitation des jeux ou des discussions retardent l'endormissement et perturbent le sommeil — alors qu'un adolescent a besoin d'environ 8 à 9 heures de sommeil par nuit pour bien apprendre.
\end{savaistu}

\courssub{Les trois dimensions de l'hygiène du numérique}

L'hygiène du numérique repose sur trois dimensions complémentaires. Si l'une d'elles est négligée, tout l'équilibre est menacé — c'est exactement ce qui arrive à Aïcha, qui est touchée sur les trois plans à la fois.

\begin{popfigure}
\begin{center}
\begin{center}\tcbox[colback=popblue,colframe=popblue,coltext=white,arc=9pt,boxsep=4pt,left=6pt,right=6pt]{\bfseries\small Hygiène du numérique}\end{center}
\vspace{-2pt}
\begin{tcbraster}[raster columns=2,raster equal height=rows,raster column skip=6pt,raster row skip=6pt,enhanced,blankest]
\begin{tcolorbox}[enhanced,colback=white,colframe=poporange,boxrule=0.9pt,arc=3pt,title={Santé physique},fonttitle=\bfseries\scriptsize,coltitle=white,colbacktitle=poporange,left=4pt,right=4pt,top=2pt,bottom=2pt,boxsep=1pt,toptitle=1.5pt,bottomtitle=1.5pt,halign=left]
\begin{itemize}[nosep,leftmargin=*,itemsep=1pt,font=\scriptsize]
\item Protéger les yeux et le dos
\item Préserver le sommeil
\end{itemize}
\end{tcolorbox}
\begin{tcolorbox}[enhanced,colback=white,colframe=popgreen,boxrule=0.9pt,arc=3pt,title={Équilibre psychologique},fonttitle=\bfseries\scriptsize,coltitle=white,colbacktitle=popgreen,left=4pt,right=4pt,top=2pt,bottom=2pt,boxsep=1pt,toptitle=1.5pt,bottomtitle=1.5pt,halign=left]
\begin{itemize}[nosep,leftmargin=*,itemsep=1pt,font=\scriptsize]
\item Éviter la dépendance aux jeux
\item Gérer le stress des réseaux
\end{itemize}
\end{tcolorbox}
\begin{tcolorbox}[enhanced,colback=white,colframe=poppurple,boxrule=0.9pt,arc=3pt,title={Bien-être social},fonttitle=\bfseries\scriptsize,coltitle=white,colbacktitle=poppurple,left=4pt,right=4pt,top=2pt,bottom=2pt,boxsep=1pt,toptitle=1.5pt,bottomtitle=1.5pt,halign=left]
\begin{itemize}[nosep,leftmargin=*,itemsep=1pt,font=\scriptsize]
\item Respecter les autres en ligne
\item Garder du temps pour la famille
\end{itemize}
\end{tcolorbox}
\end{tcbraster}

\end{center}
\legende{Les trois dimensions de l'hygiène du numérique, avec deux exemples de bonnes pratiques par dimension.}
\end{popfigure}

\begin{aretenir}[Les trois dimensions]
\begin{enumerate}
\item \textbf{Santé physique} : protéger ses yeux, son dos, ses poignets et son sommeil face aux écrans.
\item \textbf{Équilibre psychologique} : éviter la dépendance, l'anxiété et la fatigue mentale liées à un usage excessif.
\item \textbf{Bien-être social et relationnel} : respecter les autres en ligne, protéger sa vie privée et garder des relations réelles (famille, amis, camarades de classe).
\end{enumerate}
\end{aretenir}

\courssub{Le numérique : un outil formidable quand il est bien utilisé}

Attention à ne pas tout mélanger : la leçon ne dit pas que le numérique est mauvais ! Bien utilisé, il apporte des bénéfices considérables :

\begin{itemize}
\item \textbf{Apprentissage} : rechercher des informations pour un exposé, regarder un cours en vidéo, s'entraîner avec des exercices en ligne, utiliser un dictionnaire numérique ;
\item \textbf{Communication} : échanger avec ses camarades sur les devoirs, appeler sa famille restée au village, contacter un professeur ;
\item \textbf{Information} : suivre l'actualité, connaître les résultats d'examens, découvrir les démarches administratives en ligne ;
\item \textbf{Loisirs} : écouter de la musique, regarder un film, jouer à un jeu... avec modération.
\end{itemize}

Le problème ne vient donc pas de l'outil, mais de l'\textbf{usage} qu'on en fait : la durée, le moment de la journée et le contenu consulté.

\courssub{Pourquoi les adolescents sont-ils particulièrement exposés ?}

Les élèves de ton âge (14-15 ans) sont une population particulièrement fragile face aux excès du numérique, pour plusieurs raisons :

\begin{itemize}
\item un \textbf{temps d'écran élevé} : l'adolescence est la période où l'on passe le plus de temps sur les écrans, souvent sans surveillance ;
\item la \textbf{curiosité} naturelle : envie de tout essayer, de tout voir, y compris des contenus inadaptés ;
\item l'\textbf{influence du groupe} : peur d'être exclu si l'on n'est pas sur le même réseau social ou le même jeu que les copains ;
\item un \textbf{sommeil encore fragile} : le cerveau de l'adolescent a besoin de plus de sommeil que celui de l'adulte, et les écrans du soir le perturbent fortement ;
\item une \textbf{maîtrise de soi en construction} : à cet âge, il est plus difficile de s'arrêter de soi-même (« encore une vidéo, encore une partie... »).
\end{itemize}

\courssub{Usage raisonné ou usage excessif : comment faire la différence ?}

\begin{definition}[Usage raisonné et usage excessif]
Un \cle{usage raisonné} est une utilisation des outils numériques qui reste compatible avec les études, le sommeil, la santé et la vie sociale. Un \cle{usage excessif} est une utilisation dont la durée ou le moment nuit à l'une de ces activités essentielles.
\end{definition}

Pour distinguer les deux, on retient trois critères simples :

\begin{center}
\begin{tabularx}{\linewidth}{|l|X|X|}
\hline
\rowcolor{popblueL} \textbf{Critère} & \textbf{Usage raisonné} & \textbf{Usage excessif} \\
\hline
Durée quotidienne de loisir & Environ 2 h au maximum & 4 h, 5 h et plus, sans pouvoir s'arrêter \\
\hline
Moment de la journée & Après les devoirs, jamais pendant la nuit & Tard le soir, la nuit, en classe ou à table \\
\hline
Impact sur la vie & Notes stables, sommeil suffisant, relations préservées & Notes en baisse, fatigue, disputes, isolement \\
\hline
\end{tabularx}
\end{center}

Appliquons ces critères à Aïcha : plus de 4 h par jour (critère 1 dépassé), le soir jusqu'à tard (critère 2 dépassé), notes en baisse, mauvais sommeil et dispute avec une amie (critère 3 dépassé). Conclusion : son usage est \textbf{excessif} sur les trois critères.

\begin{exemplebox}[Un exemple chiffré qui fait réfléchir]
Un élève qui passe 5 h par jour sur son téléphone y consacre :
\[ 5 \times 365 = 1\,825 \text{ heures par an, soit plus de } 1\,800 \text{ h.} \]
Une journée de classe dure environ 8 h. Ces 1\,825 h représentent donc :
\[ \frac{1\,825}{8} \approx 228 \text{ journées de classe,} \]
c'est-à-dire \textbf{plus de 200 journées d'école} passées chaque année... devant un écran de loisir ! De quoi réviser toutes ses matières, faire du sport et voir ses amis.
\end{exemplebox}

Le graphique suivant compare le temps d'écran de loisir recommandé pour un adolescent et le temps moyen réellement observé :

\begin{popfigure}
\begin{center}
\begin{tikzpicture}
\begin{axis}[
  ybar, bar width=1.6cm,
  width=0.85\linewidth, height=6.5cm,
  ymin=0, ymax=7,
  ylabel={Temps d'écran de loisir (heures par jour)},
  symbolic x coords={Recommandé, Observé chez un adolescent},
  xtick=data,
  nodes near coords,
  every node near coord/.append style={font=\bfseries},
  ytick={0,1,2,3,4,5,6},
  axis lines=left,
  enlarge x limits=0.45
]
\addplot[fill=popgreen!80, draw=popgreen] coordinates {(Recommandé,2)};
\addplot[fill=poporange!85, draw=poporange] coordinates {(Observé chez un adolescent,5)};
\end{axis}
\end{tikzpicture}
\end{center}
\legende{Temps d'écran de loisir quotidien : durée recommandée ($\leq$ 2 h) contre durée moyenne observée chez un adolescent (4 à 6 h, ici 5 h).}
\end{popfigure}

\begin{exoresolu}[Analyser une situation d'usage]
Énoncé. Martial, élève de 3ème, consulte son smartphone 1 h le matin avant les cours, 2 h à midi et 3 h le soir jusqu'à 23 h 30. Il se lève à 5 h 30. Ses professeurs signalent qu'il s'endort en classe. (1) Calcule son temps d'écran quotidien. (2) Calcule sa durée de sommeil. (3) Son usage est-il raisonné ou excessif ? Justifie avec les trois critères.
\tcblower
Solution détaillée.
\begin{enumerate}
\item Temps d'écran : $1 + 2 + 3 = 6$ h par jour.
\item Sommeil : il se couche vers 23 h 30 (au plus tôt) et se lève à 5 h 30, soit environ $6$ h de sommeil, alors qu'un adolescent a besoin de 8 à 9 h.
\item Analyse par critère :
\begin{itemize}
\item \textbf{Durée} : 6 h de loisir, très au-dessus des 2 h recommandées $\rightarrow$ excessif ;
\item \textbf{Moment} : écran jusqu'à 23 h 30 et dès le matin avant les cours $\rightarrow$ excessif ;
\item \textbf{Impact} : sommeil insuffisant, endormissement en classe, donc études et santé atteintes $\rightarrow$ excessif.
\end{itemize}
\end{enumerate}
\textbf{Conclusion justifiée :} l'usage de Martial est excessif sur les trois critères. Une bonne hygiène du numérique consisterait à réduire le temps d'écran à environ 2 h de loisir, à éteindre le téléphone au moins 30 minutes avant le coucher et à viser 8 à 9 h de sommeil.
\end{exoresolu}

\begin{exercice}[À toi de jouer]
Relève ton propre temps d'écran d'hier (téléphone, télévision, jeux). Compare-le aux 2 h recommandées, puis rédige en quelques lignes une conclusion justifiée : ton usage est-il raisonné ou excessif ? Appuie ta réponse sur les trois critères (durée, moment, impact).
\end{exercice}

\begin{corrige}[Exercice : À toi de jouer]
Éléments de correction (la réponse dépend de ton relevé personnel, mais la démarche attendue est la suivante) :
\begin{enumerate}
\item \textbf{Durée} : additionne les durées passées sur chaque écran de loisir (téléphone, télévision, jeux). Si le total dépasse environ 2 h, le critère de durée est dépassé.
\item \textbf{Moment} : vérifie si tu as utilisé un écran tard le soir, pendant la nuit, en classe ou pendant les repas. Si oui, le critère du moment est dépassé.
\item \textbf{Impact} : demande-toi si cet usage a nui à ton sommeil, à tes devoirs, à tes notes ou à tes relations (fatigue, dispute, devoir non fait). Si oui, le critère d'impact est dépassé.
\end{enumerate}
\textbf{Conclusion attendue} : rédige une phrase du type « Mon usage est raisonné car... » ou « Mon usage est excessif car... », en citant les critères concernés et en proposant au moins une bonne pratique pour t'améliorer (par exemple : éteindre le téléphone 30 minutes avant le coucher).
\end{corrige}

\begin{aretenir}[L'essentiel de la section]
L'hygiène du numérique est l'ensemble des bonnes pratiques qui permettent d'utiliser les outils numériques sans nuire à sa santé physique, à son équilibre psychologique ni à son bien-être social. Le numérique est un outil précieux pour apprendre, communiquer, s'informer et se divertir, à condition d'en faire un \textbf{usage raisonné} : durée limitée (environ 2 h de loisir par jour), bons moments (jamais la nuit), et aucun impact négatif sur les études, le sommeil et les relations. Les adolescents, très exposés, doivent être particulièrement vigilants.
\end{aretenir}

\coursec{Notion d'hygiène du numérique}

L'informatique n'est pas seulement une affaire de code, de machines ou de réseaux : c'est d'abord une activité \emph{humaine} qui engage ton corps, ton esprit et tes relations aux autres. En classe de 3ème, tu passes déjà plusieurs heures par jour devant des écrans — pour tes devoirs, tes révisions, tes loisirs, ta communication. Sans repères, cette immersion numérique peut devenir un facteur de fatigue, de stress, voire de pathologies réelles. L'\cle{hygiène du numérique} est l'ensemble des connaissances, habitudes et aménagements qui permettent d'utiliser les technologies de l'information et de la communication (TIC) en préservant ta santé physique, ton équilibre psychologique et ta vie sociale. C'est une compétence transversale, au même titre que la sécurité informatique ou l'algorithmique : elle conditionne ta capacité à apprendre, à travailler et à t'épanouir durablement dans un monde numérisé.

\begin{definition}[Hygiène du numérique]
L'\cle{hygiène du numérique} désigne l'ensemble des pratiques préventives et correctives visant à réduire les effets néfastes d'une utilisation prolongée ou inadaptée des équipements numériques (ordinateurs, smartphones, tablettes, consoles) sur la santé physique (vision, sommeil, posture, audition), la santé psychologique (attention, anxiété, dépendance, estime de soi) et le bien-être social (relations familiales, scolaires, citoyenneté numérique).
\end{definition}

\begin{aretenir}[Pourquoi l'hygiène numérique est une priorité en 3ème]
\begin{itemize}
    \item \textbf{Période charnière} : à 14-15 ans, le corps (croissance, puberté), le cerveau (myélinisation, remodelage synaptique) et l'identité sociale sont en pleine construction. Les habitudes prises maintenant s'ancrent pour l'âge adulte.
    \item \textbf{Exposition massive} : cours sur ordinateur, devoirs en ligne, révisions sur smartphone, réseaux sociaux, jeux vidéo : le temps d'écran quotidien dépasse souvent 6-8 heures hors temps scolaire.
    \item \textbf{Enjeu scolaire direct} : fatigue visuelle $\rightarrow$ baisse de concentration $\rightarrow$ erreurs bêtes $\rightarrow$ notes en baisse. Troubles du sommeil $\rightarrow$ mémoire de travail défaillante $\rightarrow$ difficultés en maths, physique, langues.
    \item \textbf{Compétence BEPC} : « Appliquer les notions de l'unité à des situations concrètes de la vie courante » et « Rédiger et justifier une démarche ». L'hygiène numérique est le terrain d'application idéal : tu observes, tu analyses, tu proposes, tu justifies.
\end{itemize}
\end{aretenir}

\coursec{Risques physiques liés aux écrans}

Après avoir posé le cadre, penchons-nous sur les conséquences concrètes, mesurables et souvent sous-estimées d'une utilisation prolongée et mal maîtrisée des écrans sur ton corps. Tu as sans doute déjà ressenti une gêne oculaire après une longue session de devoirs sur ordinateur, une raideur dans le cou après avoir « scrollé » sur ton téléphone au lit, ou une difficulté à t'endormir après une partie de jeu vidéo tardive. Ces signaux ne sont pas anodins : ce sont des messages que ton corps t'envoie pour te dire que les conditions d'utilisation ne sont pas adaptées. Dans cette section, nous allons identifier, nommer et comprendre ces risques physiques pour mieux les prévenir.

\courssub{La fatigue visuelle numérique (syndrome de la vision informatique)}

Lorsque tu fixes un écran pendant une longue période, tes yeux travaillent dans des conditions très différentes de celles pour lesquelles ils sont faits naturellement : la vision de loin, en lumière naturelle, avec des clignements fréquents. Devant un écran, la distance de mise au point est fixe (généralement 50 à 70 cm), la luminosité est artificielle, le contraste peut être mal réglé, et surtout, tu clignes des yeux beaucoup moins souvent — environ 5 à 7 fois par minute au lieu de 15 à 20. Résultat : le film lacrymal qui protège la surface de l'œil s'évapore plus vite qu'il ne se renouvelle.

\begin{definition}[Fatigue visuelle numérique]
La \cle{fatigue visuelle numérique} (aussi appelée \emph{syndrome de la vision informatique}) est un ensemble de troubles oculaires et extra-oculaires provoqués par une exposition prolongée aux écrans. Elle se manifeste par des yeux secs, des picotements, une sensation de sable dans les yeux, une vision floue transitoire, des maux de tête (souvent frontaux), une sensibilité accrue à la lumière et parfois une douleur derrière les yeux.
\end{definition}

\begin{attention}[Erreur fréquente : utiliser le téléphone dans le noir complet]
Beaucoup d'élèves consultent leur téléphone au lit, lumière éteinte, en pensant que la luminosité de l'écran suffit. C'est une \textbf{erreur majeure} : le fort contraste entre l'écran brillant et l'obscurité environnante force la pupille à se dilater et à se contracter en permanence, ce qui fatigue le muscle ciliaire et la rétine bien davantage qu'un éclairage d'ambiance doux. \textbf{Allume toujours une veilleuse ou une lampe de chevet à intensité faible} quand tu utilises un écran dans la pénombre.
\end{attention}

\courssub{Les troubles du sommeil : l'effet de la lumière bleue}

C'est sans doute le risque le plus insidieux, car ses effets ne se voient pas immédiatement mais s'accumulent nuit après nuit. Les écrans (smartphones, tablettes, ordinateurs, téléviseurs) émettent une proportion importante de \cle{lumière bleue} (longueur d'onde autour de 460-480 nm). Or, nos yeux possèdent des photorécepteurs spécifiques (cellules ganglionnaires à mélanopsine) qui sont particulièrement sensibles à cette longueur d'onde et qui envoient un signal direct à l'horloge biologique (noyau suprachiasmatique de l'hypothalamus) : « il fait jour, reste éveillé ».

\begin{savaistu}[La lumière bleue trompe le cerveau]
La lumière bleue des écrans trompe le cerveau qui « croit » qu'il fait encore jour et \textbf{retarde la production de mélatonine}, l'hormone qui déclenche l'endormissement. Une exposition de 30 minutes à un écran lumineux le soir peut décaler l'horloge interne de 30 à 60 minutes. Chez l'adolescent, dont le rythme circadien est déjà naturellement décalé vers le soir (chronotype « soir »), cet effet est amplifié et peut conduire à un véritable \emph{décalage horaire social} : on s'endort tard, on se réveille fatigué, on compense par des siestes ou du café, et le cycle s'aggrave.
\end{savaistu}

\begin{popfigure}
\begin{tikzpicture}[scale=0.9]
\begin{axis}[
    width=\linewidth,
    height=9cm,
    xlabel={Heure du soir},
    ylabel={Qualité du sommeil (échelle 0-100)},
    xmin=21, xmax=25,
    xtick={21,22,23,24,25},
    xticklabels={21h, 22h, 23h, 0h, 1h},
    ymin=0, ymax=100,
    ytick={0,20,40,60,80,100},
    grid=major,
    grid style={dashed, gray!40},
    axis lines=middle,
    enlargelimits=0.05,
    legend pos=south west,
    legend style={fill=white, draw=popdark, font=\small},
    title style={font=\bfseries, align=center},
    title={Baisse de la qualité du sommeil selon le temps d'écran le soir}
]
% Courbe sans écran (référence)
\addplot[popblue, thick, smooth, mark=*] coordinates {
    (21,95) (22,90) (23,85) (24,80) (25,75)
};
\addlegendentry{Sans écran (référence)}

% Courbe 30 min écran
\addplot[poporange, thick, smooth, mark=square*] coordinates {
    (21,90) (22,80) (23,65) (24,50) (25,40)
};
\addlegendentry{30 min d'écran}

% Courbe 1h écran
\addplot[popred, thick, smooth, mark=triangle*] coordinates {
    (21,85) (22,70) (23,50) (24,30) (25,15)
};
\addlegendentry{1h d'écran}

% Courbe 2h écran
\addplot[poppurple, thick, smooth, mark=diamond*] coordinates {
    (21,80) (22,55) (23,30) (24,15) (25,5)
};
\addlegendentry{2h d'écran}

% Zone critique
\fill[popred!15] (21,0) rectangle (25,40);
\node[popred, font=\small\bfseries] at (23,20) {Zone de sommeil très dégradé};
\end{axis}
\end{tikzpicture}
\legende{Évolution simulée de la qualité du sommeil perçue en fonction de l'heure de coucher et de la durée d'exposition aux écrans le soir. Les données sont indicatives, basées sur des méta-analyses de la littérature sur la lumière bleue et la mélatonine.}
\end{popfigure}

\courssub{Les douleurs musculaires et les mauvaises postures}

Le corps humain est fait pour bouger, pas pour rester statique dans une position voûtée pendant des heures. Quand tu travailles sur un ordinateur mal réglé (écran trop bas, chaise trop haute, pas de support lombaire) ou que tu utilises ton smartphone en marchant ou assis sur ton lit, tu imposes à tes muscles, tendons et articulations des contraintes anormales.

\begin{exemplebox}[Le « text neck » : quand la tête pèse 27 kg]
La tête d'un adulte pèse en moyenne 4,5 à 5,5 kg. Mais plus tu penches la tête vers l'avant pour regarder ton téléphone, plus la charge effective sur les vertèbres cervicales augmente, à cause de l'effet de levier :
\begin{itemize}
    \item À 0° (tête droite) : charge = \textbf{5 kg}
    \item À 15° : charge = \textbf{12 kg}
    \item À 30° : charge = \textbf{18 kg}
    \item À 45° : charge = \textbf{22 kg}
    \item À 60° (position typique « nez sur l'écran ») : charge = \textbf{27 kg} — l'équivalent d'un enfant de 8 ans posé sur ton cou !
\end{itemize}
Cette surcharge quotidienne explique les torticolis, les cervicalgies chroniques, les maux de tête de tension et, à long terme, l'arthrose cervicale précoce.
\end{exemplebox}

\begin{popfigure}
\begin{tikzpicture}[scale=0.85, every node/.style={font=\small}]
% Corps humain simplifié
\draw[popdark, line width=1.5pt] (0,0) ellipse (1.2 and 3.5); % torse
\draw[popdark, line width=1.5pt] (0,3.5) circle (1); % tête
\draw[popdark, line width=1.5pt] (-1.2,1.5) .. controls (-2.5,1) and (-2.5,-1) .. (-1.2,-1); % bras gauche
\draw[popdark, line width=1.5pt] (1.2,1.5) .. controls (2.5,1) and (2.5,-1) .. (1.2,-1); % bras droit
\draw[popdark, line width=1.5pt] (-0.5,-3.5) -- (-0.5,-6.5); % jambe gauche
\draw[popdark, line width=1.5pt] (0.5,-3.5) -- (0.5,-6.5); % jambe droite

% Flèches et légendes pour les zones à risque
% Yeux
\draw[popred, -{Stealth[length=3mm]}, thick] (0.8,4.2) -- (2.5,5);
\node[popred, align=left, anchor=west] at (2.6,5) {\textbf{Yeux} : sécheresse,\\ fatigue, maux de tête};

% Cerveau / Sommeil
\draw[poppurple, -{Stealth[length=3mm]}, thick] (-0.8,4.2) -- (-3,5);
\node[poppurple, align=right, anchor=east] at (-3.1,5) {\textbf{Cerveau} : retard\\ mélatonine, insomnie};

% Cou (text neck)
\draw[poporange, -{Stealth[length=3mm]}, thick] (-0.2,2.8) -- (-2.5,2.5);
\node[poporange, align=right, anchor=east] at (-2.6,2.5) {\textbf{Cou} : text neck,\\ 27 kg à 60°};

% Dos
\draw[popgreen, -{Stealth[length=3mm]}, thick] (-0.2,0) -- (-2.5,0);
\node[popgreen, align=right, anchor=east] at (-2.6,0) {\textbf{Dos} : lombalgies,\\ cyphose, scoliose};

% Poignets
\draw[popblue, -{Stealth[length=3mm]}, thick] (1.5,1) -- (3.5,1.5);
\node[popblue, align=left, anchor=west] at (3.6,1.5) {\textbf{Poignets} : syndrome\\ du canal carpien};

% Oreilles
\draw[poppink, -{Stealth[length=3mm]}, thick] (1.2,3.8) -- (3,4.5);
\node[poppink, align=left, anchor=west] at (3.1,4.5) {\textbf{Oreilles} : surdité\\ précoce, acouphènes};

% Smartphone dans la main droite (position typique)
\draw[popmuted, fill=popcream, rounded corners=2pt] (1.8,0.5) rectangle (2.8,-0.5);
\node[popmuted, font=\tiny] at (2.3,0) {Smartphone};

% Flèche indiquant la posture voûtée
\draw[popred, dashed, thick, -{Stealth[length=2mm]}] (0,3.5) .. controls (0.5,3) and (1,2) .. (1.5,0.5);
\node[popred, font=\tiny] at (1.8,2) {Tête penchée};

% Légende globale
\node[align=center, font=\small\bfseries, popdark] at (0,-7.5) {Zones du corps impactées par une mauvaise hygiène numérique};
\end{tikzpicture}
\legende{Schéma simplifié des zones du corps affectées par une utilisation prolongée et mal posturée des écrans. Chaque flèche colorée indique une zone à risque et les pathologies associées.}
\end{popfigure}

\courssub{La sédentarité : l'ennemi invisible}

Passer 4, 6, 8 heures par jour assis devant un écran (cours, devoirs, jeux, réseaux sociaux) réduit drastiquement le temps d'activité physique. L'OMS recommande aux adolescents de 11 à 17 ans \textbf{au moins 60 minutes d'activité physique modérée à intense par jour}. La réalité dans beaucoup de collèges et lycées camerounais (et ailleurs) est bien loin de cet objectif.

\begin{aretenir}[Conséquences de la sédentarité numérique]
\begin{itemize}
    \item \textbf{Surpoids et obésité} : baisse de la dépense énergétique, grignotage devant l'écran, perturbation des signaux de satiété.
    \item \textbf{Baisse de la condition physique} : essoufflement à l'effort, perte de force musculaire, raideur articulaire.
    \item \textbf{Risques métaboliques} : résistance à l'insuline, hausse du cholestérol, hypertension artérielle — des pathologies qu'on voyait autrefois chez l'adulte et qui apparaissent maintenant chez l'adolescent.
    \item \textbf{Impact cognitif} : l'activité physique oxygène le cerveau, stimule la neurogenèse (création de neurones) et améliore la mémoire et l'attention. La sédentarité fait l'effet inverse.
\end{itemize}
\end{aretenir}

\courssub{Les risques auditifs : le danger des écouteurs à volume élevé}

C'est un risque silencieux, indolore au début, mais \textbf{irréversible}. L'oreille interne contient des cellules ciliées (environ 15 000 par oreille) qui transforment les vibrations sonores en signaux nerveux. Ces cellules \textbf{ne se régénèrent pas}. Une exposition prolongée à un son supérieur à 85 dB les détruit progressivement.

\begin{propriete}[Seuils de dangerosité sonore (OMS / INRS)]
\begin{center}
\begin{tabularx}{\linewidth}{|l|X|X|}
\hline
\rowcolor{popblueL}
\textbf{Niveau sonore} & \textbf{Exemple} & \textbf{Durée maximale sans protection} \\
\hline
80 dB & Trafic routier, restaurant bruyant & 8 heures \\
\hline
85 dB & Tondeuse à gazon, métro & 8 heures (seuil de risque) \\
\hline
90 dB & Moto, concert (loin de la scène) & 2 heures \\
\hline
95 dB & Casque à volume moyen-élevé & 45 minutes \\
\hline
100 dB & Concert, boîte de nuit, casque à fond & 15 minutes \\
\hline
110 dB & Casque à volume maximum (smartphone) & \textbf{1 à 2 minutes} \\
\hline
\end{tabularx}
\end{center}
\end{propriete}

\begin{attention}[Piège : le volume « qui va bien » dans le bruit]
Dans un environnement bruyant (bus, rue, cantine), tu augmentes le volume pour couvrir le bruit ambiant. Tu atteins facilement 100-110 dB sans t'en rendre compte, car l'oreille s'adapte (phénomène d'accoutumance). \textbf{Règle d'or} : si ton voisin entend ta musique through tes écouteurs, le volume est trop fort. Utilise des écouteurs à réduction de bruit active ou, mieux, \textbf{pose ton téléphone et écoute le monde réel}.
\end{attention}

\courssub{Autres risques : ondes, accidents et précautions d'usage}

\begin{itemize}
    \item \textbf{Ondes électromagnétiques} : les smartphones émettent des ondes radiofréquences (0,7 à 2,6 GHz selon la 2G/3G/4G/5G). Le \cle{DAS} (Débit d'Absorption Spécifique) mesure l'énergie absorbée par le corps (en W/kg). La limite réglementaire en Europe et au Cameroun est de 2 W/kg pour la tête et le tronc, 4 W/kg pour les membres. \textbf{Principe de précaution} : ne pas dormir avec le téléphone sous l'oreiller ou sur la table de nuit en mode avion désactivé ; utiliser le kit mains-libres ou le haut-parleur pour les longs appels ; privilégier les SMS ou la messagerie instantanée (moins d'émission continue).
    \item \textbf{Accidents de la voie publique} : marcher en regardant son téléphone (\"smombie\" = smartphone zombie) réduit le champ visuel de 95\% et allonge le temps de réaction. Traverser la route avec des écouteurs coupe l'alerte auditive (klaxon, moteur). Au Cameroun, où la cohabitation piétons/voitures/motos est dense, c'est une cause majeure d'accidents graves chez les jeunes.
    \item \textbf{Brûlures et incendies} : charger son téléphone sous l'oreiller ou sur un matelas empêche la dissipation de chaleur. Des cas d'incendies de literie par surchauffe de batterie (surtout contrefaite ou défectueuse) sont recensés chaque année.
\end{itemize}

\coursec{Risques psychologiques et sociaux}

Au-delà du corps, le numérique façonne ton esprit et tes relations. Les concepteurs d'applications (réseaux sociaux, jeux vidéo, plateformes de streaming) utilisent des mécanismes psychologiques puissants pour capter et retenir ton attention le plus longtemps possible : c'est l'\cle{économie de l'attention}. Comprendre ces mécanismes, c'est reprendre le pouvoir.

\courssub{La dépendance numérique et le circuit de la récompense}

Les notifications, les « likes », les récompenses aléatoires dans les jeux (loot boxes), le défilement infini (infinite scroll) de TikTok ou Instagram activent le \textbf{circuit de la récompense} dans ton cerveau (libération de dopamine). Ce n'est pas un hasard : c'est du \emph{design persuasif} (persuasive design), inspiré des machines à sous (renforcement à ratio variable). Le cerveau associe l'écran à une source rapide de plaisir, créant une envie irrépressible (craving) d'y revenir.

\begin{definition}[Nomophobie]
La \cle{nomophobie} (contraction de « no mobile phone phobia ») est la peur irrationnelle d'être séparé de son téléphone portable, de ne pas avoir de réseau ou de batterie. Elle se manifeste par de l'anxiété, des palpitations, une vérification compulsive de l'écran (checking behavior) toutes les quelques minutes, même sans notification.
\end{definition}

\begin{attention}[Le piège du « juste 5 minutes »]
Tu te dis « je regarde 5 minutes » et tu émerges 2 heures plus tard. C'est l'effet de l'\emph{autoplay} (lecture automatique), de l'absence de signaux d'arrêt (pas de fin de page, pas de générique) et de la personnalisation algorithmique qui te sert exactement ce qui te retient. \textbf{Solution technique} : désactive l'autoplay, active les limites de temps d'écran (Bien-être numérique / Screen Time), utilise une alarme physique.
\end{attention}

\courssub{FOMO, anxiété et estime de soi}

\cle{FOMO} (Fear Of Missing Out — peur de rater quelque chose) : l'affichage permanent de la vie « parfaite » des autres (voyages, soirées, notes, corps retouchés) crée une comparaison sociale ascendante toxique. Tu as l'impression que ta vie est moins intéressante, que tu rates des événements, que tu n'es pas « assez bien ». Cela alimente l'anxiété, la déprime, et fragilise l'estime de soi, déjà instable à l'adolescence.

\begin{savaistu}[Le paradoxe de la connexion]
Plus on est « connecté » numériquement (des centaines d'amis, des groupes WhatsApp), plus on peut se sentir \emph{seul} réellement. La communication médiatisée (texte, émojis) appauvrit les signaux non verbaux (regard, ton, toucher) qui construisent l'attachement et l'empathie. Passer 3h sur WhatsApp ne remplace pas 30 min de discussion en face à face avec un ami ou un parent.
\end{savaistu}

\courssub{Cyberharcèlement et violence en ligne}

Le cyberharcèlement est une forme de harcèlement perpétrée via les outils numériques (réseaux sociaux, messageries, jeux en ligne, forums). Ses spécificités : l'anonymat possible de l'agresseur, la viralité (public large, immédiat), la permanence (traces difficiles à effacer), l'intrusion dans le domicile (plus de refuge), l'effet de meute (pile-on).

\begin{propriete}[Formes principales de cyberharcèlement]
\begin{itemize}
    \item \textbf{Harcèlement direct} : insultes, menaces, rumeurs répétées par messages privés ou publics.
    \item \textbf{Usurpation d'identité} : création de faux profils pour nuire à la réputation.
    \item \textbf{Outing / Divulgation} : publication non consentie de photos, vidéos, secrets intimes (revenge porn, sextorsion).
    \item \textbf{Exclusion délibérée} : bannissement de groupes, « unfollow » massif coordonné, ignorance systématique.
    \item \textbf{Cyberharcèlement sexiste / homophobe / raciste} : attaques ciblées sur l'identité.
\end{itemize}
\end{propriete}

\begin{methode}[Réagir face au cyberharcèlement (règle des 3 N : Ne pas répondre, Notifier, Numériser)]
\begin{enumerate}
    \item \textbf{Ne pas répondre} : la réponse alimente l'agresseur. Bloque le contact.
    \item \textbf{Numériser (Preuves)} : captures d'écran (date, heure, pseudo, URL), enregistrements. Indispensables pour la plainte.
    \item \textbf{Notifier} : signaler sur la plateforme (bouton « Signaler »), en parler à un adulte de confiance (parent, professeur, CPE, infirmier), appeler le \textbf{3011} (numéro vert national au Cameroun pour la protection de l'enfance) ou contacter l'ANTIC (Agence Nationale des Technologies de l'Information et de la Communication).
    \item \textbf{Se protéger} : vérifier ses paramètres de confidentialité (profil privé, contrôle des tags, restriction des commentaires).
\end{enumerate}
\end{methode}

\courssub{Phubbing et isolement relationnel}

Le \cle{phubbing} (phone snubbing) : snobiner une personne physiquement présente pour regarder son téléphone. C'est devenu une norme sociale, mais cela dégrade la qualité des interactions : sentiment d'être ignoré, baisse de l'empathie, conflits familiaux (repas, sorties), rupture du lien d'attachement parents-enfants.

\begin{exemplebox}[Le repas de famille « connecté »]
Une famille de 4 personnes au restaurant. Chacun a son smartphone sur la table. Les parents vérifient leurs mails professionnels, l'ado scroll TikTok, le cadet regarde des dessins animés. Personne ne se parle. Le repas dure 20 minutes. \textbf{Bilan} : 0 minute d'échange réel. \textbf{Alternative} : « Règle du téléphone au centre de la table » (premier qui le touche paie l'addition ou fait la vaisselle) ou « Zone sans écran / Heure sans écran » décidée collectivement.
\end{exemplebox}

\coursec{Les bonnes pratiques d'hygiène numérique}

Maintenant que les risques sont identifiés, passons à l'action. L'hygiène numérique n'est pas une liste d'interdits, mais une \textbf{boîte à outils} pour adapter ton environnement, tes habitudes et tes réglages à ta physiologie et à tes objectifs.

\courssub{La règle d'or : 20-20-20 et pauses actives}

\begin{propriete}[Règle 20-20-20 (Association américaine d'ophtalmologie)]
Toutes les \textbf{20 minutes}, regarde un objet situé à \textbf{20 pieds (6 mètres)} pendant \textbf{20 secondes}. Cela relâche l'accommodation (muscle ciliaire) et stimule le clignement (réhydratation).
\end{propriete}

\begin{methode}[Mise en place concrète des pauses]
\begin{enumerate}
    \item Installe une extension navigateur ou une application (ex. : \texttt{Stretchly}, \texttt{EyeCare}, \texttt{Pause} sur Android/iOS, \texttt{Workrave} sur Windows/Linux) qui verrouille l'écran ou affiche une alerte toutes les 20-25 min.
    \item Pendant la pause : \textbf{lève-toi}, étire les bras, roule les épaules, bois une gorgée d'eau, regarde par la fenêtre (loin). \textbf{Pas de téléphone} pendant la pause !
    \item Toutes les 2 heures : pause longue de 15-20 min (marche, discussion, collation saine).
\end{enumerate}
\end{methode}

\courssub{Ergonomie du poste de travail (zéro budget, efficacité max)}

\begin{center}
\begin{tabularx}{\linewidth}{|l|X|X|}
\hline
\rowcolor{popblueL}
\textbf{Élément} & \textbf{Règle} & \textbf{Astuce 0 FCFA} \\
\hline
Distance yeux-écran & 50-70 cm (bras tendu) & Reculer la chaise, avancer le clavier \\
\hline
Hauteur écran & Bord supérieur \textbf{au niveau des yeux} & Surélever l'écran/portable avec des livres, une boîte de chaussures, des dictionnaires \\
\hline
Chaise & Pieds à plat, genoux 90°, dos soutenu & Coussin lombaire (serviette roulée), cale-pieds (carton) \\
\hline
Clavier / Souris & Avant-bras parallèles au sol, poignets droits & Tapis de souris avec repose-poignet (tissu plié) \\
\hline
Éclairage & Lumière latérale, pas de reflets sur l'écran & Éteindre néon au-dessus, lampe de bureau vers clavier, store/rideau fenêtre \\
\hline
\end{tabularx}
\end{center}

\courssub{Hygiène du sommeil numérique}

\begin{itemize}
    \item \textbf{Couvre-feu numérique} : 30 à 60 min avant le coucher, \textbf{zéro écran}. Remplace par lecture papier, dessin, discussion, méditation, étirements.
    \item \textbf{Chambre = sanctuaire} : téléphone \textbf{hors de la chambre} (mode avion + chargeur dans une autre pièce). Achète un \textbf{réveil classique} (500-1000 FCFA). C'est l'investissement le plus rentable de ta scolarité.
    \item \textbf{Filtres logiciels} : active \texttt{Night Light} (Windows 10/11), \texttt{Mode Nuit} (macOS), \texttt{Protection oculaire / Mode nuit} (Android), \texttt{Night Shift} (iOS) dès le coucher du soleil (réglage automatique sur lever/coucher du soleil). Cela réduit la lumière bleue sans dégrader trop les couleurs.
    \item \textbf{Régularité} : même heure de coucher/lever week-end compris (max 1h de décalage). Le rythme circadien aime la régularité.
\end{itemize}

\courssub{Gestion du temps d'écran et attention}

\begin{itemize}
    \item \textbf{Mesure} : active « Bien-être numérique » (Android) / « Temps d'écran » (iOS) / « Activité » (Windows). Regarde les chiffres réels (souvent 2x l'estimation subjective).
    \item \textbf{Budget temps} : fixe une limite journalière par catégorie (Réseaux sociaux : 1h, Jeux : 45 min, YouTube : 1h). L'appli bloque ou alerte à la limite.
    \item \textbf{Mode monochrome} : passe l'écran en noir et blanc (Accessibilité > Filtres de couleur). L'écran devient beaucoup moins attractif, le « doomscrolling » perd de son pouvoir.
    \item \textbf{Notifications} : \textbf{désactive TOUT} sauf appels, SMS, et 1-2 apps essentielles (ex. : Pronote, messagerie parents). Pas de badge, pas de son, pas de bannière pour Instagram, TikTok, jeux, actualités.
    \item \textbf{Technique Pomodoro} : 25 min travail concentré (téléphone en mode avion, loin de toi) / 5 min pause. 4 cycles = 2h productives.
\end{itemize}

\courssub{Hygiène auditive}

\begin{itemize}
    \item \textbf{Règle 60/60} : volume max \textbf{60\%}, durée max \textbf{60 min} d'affilée, puis pause 10-15 min sans son.
    \item \textbf{Casque circum-auriculaire} (autour de l'oreille) préféré aux intra-auriculaires (dans le conduit) : meilleure isolation passive $\rightarrow$ volume plus bas pour même confort.
    \item \textbf{Réduction de bruit active (ANC)} : investissement utile si budget permet (réduit le besoin de monter le volume dans le bruit).
    \item \textbf{Silence} : accorde-toi des plages de silence total chaque jour (marche, devoirs, repas, avant sommeil). Le cerveau a besoin de traiter l'information sans flux entrant.
\end{itemize}

\begin{methode}[Auto-évaluation rapide de ton hygiène physique numérique]
Réponds honnêtement par Oui / Non / Parfois :
\begin{enumerate}
    \item Est-ce que je ressens des yeux secs ou des maux de tête après 1h d'écran ?
    \item Est-ce que je regarde mon téléphone dans les 30 minutes avant de dormir ?
    \item Est-ce que je dors avec mon téléphone à moins de 50 cm de ma tête ?
    \item Est-ce que j'ai mal au cou, au dos ou aux poignets en fin de journée ?
    \item Est-ce que j'écoute de la musique au casque à plus de 60\% du volume max ?
    \item Est-ce que je marche dans la rue en regardant mon écran ou avec des écouteurs ?
    \item Est-ce que je fais moins de 30 min d'activité physique par jour hors cours de sport ?
\end{enumerate}
\textbf{Interprétation} : 0-1 « Oui » = \textbf{Excellent} ; 2-3 = \textbf{Vigilance nécessaire} ; 4+ = \textbf{Urgent : change tes habitudes dès aujourd'hui}.
\end{methode}

\begin{experience}[TP guidé : Mesurer et améliorer ton environnement d'écran]
\textbf{Objectif} : Appliquer les règles d'ergonomie à ton poste de travail (maison ou cybercafé) et mesurer l'impact sur ton confort.
\textbf{Matériel} : Un ordinateur (ou smartphone), une règle, un minuteur, une fiche d'observation.
\textbf{Manipulation} :
\begin{enumerate}
    \item \textbf{Mesure la distance yeux-écran} : assis naturellement, mesure la distance entre tes yeux et l'écran. Note-la. Recommandation : 50-70 cm (longueur du bras tendu).
    \item \textbf{Vérifie la hauteur de l'écran} : le bord supérieur de l'écran doit être \textbf{au niveau des yeux} ou légèrement en dessous. Si l'écran est trop bas (cas fréquent des portables), surélève-le avec des livres ou un support.
    \item \textbf{Teste la règle 20-20-20} : règle un minuteur sur 20 minutes. Toutes les 20 min, regarde un objet à \textbf{20 pieds (6 m)} pendant \textbf{20 secondes}. Note ta sensation oculaire avant/après 1h de test.
    \item \textbf{Éclairage} : éteins la lumière principale, allume une lampe de bureau orientée vers le clavier (pas vers l'écran). Compare la fatigue visuelle après 30 min.
    \item \textbf{Posture} : pieds à plat au sol, genoux à 90°, dos droit soutenu par le dossier, avant-bras parallèles au sol, poignets droits (pas pliés vers le haut).
\end{enumerate}
\textbf{Observations attendues} : Après 1h avec ces réglages, tu devrais noter une diminution de la fatigue oculaire, moins de tensions cervicales, une meilleure concentration.
\textbf{Interprétation} : L'ergonomie n'est pas du luxe, c'est de la \textbf{prévention primaire}. Un poste bien réglé coûte 0 FCFA (quelques livres, un coussin) et t'évite des douleurs chroniques.
\end{experience}

\begin{exoresolu}[Analyse d'une situation réelle : Le cas de Kevin, élève en 3ème]
\textbf{Énoncé} : Kevin, 14 ans, passe en moyenne 5h/jour sur son smartphone (TikTok, WhatsApp, jeux) et 2h sur l'ordinateur familial pour les devoirs. Il dort avec son téléphone sous l'oreiller « pour l'alarme ». Il se plaint de maux de tête quotidiens, de difficultés à s'endormir avant minuit, de douleurs au cou et d'une baisse de ses notes au 2ème trimestre. Ses parents pensent qu'il « manque de sérieux ». Analyse la situation et propose un plan d'action concret sur 4 semaines.

\tcblower
\textbf{Solution détaillée} :

\textbf{1. Diagnostic des risques identifiés} :
\begin{itemize}
    \item \textbf{Sommeil} : téléphone sous l'oreiller = ondes + lumière bleue au réveil nocturne + alarme stressante. Endormissement vers minuit au lieu de 21h30-22h (besoin : 9h de sommeil à 14 ans).
    \item \textbf{Vision} : 7h d'écran sans pause 20-20-20 = syndrome de vision informatique certain.
    \item \textbf{Posture} : smartphone tenu bas = text neck (27 kg sur cervicales). Ordinateur familial probablement mal réglé (table de salle à manger).
    \item \textbf{Sédentarité} : 7h assises + cours + transport = quasi 0 activité physique.
    \item \textbf{Cognitif} : fragmentation de l'attention (notifications), sommeil insuffisant = baisse mémoire, concentration, humeur.
\end{itemize}

\textbf{2. Plan d'action progressif sur 4 semaines (méthode des petits pas)} :

\begin{center}
\begin{tabularx}{\linewidth}{|c|X|X|}
\hline
\rowcolor{popblueL}
\textbf{Semaine} & \textbf{Objectif principal} & \textbf{Actions concrètes} \\
\hline
1 & \textbf{Sommeil} & Acheter un réveil classique (500-1000 FCFA). Téléphone en mode avion \textbf{hors de la chambre} à 21h30. Lumière tamisée 30 min avant. Coucher 22h max. \\
\hline
2 & \textbf{Vision + Pauses} & Installer appli « Pause écran » (gratuite) : alerte 20-20-20. Lunettes anti-lumière bleue si possible (ou filtre logiciel « Night Light » / f.lux). \\
\hline
3 & \textbf{Posture + Mouvement} & Support portable (livres) pour écran à hauteur yeux. Chaise avec coussin lombaire. 15 min marche rapide / jour (aller-retour école, course pour maman). \\
\hline
4 & \textbf{Gestion du temps} & « Screen Time » / « Bien-être numérique » : limite 2h/jour réseaux sociaux + jeux. Remplacer 30 min écran par lecture, dessin, discussion famille, sport. \\
\hline
\end{tabularx}
\end{center}

\textbf{3. Indicateurs de suivi (tableau de bord hebdo)} :
\begin{itemize}
    \item Heure d'endormissement / heure de réveil / qualité sommeil (1-5)
    \item Présence maux de tête (O/N), douleurs cou/dos (1-5)
    \item Temps d'écran loisir (appli intégrée)
    \item Minutes d'activité physique volontaire
    \item Humeur / concentration en classe (auto-évaluation 1-5)
\end{itemize}

\textbf{4. Justification pédagogique} : Ce plan ne vise pas l'interdiction (contre-productive) mais la \textbf{reprise de contrôle}. Chaque semaine ajoute une brique sans tout changer d'un coup. Les parents passent du rôle de « policiers » à « alliés » (achat réveil, marche ensemble, écoute). Au bout d'un mois, Kevin a retrouvé du sommeil, moins de douleurs, et ses notes remontent car sa mémoire de travail fonctionne à nouveau. C'est ça, l'hygiène numérique : \textbf{un investissement sur soi-même}.
\end{exoresolu}

\coursec{Usage responsable des réseaux sociaux et bien-être social}

Les réseaux sociaux (WhatsApp, Facebook, Instagram, TikTok, Snapchat, X/Twitter) sont devenus la place du village numérique. Ils offent des opportunités (information, entraide, expression, citoyenneté) mais posent des défis éthiques, juridiques et psychologiques majeurs. En 3ème, tu construis ton identité numérique : elle te suivra toute ta vie.

\courssub{Empreinte numérique et vie privée : ce que tu publies ne s'efface pas}

\begin{definition}[Empreinte numérique]
L'\cle{empreinte numérique} (ou \emph{trace numérique}) est l'ensemble des informations que tu laisses volontairement (posts, photos, commentaires, likes) ou involontairement (métadonnées de géolocalisation, historique de navigation, cookies, adresses IP) sur Internet. Elle est \textbf{persistante, réplicable, traçable et potentiellement accessible à tous} (recruteurs, administrations, pairs, prédateurs).
\end{definition}

\begin{propriete}[Droits numériques fondamentaux (Réf. Loi camerounaise sur la cybercriminalité, RGPD inspiration)]
\begin{itemize}
    \item \textbf{Droit à l'oubli} : demander l'effacement de données te concernant (démarche auprès de la plateforme, puis ANTIC / justice si refus).
    \item \textbf{Droit à la portabilité} : récupérer tes données dans un format lisible.
    \item \textbf{Consentement} : tes données ne peuvent être collectées/traitées sans ton accord libre, éclairé, spécifique (sauf exceptions légales).
    \item \textbf{Protection des mineurs} : interdiction de profilage publicitaire ciblé pour les < 18 ans sur la plupart des plateformes responsables.
\end{itemize}
\end{propriete}

\begin{attention}[Le mythe de la suppression]
Supprimer un post ou une story ne l'efface pas des serveurs de la plateforme, ni des captures d'écran faites par d'autres, ni des caches des moteurs de recherche (Google, Bing), ni des archives (Wayback Machine). \textbf{Règle} : ne publie \textbf{jamais} ce que tu ne voudras pas voir sur la une d'un journal dans 10 ans (photo en maillot, propos colériques, blague douteuse, document d'identité).
\end{attention}

\courssub{Netiquette et communication non violente en ligne}

L'écran désinhibe : on écrit des choses qu'on n'oserait pas dire en face. La \cle{netiquette} (étiquette du net) est le code de conduite pour des échanges respectueux.

\begin{center}
\begin{tabularx}{\linewidth}{|l|X|}
\hline
\rowcolor{popblueL}
\textbf{Principes de netiquette} & \textbf{Application concrète} \\
\hline
Humanité & Derrière chaque pseudo, une personne réelle. Lis ton message avant d'envoyer : le dirais-tu en face ? \\
\hline
Clarté & Évite le langage SMS excessif, les MAJUSCULES (crier), l'ironie mal perçue. Relis-toi. \\
\hline
Respect vie privée & Ne partage pas photo/infos d'autrui sans accord. Ne transfère pas de messages privés. \\
\hline
Légalité & Pas d'insultes, menaces, incitation à la haine, diffamation, apologie de crimes, contenus pédopornographiques (sanctions pénales lourdes au Cameroun : Loi 2010/012). \\
\hline
Source & Vérifie avant de partager (Fake news). Cite la source. Signale les contenus illicites. \\
\hline
\end{tabularx}
\end{center}

\courssub{Gestion de l'image de soi et pression sociale}

\begin{itemize}
    \item \textbf{Profil public vs privé} : par défaut, mets ton profil en \textbf{privé} (amis abonnés validés). Sépare « amis proches » (story privée) et « connaissances » (profil public léger).
    \item \textbf{Authenticité} : ne te compare pas aux « highlights » (meilleurs moments) retouchés des influenceurs. La vraie vie a des boutons, des échecs, des matins difficiles.
    \item \textbf{Validation externe} : ne lie pas ton estime de toi au nombre de likes/vues. Ton identité ne se réduit pas à ton profil Instagram.
    \item \textbf{Sexting / Nudes} : \textbf{Jamais}. Risque de sextorsion (chantage), diffusion virale, poursuites pénales (détention d'images pédopornographiques si < 18 ans, même consenties). Si tu reçois : \textbf{supprime immédiatement, ne transfère jamais, signale, en parle à un adulte}.
\end{itemize}

\courssub{Citoyenneté numérique : de consommateur à acteur responsable}

\begin{itemize}
    \item \textbf{Signalement} : utilise les outils de signalement des plateformes (harcèlement, haine, désinformation, contenu illégal). C'est un devoir civique.
    \item \textbf{Esprit critique} : vérifie les sources (site officiel, média reconnu, date, auteur), méfie-toi des titres racoleurs (clickbait), des images sorties de contexte, des deepfakes (IA). Outils : \texttt{Africa Check}, \texttt{Dubawa}, \texttt{Google Lens} (recherche image inversée).
    \item \textbf{Engagement positif} : utilise le numérique pour t'informer (MOOCs, Khan Academy, Wikipédia), t'engager (pétitions, associations, veille citoyenne), créer (code, blog, art, vidéo éducative), t'entraider (groupes de révision, tutorat pair-à-pair).
    \item \textbf{Déconnexion choisie} : instaure des « temps morts » sacrés : repas, prière/méditation, devoirs profonds, sommeil, moments famille/amis. La \textbf{disponibilité permanente n'est pas une vertu}.
\end{itemize}

\coursec{Élaborer et justifier une démarche d'hygiène numérique}

Cette dernière section met en œuvre la compétence « Rédiger et justifier une démarche ». Tu dois être capable d'analyser une situation réelle, d'identifier les risques multiples (physiques, psychologiques, sociaux), de proposer un plan d'action étayé, chiffré, progressif, et de définir des indicateurs de suivi. C'est exactement ce que tu as fait avec le cas de Kevin. Maintenant, à toi de jouer sur un cas professionnel.

\begin{exercice}[Mise en situation : Le cybercafé de ton quartier]
Tu gères un petit cybercafé à Yaoundé (quartier Mvog-Ada). Tes clients sont majoritairement des élèves et étudiants qui y passent 2 à 6 heures d'affilée. Le mobilier est basique : chaises en plastique, tables fixes, écrans posés directement sur la table, éclairage néon au plafond, climatisation bruyante. Tu veux améliorer l'hygiène numérique de tes clients \textbf{sans gros budget}. Propose 5 aménagements concrets, peu coûteux, avec pour chacun : le risque ciblé, la solution, le coût estimé, l'argument pour convaincre le gérant (ton patron).
\end{exercice}

\begin{corrige}[Exercice : Le cybercafé de ton quartier]
\begin{enumerate}
    \item \textbf{Risque : Fatigue visuelle (éblouissement néon)}. \textbf{Solution} : Éteindre 1 néon sur 2 + poser des stores ou rideaux légers aux fenêtres pour éviter reflets. \textbf{Coût} : 0 à 5 000 FCFA (rideaux récupérés). \textbf{Argument} : « Les clients restent plus longtemps si leurs yeux ne piquent pas = plus de revenus. »
    \item \textbf{Risque : Mauvaise posture (écran trop bas, chaise inadaptée)}. \textbf{Solution} : Surélever les écrans avec des blocs de béton peints / planches + coussins lombaires (mousse découpée dans du tissu). \textbf{Coût} : ~2 000 FCFA / poste. \textbf{Argument} : « Moins de douleurs = clients fidèles, bouche-à-oreille positif. »
    \item \textbf{Risque : Sédentarité / raideur}. \textbf{Solution} : Afficher une affiche « Pause 20-20-20 » + minuteur visible + espace « étirement » (1 m² libre au fond). \textbf{Coût} : 500 FCFA (affiche imprimée). \textbf{Argument} : « Image moderne, responsable, attire parents et écoles. »
    \item \textbf{Risque : Troubles du sommeil (sessions tardives)}. \textbf{Solution} : Activer « Night Light » (Windows) / « Mode nuit » (Android) sur tous les postes + affiche « Après 22h : lumière chaude activée ». \textbf{Coût} : 0 FCFA. \textbf{Argument} : « Service gratuit qui montre qu'on prend soin des élèves. »
    \item \textbf{Risque : Hygiène auditive (casques partagés, volume fort)}. \textbf{Solution} : Fournir des embouts jetables (hygiène) + limiter volume max dans Windows (Propriétés haut-parleur > Niveaux > 70\%). \textbf{Coût} : 1 000 FCFA / mois (embouts). \textbf{Argument} : « Prévention surdité = responsabilité sociale, évite plaintes parents. »
\end{enumerate}
\textbf{Bilan} : Investissement total < 15 000 FCFA pour 10 postes. Retour sur investissement : fidélisation, réputation, différenciation face aux concurrents.
\end{corrige}

\coursec{Risques psychologiques et sociaux}

Dans la section précédente, nous avons vu que les écrans peuvent fatiguer nos yeux, notre dos et notre sommeil : ce sont les risques \emph{physiques}. Mais l'ordinateur, la tablette et le téléphone n'agissent pas seulement sur notre corps. Ils agissent aussi sur notre \emph{esprit} (nos émotions, notre concentration, notre estime de nous-mêmes) et sur nos \emph{relations avec les autres}. Ce sont les risques psychologiques et sociaux, souvent moins visibles qu'un mal de dos, mais tout aussi réels. Dans cette section, nous allons apprendre à les reconnaître, à les comprendre et à adopter les bons réflexes pour s'en protéger.

\courssub{L'addiction aux écrans et aux jeux}

\textbf{Une situation pour comprendre.} Alain, élève de 3ème, pose son téléphone sur la table pendant ses devoirs. Au bout de cinq minutes, il le déverrouille « juste pour voir ». Puis encore cinq minutes après. Le soir, il joue en ligne jusqu'à 1 heure du matin et s'énerve quand sa mère lui demande d'arrêter. Petit à petit, il délaisse le football avec ses amis et ses notes baissent. Alain présente les signes classiques d'une \cle{addiction aux écrans}.

\begin{definition}[Addiction aux écrans]
L'\cle{addiction aux écrans} est une dépendance comportementale : la personne ressent un besoin irrésistible et répété d'utiliser son téléphone, sa console ou son ordinateur, même quand cela nuit à sa vie (études, sommeil, relations). Ce n'est pas l'écran lui-même qui crée la dépendance, mais l'usage excessif et incontrôlé qu'on en fait.
\end{definition}

\textbf{Les signes qui doivent alerter :}
\begin{itemize}
  \item le \textbf{besoin constant} de consulter son téléphone, même sans raison (aucune notification reçue) ;
  \item l'\textbf{irritabilité} ou l'énervement quand on ne peut pas utiliser l'appareil (panne de batterie, interdiction des parents, coupure réseau) ;
  \item l'\textbf{abandon progressif des autres activités} : sport, lecture, sorties entre amis, repas en famille ;
  \item le \textbf{mensonge} sur le temps passé devant l'écran (« je n'ai joué qu'une heure » alors qu'il s'est écoulé quatre heures) ;
  \item la \textbf{perte de la notion du temps} : on prévoit dix minutes, on y reste trois heures.
\end{itemize}

\textbf{Conséquences sur les études.} Le cerveau fatigué par des nuits trop courtes mémorise mal. L'élève qui consulte son téléphone toutes les cinq minutes ne peut pas se concentrer assez longtemps pour comprendre une leçon ou résoudre un exercice. Résultat : baisse des notes, retards dans les devoirs, découragement. L'addiction aux écrans est donc un vrai problème scolaire, pas seulement un problème de loisir.

\begin{savaistu}[Pourquoi est-ce si difficile de lâcher son téléphone ?]
Le besoin de vérifier sans cesse ses notifications active les mêmes \emph{circuits de récompense} du cerveau que certains jeux de hasard. Chaque notification est comme une petite loterie : « peut-être un message intéressant ? ». Le cerveau libère alors une substance de plaisir (la dopamine), ce qui donne envie de recommencer. Les applications sont d'ailleurs \emph{conçues} pour cela : couleurs vives, compteurs de messages, vidéos qui s'enchaînent automatiquement. Comprendre ce mécanisme, c'est déjà reprendre le contrôle.
\end{savaistu}

\courssub{Le cyberharcèlement}

\textbf{Une situation pour comprendre.} Après une dispute, des camarades créent un groupe de discussion pour se moquer de Sandrine : surnoms blessants, photo truquée, rumeurs. Chaque soir, les messages continuent. Sandrine n'ose plus venir au collège. Contrairement à une dispute dans la cour, ce harcèlement la suit \emph{partout}, même dans sa chambre, jour et nuit.

\begin{definition}[Cyberharcèlement]
Le \cle{cyberharcèlement} est un harcèlement \textbf{répété} commis par des moyens numériques : messages, appels, réseaux sociaux, photos ou vidéos. Comme tout harcèlement, il se caractérise par la \textbf{répétition} des agressions et par une relation de domination (un groupe contre un élève isolé, par exemple). Une simple moquerie une fois n'est pas du harcèlement ; des moqueries qui reviennent chaque jour en sont.
\end{definition}

\textbf{Les formes du cyberharcèlement :}
\begin{itemize}
  \item les \textbf{moqueries et insultes} répétées par messages ou en commentaires ;
  \item les \textbf{rumeurs} propagées dans des groupes de discussion ;
  \item la publication de \textbf{photos ou vidéos humiliantes}, parfois truquées ;
  \item l'\textbf{exclusion volontaire} : créer un groupe de classe en y refusant un élève, ou l'y faire entrer pour mieux l'attaquer ;
  \item l'\textbf{usurpation d'identité} : créer un faux compte au nom de la victime pour publier des propos embarrassants.
\end{itemize}

\begin{methode}[Face au cyberharcèlement : les 4 réflexes]
Si tu es victime ou témoin de cyberharcèlement, applique dans l'ordre les 4 réflexes :
\begin{enumerate}
  \item \textbf{Ne pas répondre} aux attaques : répondre nourrit le conflit et donne au harceleur ce qu'il cherche (une réaction).
  \item \textbf{Garder les preuves} : faire des captures d'écran des messages, photos et commentaires, avec les dates et les pseudos.
  \item \textbf{En parler à un adulte} de confiance : parent, enseignant, surveillant, directeur. Le harcèlement ne s'arrête presque jamais tout seul.
  \item \textbf{Bloquer et signaler} le compte du harceleur sur la plateforme, et si nécessaire porter plainte : le cyberharcèlement est puni par la loi.
\end{enumerate}
\end{methode}

\begin{attention}[Erreur fréquente : supprimer les messages]
Beaucoup de victimes suppriment immédiatement les messages de harcèlement parce qu'ils font mal. C'est une erreur : une fois supprimés, il devient très difficile de prouver les faits. \textbf{On conserve d'abord les preuves} (captures d'écran), et on ne supprime ou ne bloque qu'ensuite.
\end{attention}

\begin{popfigure}
\begin{tikzpicture}[
  node distance=0.55cm and 1.1cm,
  etape/.style={rectangle, rounded corners=3pt, draw=popblue, fill=popblueL, text width=4.6cm, align=center, font=\small, inner sep=5pt},
  victime/.style={rectangle, rounded corners=3pt, draw=poppink, fill=poppinkL, text width=4.6cm, align=center, font=\small\bfseries, inner sep=5pt},
  fin/.style={rectangle, rounded corners=3pt, draw=popgreen, fill=popgreenL, text width=4.6cm, align=center, font=\small, inner sep=5pt},
  fleche/.style={-{Stealth[length=2.5mm]}, thick, popdark}
]
\node[victime] (v) {Je suis victime ou témoin de cyberharcèlement};
\node[etape, below=of v] (r1) {\textbf{1.} Je ne réponds pas aux attaques};
\node[etape, below=of r1] (r2) {\textbf{2.} Je conserve les preuves (captures d'écran, dates, pseudos)};
\node[etape, below=of r2] (r3) {\textbf{3.} J'en parle à un adulte de confiance (parent, enseignant)};
\node[etape, below=of r3] (r4) {\textbf{4.} Je bloque et je signale le compte sur la plateforme};
\node[fin, below=of r4] (f) {La situation est prise en charge ; si nécessaire, plainte auprès des autorités};
\draw[fleche] (v) -- (r1);
\draw[fleche] (r1) -- (r2);
\draw[fleche] (r2) -- (r3);
\draw[fleche] (r3) -- (r4);
\draw[fleche] (r4) -- (f);
\end{tikzpicture}
\legende{Arbre de décision : que faire face au cyberharcèlement ? Les quatre réflexes s'appliquent dans l'ordre.}
\end{popfigure}

\courssub{L'isolement social}

\textbf{Intuition.} Les réseaux sociaux promettent de nous « connecter » aux autres. Pourtant, un élève peut avoir 500 « amis » en ligne et se sentir profondément seul. Pourquoi ? Parce que les relations humaines ont besoin de la présence réelle : le regard, la voix, les rires partagés, les jeux dans la cour.

\begin{definition}[Isolement social numérique]
L'\cle{isolement social} lié au numérique est le remplacement progressif des relations réelles (famille, camarades, voisinage) par des relations uniquement virtuelles. La personne se replie sur elle-même : elle passe ses journées seule avec son écran, perd l'habitude de parler aux autres en face à face et finit par craindre les contacts réels.
\end{definition}

\textbf{Signes et conséquences :} refus des sorties et des activités de groupe, repas pris seul avec le téléphone, difficulté à soutenir une conversation, sentiment de solitude malgré les nombreux contacts en ligne. À long terme, l'isolement fragilise la santé mentale et les compétences sociales indispensables dans la vie (travail en groupe, entretien, amitié).

\textbf{Le bon réflexe :} les outils numériques doivent \emph{servir} les relations réelles (organiser un match, préparer un exposé ensemble), pas les \emph{remplacer}. Une règle simple : pour chaque heure d'écran de loisir, prévoir du temps avec des personnes en vrai.

\courssub{La comparaison sociale et l'estime de soi}

\textbf{Une situation pour comprendre.} En faisant défiler les publications, Nadia voit des photos de camarades toujours bien habillés, en voyage, entourés d'amis, avec des téléphones coûteux. Elle regarde sa propre vie et se sent « nulle », « pauvre », « moins intéressante ». Pourtant, elle ne voit que ce que les autres \emph{choisissent} de montrer.

Sur les réseaux sociaux, chacun publie ses meilleurs moments : la photo réussie sur trente essais, la fête et non les devoirs, le sourire et non les soucis. On obtient une \cle{image idéalisée} de la vie des autres. Comparer sa vie réelle, avec ses difficultés normales, à ces vitrines parfaites produit un \textbf{sentiment d'infériorité}, de l'\textbf{anxiété} et une baisse de l'\cle{estime de soi} (la valeur qu'on s'accorde à soi-même).

\begin{popfigure}
\begin{tikzpicture}[
  col/.style={rectangle, rounded corners=4pt, text width=6.4cm, align=left, font=\small, inner sep=7pt},
  titre/.style={font=\small\bfseries}
]
\node[col, draw=poppurple, fill=poppurpleL, align=center] (g){%
  \textbf{Ce qu'on voit sur les réseaux (image idéalisée)}\\[3pt]
  $\bullet$ La meilleure photo, choisie parmi 30 essais\\
  $\bullet$ Filtres et retouches sur le visage\\
  $\bullet$ Seulement les fêtes, les cadeaux, les voyages\\
  $\bullet$ Des « vies parfaites » sans problème\\
  $\bullet$ Des vêtements parfois empruntés pour la photo};
\node[col, draw=popgreen, fill=popgreenL, right=0.8cm of g, align=center] (d){%
  \textbf{La réalité quotidienne (qu'on ne voit pas)}\\[3pt]
  $\bullet$ Des journées ordinaires, comme tout le monde\\
  $\bullet$ Des boutons, de la fatigue, des soucis\\
  $\bullet$ Des devoirs, des corvées, des disputes\\
  $\bullet$ Des échecs et des notes décevantes\\
  $\bullet$ Des difficultés d'argent comme beaucoup de familles};
\draw[-{Stealth[length=3mm]}, very thick, poporange] (g.east) -- (d.west)
  node[midway, above, font=\small\bfseries, popdark]{Ne pas confondre !};
\end{tikzpicture}
\legende{Comparer sa vie réelle aux images idéalisées des réseaux sociaux, c'est comparer les coulisses d'un film avec l'affiche du film.}
\end{popfigure}

\begin{aretenir}[Le réflexe anti-comparaison]
Quand une publication te donne un sentiment d'infériorité, rappelle-toi : « je vois la vitrine, pas l'arrière-boutique ». Personne ne publie ses échecs. Ta valeur ne se mesure ni en « j'aime », ni en abonnés, ni en comparaison avec des images retouchées.
\end{aretenir}

\courssub{Le stress et la surcharge informationnelle}

\textbf{Intuition.} Essaie de résoudre un problème de mathématiques pendant que quelqu'un t'appelle par ton nom toutes les deux minutes. Impossible de te concentrer ! C'est exactement ce que font les notifications : elles interrompent le cerveau en permanence.

\begin{definition}[Surcharge informationnelle et FOMO]
La \cle{surcharge informationnelle} est l'état d'un cerveau qui reçoit plus d'informations (messages, notifications, vidéos, actualités) qu'il ne peut en traiter. Elle provoque du \textbf{stress}, de la fatigue mentale et des \textbf{difficultés de concentration}. Elle est aggravée par la \cle{FOMO} (\emph{Fear Of Missing Out}, « peur de rater quelque chose ») : l'angoisse de manquer un message, une nouvelle ou une tendance, qui pousse à vérifier son téléphone sans arrêt.
\end{definition}

\textbf{Conséquences concrètes sur les études :} un devoir qui devrait prendre 30 minutes en prend 2 heures ; on relit trois fois la même phrase ; on retient mal la leçon. Le cerveau a besoin de périodes \emph{sans interruption} pour apprendre.

\textbf{Les bons réflexes :}
\begin{itemize}
  \item désactiver les notifications non essentielles pendant les devoirs et la nuit ;
  \item poser le téléphone dans une autre pièce pendant les révisions ;
  \item consulter ses messages à des moments choisis (par exemple trois fois par jour) au lieu de réagir à chaque vibration ;
  \item se rappeler que rien d'important ne se perd : un vrai message urgent attendra une heure.
\end{itemize}

\courssub{Vie privée, réputation et traces numériques}

\textbf{Intuition.} Écrire sur Internet, c'est comme écrire au marqueur indélébile sur le mur du marché : tout le monde peut le voir, le photographier, et on ne peut plus l'effacer complètement.

\begin{definition}[E-réputation et trace numérique]
Toute publication en ligne (photo, commentaire, vidéo) laisse une \cle{trace numérique} durable : elle peut être copiée, capturée et partagée, même si on la supprime ensuite. L'ensemble des informations qu'on trouve sur une personne en ligne forme son \cle{e-réputation} (réputation numérique), consultée par les camarades, mais aussi plus tard par les employeurs, les écoles ou les administrations.
\end{definition}

\textbf{Les dangers principaux :}
\begin{itemize}
  \item publier une photo de quelqu'un \textbf{sans son accord} (c'est une atteinte à sa vie privée, punie par la loi) ;
  \item publier des propos regrettables (insultes, plaisanteries blessantes) qui ressurgissent des années plus tard ;
  \item partager des informations personnelles (adresse, école, numéro de téléphone) accessibles à des inconnus mal intentionnés.
\end{itemize}

\begin{exemplebox}[La photo « pour rire »]
En classe de 3ème, Brice publie « pour rire » une photo ridicule de son camarade. Un ami la copie, la partage dans un groupe, un autre la republie ailleurs. Trois ans plus tard, le camarade postule à un stage : en recherchant son nom, le responsable tombe sur la photo. Ce qui était une blague de dix secondes est devenu une trace durable. \textbf{Règle d'or : avant de publier, se demander si l'on accepterait que ses parents, son directeur ou un futur employeur voient cette publication dans cinq ans.}
\end{exemplebox}

\courssub{Les contenus inappropriés ou choquants}

En naviguant ou en discutant en ligne, on peut tomber — parfois sans l'avoir cherché — sur des \cle{contenus inappropriés} : images violentes, vidéos choquantes, propos haineux, ou sollicitations d'inconnus (demandes de photos, propositions étranges). Ces contenus peuvent perturber durablement, surtout à ton âge.

\textbf{La conduite à tenir :}
\begin{enumerate}
  \item \textbf{Fermer} immédiatement la page ou quitter la conversation, sans partager le contenu ;
  \item \textbf{Ne pas se sentir coupable} : tomber sur un tel contenu n'est pas une faute ;
  \item \textbf{En parler à un adulte de confiance} (parent, enseignant) : en parler permet de dédramatiser, de comprendre et d'être protégé ;
  \item \textbf{Bloquer} l'expéditeur et \textbf{signaler} le contenu à la plateforme ;
  \item ne jamais répondre aux sollicitations d'inconnus et ne jamais envoyer de photos personnelles.
\end{enumerate}

\begin{exoresolu}[Analyse d'une situation de vie]
\textbf{Énoncé.} Chaque soir depuis deux semaines, Merveille reçoit sur son téléphone des messages insultants d'un groupe de camarades de classe. Ils ont publié une photo truquée d'elle. Choquée, elle supprime tous les messages, ne répond plus, s'isole dans sa chambre, passe ses nuits sur son téléphone pour surveiller ce qui se dit, et ses notes chutent. Elle n'ose rien dire à ses parents.\\[2pt]
1) Identifier les risques psychologiques et sociaux présents dans cette situation.\\
2) Relever les erreurs commises par Merveille.\\
3) Rédiger la démarche correcte qu'elle aurait dû suivre, en la justifiant.
\tcblower
\textbf{Solution détaillée.}

\textbf{1) Les risques présents :}
\begin{itemize}
  \item le \textbf{cyberharcèlement} : insultes répétées et photo humiliante diffusée par des moyens numériques ;
  \item l'\textbf{isolement social} : Merveille se replie sur elle-même et s'enferme ;
  \item le \textbf{stress et la surcharge informationnelle} : elle passe ses nuits à surveiller son téléphone (FOMO), ce qui nuit à son sommeil et à sa concentration ;
  \item les \textbf{conséquences scolaires} : la chute des notes est la conséquence directe de ce stress et de cette fatigue.
\end{itemize}

\textbf{2) Les erreurs commises :}
\begin{itemize}
  \item elle a \textbf{supprimé les messages} : elle a détruit les preuves du harcèlement, alors qu'il fallait d'abord faire des captures d'écran ;
  \item elle \textbf{n'en a parlé à personne} : garder le secret laisse le harcèlement continuer et aggrave l'isolement ;
  \item elle \textbf{surveille son téléphone toute la nuit} : cela entretient le stress et la fatigue au lieu de la protéger.
\end{itemize}

\textbf{3) La démarche correcte (les 4 réflexes) :}
\begin{enumerate}
  \item \textbf{Ne pas répondre} aux insultes : répondre nourrit le conflit et encourage les harceleurs ;
  \item \textbf{Conserver les preuves} : captures d'écran des messages et de la photo, avec dates et pseudos, afin de pouvoir établir les faits ;
  \item \textbf{En parler à un adulte de confiance} (parents, enseignant, surveillant) : un adulte peut faire cesser la situation, alerter l'établissement et soutenir la victime ;
  \item \textbf{Bloquer et signaler} les comptes concernés sur la plateforme, et si nécessaire déposer plainte, car le cyberharcèlement est un délit puni par la loi.
\end{enumerate}
En complément, Merveille devrait limiter son téléphone le soir (notifications coupées, téléphone hors de la chambre) pour retrouver le sommeil et la concentration, et maintenir des activités réelles avec des amis de confiance pour rompre l'isolement.
\end{exoresolu}

\begin{exercice}[À toi de jouer]
Ton petit frère passe six heures par jour sur des jeux en ligne, s'énerve quand on le lui retire et a abandonné ses amis. Par ailleurs, une camarade te montre des messages moqueurs qu'un groupe publie sur elle.\\[2pt]
1) Quels signes d'addiction aux écrans observes-tu chez ton frère ? Propose deux solutions concrètes.\\
2) Explique à ta camarade les 4 réflexes à appliquer, dans l'ordre, en justifiant chacun.\\
3) Pourquoi ne faut-il jamais publier la photo de quelqu'un sans son accord ? Cite deux raisons.
\end{exercice}

\begin{corrige}[Exercice]
\textbf{1)} Signes d'addiction : temps de jeu excessif (six heures), irritabilité quand il ne peut pas jouer, abandon des autres activités et de ses amis. Solutions : fixer des horaires de jeu limités et respectés (par exemple une heure après les devoirs), proposer des activités de remplacement attractives (sport, sorties) et placer la console hors de sa chambre le soir.\\[3pt]
\textbf{2)} Les 4 réflexes : (1) ne pas répondre pour ne pas nourrir le conflit ; (2) conserver les preuves par captures d'écran afin de pouvoir prouver les faits ; (3) en parler à un adulte de confiance car le harcèlement ne cesse presque jamais seul ; (4) bloquer et signaler les comptes, et si nécessaire porter plainte car le cyberharcèlement est puni par la loi.\\[3pt]
\textbf{3)} Publier une photo sans accord est une atteinte à la vie privée, punie par la loi ; de plus, la photo devient une trace numérique durable qui peut être copiée, partagée et nuire à la réputation (e-réputation) de la personne pendant des années.
\end{corrige}

\coursec{Les bonnes pratiques d'hygiène numérique}

Dans la section précédente, nous avons vu que l'usage excessif ou mal maîtrisé des écrans peut entraîner des troubles psychologiques (stress, dépendance, anxiété) et sociaux (isolement, cyberharcèlement). Face à ces risques, la bonne nouvelle est simple : il existe des \cle{bonnes pratiques} concrètes, faciles à appliquer, qui permettent de profiter du numérique sans en subir les méfaits. C'est l'objet de cette section : transformer les dangers étudiés en habitudes de vie saines.

\courssub{Gérer son temps d'écran}

\textbf{Pourquoi limiter le temps d'écran ?}\par
Le temps passé devant un écran « file » sans qu'on s'en rende compte : une vidéo en entraîne une autre, une partie de jeu en appelle une deuxième, et voilà trois heures envolées. Or ce temps est pris sur le sommeil, les devoirs, le sport et la vie en famille. La première bonne pratique consiste donc à \cle{fixer des limites quotidiennes} et à s'y tenir.

Concrètement, trois outils simples existent :
\begin{itemize}
  \item les \cle{minuteurs} (chronomètre du téléphone, réveil de cuisine) qui sonnent quand le temps est écoulé ;
  \item les \cle{applications de contrôle} intégrées aux smartphones, qui mesurent et bloquent l'usage au-delà d'une durée choisie ;
  \item les \cle{plages sans écran} planifiées à l'avance : pendant les repas, pendant les devoirs, et dans l'heure qui précède le coucher.
\end{itemize}

\begin{methode}[Limiter son temps d'écran en 4 étapes]
Pour reprendre le contrôle de ton temps d'écran, suis cette démarche :
\begin{enumerate}
  \item \textbf{Mesurer} son temps actuel : pendant quelques jours, note (ou lis dans les réglages du téléphone) le temps réellement passé chaque jour devant les écrans de loisir.
  \item \textbf{Fixer un objectif réaliste} : par exemple, passer de 4 heures à 2 heures de loisir numérique par jour, pas zéro du premier coup.
  \item \textbf{Utiliser un minuteur} : programme une alarme ou une limite d'application qui t'avertit quand l'objectif est atteint.
  \item \textbf{Évaluer chaque semaine} : compare le temps réel à l'objectif, félicite-toi des progrès et ajuste si nécessaire.
\end{enumerate}
\end{methode}

\begin{savaistu}[Le bien-être numérique est déjà dans ta poche]
La plupart des smartphones récents intègrent \emph{gratuitement} un outil de « \cle{bien-être numérique} » (parfois appelé « Temps d'écran » ou « Digital Wellbeing »). Il mesure automatiquement le temps passé sur chaque application, permet de fixer des limites quotidiennes et propose un « mode coucher » qui grise l'écran le soir. Inutile d'installer quoi que ce soit : il suffit d'ouvrir les réglages du téléphone !
\end{savaistu}

\courssub{Protéger ses yeux : la règle des 20-20-20}

\textbf{Le problème.} Quand on fixe un écran de près pendant longtemps, les muscles des yeux restent contractés et l'on cligne moins des paupières : d'où les picotements, les maux de tête et la vision floue étudiés dans la section 2.

\textbf{La solution.} Les ophtalmologues recommandent une règle très simple à mémoriser : la \cle{règle des 20-20-20}.

\begin{propriete}[Règle des 20-20-20 (admise)]
Toutes les \textbf{20 minutes} passées devant un écran, regarder pendant \textbf{20 secondes} un objet situé à environ \textbf{20 pieds} (soit environ 6 mètres) : par exemple un arbre par la fenêtre ou le mur du fond de la cour. Cette règle, issue des recommandations des ophtalmologues, réduit significativement la fatigue visuelle. On l'applique telle quelle, sans démonstration.
\end{propriete}

\begin{exemplebox}[Appliquer la règle 20-20-20 pendant les devoirs]
Mariam fait ses devoirs sur l'ordinateur du lycée pendant 1 heure. Elle programme trois alarmes : à 20 min, 40 min et 60 min. À chaque alarme, elle lève les yeux et regarde le flamboyant de la cour, à une dizaine de mètres, en comptant jusqu'à 20. Résultat : elle termine sa séance sans maux de tête, alors qu'auparavant ses yeux piquaient dès la demi-heure.
\end{exemplebox}

\courssub{L'ergonomie du poste de travail}

L'\cle{ergonomie} est l'art d'aménager son poste de travail pour que le corps reste dans une position confortable et sans danger. Un mauvais réglage (écran trop bas, chaise trop haute) provoque à la longue des douleurs au cou, au dos et aux poignets.

\begin{aretenir}[Les 5 règles d'un poste de travail ergonomique]
\begin{enumerate}
  \item Le \textbf{haut de l'écran} arrive à \textbf{hauteur des yeux} (regard légèrement vers le bas).
  \item La \textbf{distance écran-yeux} est de \textbf{50 à 70 cm} (environ la longueur d'un bras).
  \item Le \textbf{dos est droit}, appuyé contre le dossier, avec un angle d'environ 90 à 100° au niveau des hanches.
  \item Les \textbf{pieds sont à plat au sol} (ou sur un repose-pieds), les cuisses à l'horizontale.
  \item La pièce est \textbf{bien éclairée}, sans reflet sur l'écran ni fenêtre juste derrière lui.
\end{enumerate}
\end{aretenir}

\begin{popfigure}
\begin{center}
\begin{tikzpicture}[scale=0.9, every node/.style={font=\small}]
  % Sol
  \draw[thick, popdark] (0,0) -- (10.5,0);
  % Chaise
  \draw[fill=popblueL, draw=popblue] (2.2,0) rectangle (2.45,1.5);
  \draw[fill=popblueL, draw=popblue] (2.2,1.5) rectangle (4.4,1.75);
  \draw[fill=popblueL, draw=popblue] (4.15,0) rectangle (4.4,1.5);
  % Personnage (simplifié, assis)
  \draw[fill=popgoldL, draw=popgold, thick] (2.9,1.75) -- (3.0,3.3); % dos
  \draw[fill=popgoldL, draw=popgold, thick] (2.95,3.5) circle (0.32); % tête
  \draw[popgold, thick] (2.9,1.75) -- (4.15,1.8); % cuisse
  \draw[popgold, thick] (4.15,1.8) -- (4.2,0.15); % jambe
  \draw[popgold, thick] (4.2,0.15) -- (4.75,0.12); % pied
  \draw[popgold, thick] (3.0,3.1) -- (5.6,2.35); % bras
  % Angle dos
  \draw[poporange, thick, -{Stealth}] (2.95,2.6) arc (90:8:0.85);
  \node[poporange, font=\footnotesize] at (3.85,2.95) {90--100°};
  % Bureau
  \draw[fill=popcream, draw=popdark, thick] (5.4,2.2) rectangle (9.6,2.4);
  \draw[fill=popcream, draw=popdark, thick] (9.2,0) rectangle (9.5,2.2);
  % Écran
  \draw[fill=popblue, draw=popdark, thick] (6.4,2.4) rectangle (8.3,4.1);
  \draw[fill=popblueL] (6.5,2.5) rectangle (8.2,4.0);
  \draw[popdark, thick] (7.2,2.4) -- (7.5,2.4) -- (7.35,2.15) -- cycle;
  % Clavier
  \draw[fill=popgreenL, draw=popgreen, thick] (5.7,2.4) rectangle (6.9,2.75);
  % Ligne de regard
  \draw[dashed, poppink, thick] (3.3,3.55) -- (8.3,3.9);
  \node[poppink, font=\footnotesize, align=center] at (5.6,4.35) {regard à hauteur\\du haut de l'écran};
  % Cote distance yeux-écran
  \draw[{Stealth}-{Stealth}, poppurple, thick] (3.3,1.1) -- (6.4,1.1);
  \node[poppurple, font=\footnotesize] at (4.85,0.8) {distance yeux--écran : 50 à 70 cm};
  % Lampe / éclairage
  \draw[fill=poppinkL, draw=poppink, thick] (9.0,2.4) -- (9.0,3.3) -- (8.55,3.55);
  \draw[fill=poppinkL, draw=poppink, thick] (8.35,3.45) circle (0.18);
  \node[poppink, font=\footnotesize, align=center] at (9.35,3.85) {bon éclairage\\de la pièce};
  % Pieds au sol
  \node[popgreen, font=\footnotesize] at (5.6,0.35) {pieds à plat au sol};
\end{tikzpicture}
\end{center}
\legende{Schéma coté d'un poste de travail ergonomique : écran à hauteur des yeux, distance de 50 à 70 cm, dos à 90--100°, pieds au sol, pièce bien éclairée.}
\end{popfigure}

\begin{attention}[Le téléphone qui charge sous l'oreiller]
Erreur fréquente : laisser son téléphone \textbf{charger sous l'oreiller ou sur le lit} pendant la nuit. C'est dangereux à deux titres : l'appareil mal ventilé peut \textbf{surchauffer} (risque de brûlure, voire d'incendie), et la tentation de consulter l'écran à chaque réveil dégrade fortement la \textbf{qualité du sommeil}. La bonne habitude : charger le téléphone sur une table, loin du lit, ou le laisser dans une autre pièce.
\end{attention}

\courssub{Protéger son sommeil et son audition}

\textbf{Le sommeil.} La lumière bleue des écrans trompe le cerveau, qui croit qu'il fait encore jour : l'endormissement devient difficile. Deux règles simples suffisent :
\begin{itemize}
  \item \textbf{arrêter les écrans} (téléphone, tablette, télévision, console) au moins \textbf{30 minutes à 1 heure avant le coucher} ;
  \item laisser le téléphone \textbf{hors de la chambre} ou le passer en \textbf{mode avion} pendant la nuit.
\end{itemize}

\textbf{L'audition.} Écouter de la musique trop fort et trop longtemps au casque endommage l'oreille de façon \emph{irréversible}. On applique la \cle{règle du 60/60}.

\begin{propriete}[Règle du 60/60]
Au casque ou avec des écouteurs : régler le volume à \textbf{60 \% du maximum} au plus, et ne pas dépasser \textbf{60 minutes d'écoute d'affilée} avant de faire une pause.
\end{propriete}

\begin{exemplebox}[La règle 60/60 en pratique]
Appliquer la règle 60/60 revient concrètement à écouter sa musique \textbf{à mi-volume} (curseur réglé à 60 \%) par \textbf{séances d'une heure maximum}, entrecoupées de pauses. Par exemple, Junior écoute ses chansons de 17 h à 18 h à mi-volume pendant ses révisions, puis enlève le casque. À ce rythme, il préserve son audition sur le long terme, alors que son ami qui écoute à volume maximal toute la soirée s'expose à une baisse définitive de l'ouïe avant l'âge adulte.
\end{exemplebox}

\courssub{Alterner les activités et soigner ses relations en ligne}

L'hygiène numérique ne consiste pas seulement à « moins » utiliser les écrans : elle consiste aussi à \textbf{mieux vivre à côté}. Deux familles d'habitudes y contribuent.

\textbf{1. L'alternance des activités.} À chaque heure d'écran doit correspondre du temps pour d'autres activités : pratiquer un \textbf{sport} (football, course, danse), \textbf{lire} un livre, participer aux \textbf{tâches familiales} (cuisine, ménage, courses au marché) et aux \textbf{jeux de plein air} avec les amis du quartier. Ces activités reposent les yeux, font bouger le corps et entretiennent les liens réels.

\textbf{2. L'hygiène relationnelle en ligne.} Sur Internet, on applique la \cle{netiquette}, c'est-à-dire les règles de politesse et de respect du vivre-ensemble numérique :
\begin{itemize}
  \item rester \textbf{poli et respectueux} dans les messages et commentaires, comme en face à face ;
  \item \textbf{demander l'autorisation} avant de publier la photo de quelqu'un ;
  \item \textbf{vérifier une information} avant de la partager, pour ne pas propager de fausses nouvelles (\emph{fake news}).
\end{itemize}

\begin{popfigure}
\begin{center}
\begin{tikzpicture}[scale=1, every node/.style={font=\small}]
  % Axe horaire de 6h à 22h
  \draw[thick, popdark, -{Stealth}] (0,0) -- (16.6,0);
  \foreach \x/\h in {0/6h, 2/8h, 4/10h, 6/12h, 8/14h, 10/16h, 12/18h, 14/20h, 16/22h}{
    \draw[popdark] (\x,0.08) -- (\x,-0.08);
    \node[below, font=\footnotesize, popmuted] at (\x,-0.1) {\h};
  }
  % Blocs d'activités (1 unité = 1 heure)
  \draw[fill=popblueL, draw=popblue] (1,0.35) rectangle (7,1.05);
  \node[popdark, font=\footnotesize] at (4,0.7) {École (7h--13h)};
  \draw[fill=popgreenL, draw=popgreen] (7,0.35) rectangle (8,1.05);
  \node[popdark, font=\footnotesize] at (7.5,0.7) {Repas};
  \draw[fill=popgoldL, draw=popgold] (8,0.35) rectangle (10,1.05);
  \node[popdark, font=\footnotesize] at (9,0.7) {Devoirs};
  \draw[fill=poporangeL, draw=poporange] (10,0.35) rectangle (12,1.05);
  \node[popdark, font=\footnotesize] at (11,0.7) {Sport / jeux};
  \draw[fill=poppurpleL, draw=poppurple] (12,0.35) rectangle (13,1.05);
  \node[popdark, font=\footnotesize] at (12.5,0.7) {Écran};
  \draw[fill=popgreenL, draw=popgreen] (13,0.35) rectangle (14,1.05);
  \node[popdark, font=\footnotesize] at (13.5,0.7) {Famille};
  \draw[fill=poppinkL, draw=poppink] (14,0.35) rectangle (15,1.05);
  \node[popdark, font=\footnotesize, align=center] at (14.5,0.7) {Lecture};
  \draw[fill=popcream, draw=popmuted] (15,0.35) rectangle (16,1.05);
  \node[popmuted, font=\footnotesize] at (15.5,0.7) {Coucher};
  % Annotations
  \node[poppurple, font=\footnotesize, align=center] at (12.5,1.55) {écran limité : 1 h max};
  \draw[poppurple, -{Stealth}] (12.5,1.35) -- (12.5,1.08);
  \node[popmuted, font=\footnotesize, align=center] at (15.5,1.55) {sans téléphone};
  \draw[popmuted, -{Stealth}] (15.5,1.35) -- (15.5,1.08);
\end{tikzpicture}
\end{center}
\legende{Frise horaire d'une journée équilibrée d'un élève de 3ème : école, repas en famille, devoirs, sport, écran de loisir limité à une heure, puis coucher sans téléphone vers 21 h.}
\end{popfigure}

\courssub{Hygiène matérielle : un matériel propre et sain}

Dernier réflexe, souvent oublié : l'\cle{hygiène matérielle}. Le clavier et l'écran accumulent poussières et microbes, surtout dans les salles multimédias partagées du collège. Deux habitudes suffisent :
\begin{itemize}
  \item \textbf{nettoyer régulièrement} l'écran (chiffon doux légèrement humide, appareil éteint) et le clavier (chiffon, souffle d'air pour les touches) ;
  \item \textbf{ne pas manger au-dessus du clavier} : les miettes bloquent les touches et attirent les insectes.
\end{itemize}

\begin{exoresolu}[Construire sa journée équilibrée]
\textbf{Énoncé.} Alain, élève de 3ème, rentre de l'école à 14 h et se couche à 21 h 30. Il passe actuellement 5 h par soirée sur son téléphone, s'endort difficilement et a souvent mal aux yeux. Propose-lui un emploi du temps du soir qui applique les bonnes pratiques d'hygiène numérique, et justifie chaque choix.
\tcblower
\textbf{Solution détaillée.} On applique les règles de la leçon : limitation du temps d'écran, alternance des activités, protection du sommeil (écrans arrêtés au moins 30 min à 1 h avant le coucher).
\begin{itemize}
  \item \textbf{14 h -- 14 h 30 : repas et repos, sans écran} (plage sans écran pendant le repas).
  \item \textbf{14 h 30 -- 16 h : devoirs.} Si l'ordinateur est utilisé, Alain applique la règle des 20-20-20 et règle son poste (écran à hauteur des yeux, 50--70 cm).
  \item \textbf{16 h -- 17 h 30 : sport ou jeux de plein air} avec les amis (alternance des activités, bon pour le corps et le moral).
  \item \textbf{17 h 30 -- 18 h 30 : temps d'écran de loisir limité à 1 h}, contrôlé par un minuteur (méthode en 4 étapes : mesurer, fixer, minuter, évaluer).
  \item \textbf{18 h 30 -- 20 h : tâches familiales, dîner, lecture} (plages sans écran).
  \item \textbf{20 h 30 : téléphone posé en mode avion hors de la chambre}, soit 1 h avant le coucher (protection du sommeil, pas de charge sous l'oreiller).
  \item \textbf{21 h 30 : coucher.}
\end{itemize}
\textbf{Justification.} Le temps d'écran passe de 5 h à 1 h, ce qui libère du temps pour le sommeil et le sport ; l'arrêt des écrans 1 h avant le coucher facilite l'endormissement ; la règle 20-20-20 et l'ergonomie suppriment les maux de tête et les douleurs. Chaque choix répond donc à un risque précis identifié dans les sections précédentes.
\end{exoresolu}

\begin{exercice}[À toi de jouer]
\begin{enumerate}
  \item Énumère les cinq règles d'un poste de travail ergonomique.
  \item Explique, avec tes mots, la règle des 20-20-20 et la règle du 60/60.
  \item Pendant une semaine, mesure ton temps d'écran quotidien, puis applique la méthode en 4 étapes pour le réduire de moitié. Rédige un court bilan justifiant ta démarche.
\end{enumerate}
\end{exercice}

\begin{corrige}[Exercice]
\begin{enumerate}
  \item Écran à hauteur des yeux ; distance yeux-écran de 50 à 70 cm ; dos droit (angle de 90 à 100°) ; pieds à plat au sol ; pièce bien éclairée sans reflet sur l'écran.
  \item Règle 20-20-20 : toutes les 20 minutes d'écran, regarder 20 secondes un objet à environ 6 mètres pour reposer les yeux. Règle 60/60 : au casque, volume à 60 \% maximum et pas plus de 60 minutes d'affilée.
  \item Le bilan doit suivre la démarche : mesure du temps initial (par exemple 4 h/jour), objectif réaliste (2 h/jour), utilisation d'un minuteur ou de l'outil de bien-être numérique, puis évaluation hebdomadaire comparant le temps réel à l'objectif, avec une conclusion justifiée.
\end{enumerate}
\end{corrige}

\begin{aretenir}[L'essentiel de la section]
L'hygiène numérique repose sur des gestes simples et quotidiens : \textbf{limiter} son temps d'écran (minuteur, plages sans écran), \textbf{ménager ses yeux} (règle 20-20-20), \textbf{aménager son poste} (écran à hauteur des yeux, 50--70 cm, dos droit), \textbf{protéger son sommeil} (écrans arrêtés 1 h avant le coucher, téléphone hors de la chambre), \textbf{ménager ses oreilles} (règle 60/60), \textbf{alterner} avec le sport et la vie familiale, respecter la \textbf{netiquette} en ligne et garder un \textbf{matériel propre}. Ces habitudes, appliquées ensemble, permettent de profiter pleinement du numérique tout en préservant sa santé.
\end{aretenir}

\coursec{Usage responsable des réseaux sociaux et bien-être social}

Dans la section précédente, nous avons appris à adopter de bonnes pratiques d'hygiène numérique : protéger nos yeux et notre dos devant l'écran, faire des pauses régulières, limiter le temps passé en ligne. Mais l'hygiène du numérique ne concerne pas seulement notre corps : elle concerne aussi nos \textbf{relations avec les autres} et notre \textbf{équilibre de vie}. Un adolescent peut avoir une excellente posture devant son téléphone et pourtant souffrir, ou faire souffrir, à cause de ce qu'il publie sur les réseaux sociaux. C'est ce que nous allons étudier maintenant : comment utiliser WhatsApp, Facebook, TikTok ou Instagram de manière \cle{responsable}, pour que le numérique reste un outil de bonheur et non une source de conflits.

\courssub{La netiquette : être poli et respectueux en ligne}

\textbf{Intuition.} Dans la cour de récréation, tu ne cries pas sur les autres, tu n'insultes pas un camarade qui n'est pas d'accord avec toi, tu attends ton tour pour parler. Pourquoi ? Parce que vivre ensemble demande des règles de politesse. Sur Internet, c'est exactement pareil : des millions de personnes discutent ensemble, et sans règles de respect, les échanges deviennent vite des disputes. L'ensemble de ces règles de bonne conduite en ligne porte un nom : la \cle{netiquette} (mot formé de \emph{net}, le réseau, et \emph{étiquette}, les bonnes manières).

\begin{definition}[La netiquette]
La \cle{netiquette} est l'ensemble des règles de politesse et de respect à observer lorsqu'on communique sur Internet (réseaux sociaux, messageries, groupes de discussion, jeux en ligne).
\end{definition}

Voici les règles essentielles de la netiquette :

\begin{itemize}
\item \textbf{Ne pas insulter ni humilier}, même en plaisantant : derrière chaque écran se trouve une personne réelle, avec des sentiments réels.
\item \textbf{Ne pas écrire en MAJUSCULES} : sur Internet, écrire en majuscules équivaut à \emph{crier}. « TU VIENS OU PAS ? » donne l'impression d'une agression, alors que « Tu viens ou pas ? » est une simple question.
\item \textbf{Respecter les opinions des autres} : on peut ne pas être d'accord avec quelqu'un et le dire calmement, sans moquerie ni attaque personnelle.
\item \textbf{Relire son message avant de l'envoyer} : un message écrit sous le coup de la colère ne peut plus être repris une fois envoyé.
\item \textbf{Ne pas envoyer de messages en masse} (chaînes, publicités non demandées) qui encombrent les téléphones des autres.
\end{itemize}

\begin{attention}[Erreur fréquente : « Ce n'est que du virtuel »]
Beaucoup d'élèves pensent que ce qui se passe en ligne « n'est pas grave » parce que ce n'est « pas la vraie vie ». C'est faux : une insulte publiée dans un groupe de classe blesse autant qu'une insulte criée dans la cour, et elle reste visible par des dizaines de personnes. Les paroles s'envolent, mais les écrits restent.
\end{attention}

\courssub{Protéger sa vie privée}

\textbf{Intuition.} Tu ne donnerais pas l'adresse de ta maison ni le nom de ton école à un inconnu croisé dans la rue à Mokolo ou à Akwa. Sur Internet, les inconnus sont encore plus nombreux et parfois mal intentionnés : il faut donc appliquer la même prudence, et même davantage.

\begin{propriete}[Règles de protection de la vie privée en ligne]
Pour se protéger sur les réseaux sociaux, un élève doit :
\begin{enumerate}
\item \textbf{Régler les paramètres de confidentialité} de ses comptes : profil privé, publications visibles uniquement par les amis, localisation désactivée.
\item \textbf{Ne jamais partager publiquement ses informations personnelles} : adresse du domicile, nom de l'école, numéro de téléphone, horaires habituels, photos intimes.
\item \textbf{Être prudent avec les inconnus} : ne pas accepter toutes les demandes d'ami, se méfier des personnes qui posent trop de questions personnelles, ne jamais accepter de rencontrer seul un « ami » connu uniquement en ligne.
\item \textbf{Utiliser des mots de passe solides et différents} pour chaque compte, et ne jamais les communiquer, même à un ami proche.
\end{enumerate}
\end{propriete}

\begin{attention}[Ce qui est supprimé ne disparaît pas toujours]
Erreur très fréquente : croire qu'une publication supprimée disparaît définitivement. En réalité, n'importe qui peut faire une \textbf{capture d'écran} avant la suppression, et cette image peut circuler pour toujours. Une photo intime envoyée « pour une minute » peut être conservée, copiée et diffusée. Règle d'or : \textbf{ne publie jamais ce que tu ne voudrais pas voir affiché au tableau de ta classe}.
\end{attention}

\courssub{L'e-réputation : l'image que tu laisses sur Internet}

\textbf{Intuition.} Quand tu arrives dans un nouveau quartier, les gens se font une idée de toi à partir de ce qu'ils voient et entendent : ta tenue, tes paroles, tes actes. Sur Internet, c'est pareil, mais à une échelle beaucoup plus grande : tout ce que tu publies (photos, commentaires, vidéos) construit une image de toi, visible par beaucoup de monde, pendant très longtemps.

\begin{definition}[L'e-réputation]
L'\cle{e-réputation} (ou réputation numérique) est l'image qu'une personne laisse d'elle-même sur Internet à travers ses publications, ses photos, ses vidéos et ses commentaires.
\end{definition}

Cette image a des conséquences réelles. Un futur employeur, un responsable d'école ou même une famille pour un mariage peuvent rechercher ton nom sur Internet. Une vieille vidéo moqueuse ou un commentaire insultant peut alors te fermer des portes, des années plus tard. Avant de publier, pose-toi la question simple : \emph{« Est-ce que je montrerais cela à mes parents ou à mon futur employeur ? »} Si la réponse est non, ne publie pas.

\begin{methode}[Le test des 4 questions avant de publier]
Avant d'appuyer sur « Publier » ou « Envoyer », pose-toi ces quatre questions, dans l'ordre :
\begin{enumerate}
\item \textbf{Est-ce vrai ?} (Ai-je vérifié l'information ?)
\item \textbf{Est-ce utile ?} (Cette publication apporte-t-elle quelque chose de positif ?)
\item \textbf{Est-ce respectueux ?} (Cela ne blesse-t-il personne ?)
\item \textbf{Ai-je l'autorisation ?} (Les personnes visibles ou concernées sont-elles d'accord ?)
\end{enumerate}
Si une seule réponse est \textbf{non}, alors \textbf{on s'abstient de publier}.
\end{methode}

\begin{popfigure}
\begin{center}
\begin{tikzpicture}[
  font=\small,
  quest/.style={draw=popblue, fill=popblueL, diamond, aspect=2.2, align=center, inner sep=1pt, text width=2.6cm},
  action/.style={draw=popgreen, fill=popgreenL, rounded corners, align=center, minimum width=2.8cm, minimum height=0.9cm},
  stopb/.style={draw=poppink, fill=poppinkL, rounded corners, align=center, minimum width=2.8cm, minimum height=0.9cm},
  fleche/.style={-{Stealth[length=2.5mm]}, thick, popdark}
]
\node[quest] (q1) at (0,0) {Est-ce vrai ?};
\node[quest] (q2) at (0,-2) {Est-ce utile ?};
\node[quest] (q3) at (0,-4) {Est-ce respectueux ?};
\node[quest] (q4) at (0,-6) {Ai-je l'autorisation ?};
\node[action] (ok) at (0,-8.2) {\textbf{PUBLIER}};
\node[stopb, align=center] (no) at (5,-3){\textbf{S'ABSTENIR}\\ (ne pas publier)};
\draw[fleche] (q1) -- node[right]{oui} (q2);
\draw[fleche] (q2) -- node[right]{oui} (q3);
\draw[fleche] (q3) -- node[right]{oui} (q4);
\draw[fleche] (q4) -- node[right]{oui} (ok);
\draw[fleche] (q1.east) -- node[above, pos=0.4]{non} (no.north west);
\draw[fleche] (q2.east) -- node[above, pos=0.4]{non} (no.west);
\draw[fleche] (q3.east) -- node[above, pos=0.4]{non} (no.south west);
\draw[fleche] (q4.east) -- ++(2,0) |- node[right, pos=0.25]{non} (no.south);
\end{tikzpicture}
\end{center}
\legende{Organigramme « Réfléchir avant de publier » : quatre questions successives. Une seule réponse « non » suffit à s'abstenir.}
\end{popfigure}

\courssub{Le respect d'autrui et la lutte contre les fausses informations}

\textbf{Respecter les autres, c'est aussi respecter leur image.} Deux règles sont absolues :

\begin{itemize}
\item \textbf{Ne jamais publier la photo ou la vidéo de quelqu'un sans son autorisation.} Même une photo amusante prise à la cantine peut humilier un camarade si elle circule dans tous les groupes de l'école.
\item \textbf{Ne jamais transférer les messages privés} qu'on t'a envoyés en confiance : une conversation privée doit rester privée.
\end{itemize}

\textbf{Les fausses informations (fake news).} Sur les réseaux sociaux, les rumeurs circulent très vite : fausse alerte à la maladie, fausse annonce de concours, photo truquée. Partager une fausse information, même sans mauvaise intention, peut créer la panique ou détruire la réputation de quelqu'un.

\begin{methode}[Vérifier une information avant de la partager]
\begin{enumerate}
\item \textbf{Regarder la source} : qui a publié ? Un site connu et sérieux, ou un compte anonyme ?
\item \textbf{Croiser l'information} : est-ce que d'autres sources fiables (radio nationale, sites officiels, journaux reconnus) disent la même chose ?
\item \textbf{Vérifier la date} : une vieille information ressortie de son contexte peut devenir fausse.
\item \textbf{En cas de doute, ne pas partager} : ne jamais relayer une rumeur, même « au cas où ».
\end{enumerate}
\end{methode}

\begin{exemplebox}[Une rumeur peut devenir un cyberharcèlement]
Dans un groupe WhatsApp de classe, un élève lit que « Marlyse a triché au dernier devoir ». Sans vérifier, il transfère le message dans deux autres groupes. En quelques heures, des dizaines d'élèves se moquent de Marlyse, certains l'insultent, elle n'ose plus venir en classe. Pourtant, l'information était fausse. L'élève qui a transféré n'avait \emph{aucune mauvaise intention} : il a pourtant participé à un \cle{cyberharcèlement} collectif. Moralité : \textbf{celui qui relaie une rumeur est aussi responsable que celui qui l'a inventée}.
\end{exemplebox}

\begin{savaistu}[La loi camerounaise te protège... et te responsabilise]
Au Cameroun, la \textbf{loi n° 2010/012 du 21 décembre 2010} relative à la cybersécurité et à la cybercriminalité punit certains abus commis en ligne : diffusion de fausses informations, atteinte à la vie privée, harcèlement par voie électronique, usurpation d'identité. Ce que l'on fait derrière un écran peut donc avoir des conséquences devant un tribunal. Être responsable en ligne, c'est aussi être un bon citoyen.
\end{savaistu}

\courssub{L'équilibre entre vie en ligne et vie réelle}

\textbf{Intuition.} Un repas équilibré contient plusieurs aliments ; une vie équilibrée contient plusieurs activités. Si tout ton temps libre est absorbé par l'écran, tu perds ce que le téléphone ne peut pas te donner : l'amitié réelle, la santé du corps, la complicité de la famille.

Le bien-être social repose sur des \textbf{rencontres réelles} : jouer au football avec les voisins, discuter en famille le soir, participer à une association, à un mouvement de jeunesse ou au club d'informatique de l'école. Ces activités développent des qualités qu'aucun réseau social ne remplace : le courage, la solidarité, le sens du collectif.

\begin{popfigure}
\begin{center}
\begin{tikzpicture}[scale=0.85, transform shape]
  \def\maxwidth{11.5cm}
  \def\barheight{0.55cm}
  \def\gap{0.25cm}
  \foreach \i/\label/\val/\color in {
    1/{Écran (réseaux, jeux)}/30/popblue,
    2/Sport et activités physiques/20/popgreen,
    3/Famille et discussions/20/poporange,
    4/Lecture et devoirs personnels/15/poppurple,
    5/Amis et associations/15/popgold} {
    \pgfmathsetmacro{\y}{-(\i-1)*(\barheight+\gap)}
    \pgfmathsetmacro{\w}{\val/30*\maxwidth}
    \draw[fill=\color, draw=\color!50!popdark] (0, \y) rectangle (\w, \y+\barheight);
    \node[anchor=east, font=\small] at (-0.2, \y+\barheight/2) {\label};
    \node[anchor=west, font=\small] at (\w+0.2, \y+\barheight/2) {\val\%};
  }
  \node[anchor=south, font=\bfseries\small] at (\maxwidth/2, \barheight+\gap) {Répartition équilibrée du temps libre (exemple)};
\end{tikzpicture}
\end{center}
\legende{Exemple de répartition équilibrée du temps libre d'un élève : l'écran garde une place, mais ne prend pas toute la place.}
\end{popfigure}

\begin{aretenir}[Les signes d'un déséquilibre]
Tu dois réagir si : tu négliges tes devoirs pour rester en ligne, tu dors moins à cause du téléphone, tu t'énerves quand on te coupe Internet, tu préfères toujours l'écran aux sorties avec les amis. Dans ce cas, parle-en à un adulte et réduis progressivement ton temps d'écran.
\end{aretenir}

\courssub{Le rôle de la famille et de l'école}

L'hygiène numérique ne se construit pas seul. La \textbf{famille} et l'\textbf{école} sont tes deux meilleurs alliés :

\begin{itemize}
\item \textbf{Dialoguer avec ses parents} : leur montrer ce que tu fais en ligne, accepter de discuter de tes usages sans cachette. La confiance se construit par la transparence.
\item \textbf{Convenir ensemble de règles d'usage} : par exemple, pas de téléphone pendant les repas et après 21 h, temps d'écran limité les jours d'école, téléphone posé hors de la chambre la nuit. Des règles discutées ensemble sont mieux acceptées que des interdictions subies.
\item \textbf{Recourir aux adultes de confiance en cas de problème} : si tu es victime de moqueries, de menaces ou de harcèlement en ligne, ne reste jamais seul avec ton problème. Parle-en à tes parents, à ton professeur principal, au censeur ou à un enseignant de confiance. Garder les captures d'écran comme preuves.
\end{itemize}

\begin{experience}[Élaborer une charte d'hygiène numérique familiale]
\textbf{Objectif.} Mettre en pratique les notions de netiquette, de vie privée et d'équilibre en rédigeant collectivement une charte d'usage du téléphone à la maison.
\textbf{Matériel.} Papier, stylos, ou un traitement de texte.
\textbf{Démarche.}
\begin{enumerate}
\item \textbf{En famille}, lister les usages actuels (heures, applications, conflits fréquents).
\item \textbf{Négocier} 3 à 5 règles claires (ex. : « Téléphone dans le salon pour la charge nocturne », « Pas d'écran pendant les repas », « Demander l'autorisation avant de poster une photo de famille »).
\item \textbf{Rédiger} la charte en justifiant chaque règle par un bénéfice (santé, respect, réussite scolaire, sécurité).
\item \textbf{Signer} la charte et l'afficher dans un lieu visible. Prévoir un bilan dans un mois.
\end{enumerate}
\textbf{Observation attendue.} Les règles acceptées par tous sont mieux respectées. Le dialogue réduit les conflits liés au numérique.
\end{experience}

\begin{exoresolu}[Analyser une situation et proposer une démarche responsable]
\textbf{Énoncé.} Junior, élève de 3\textsuperscript{e}, reçoit dans le groupe WhatsApp de sa classe une photo montée d'un camarade, présenté de façon ridicule, avec le message : « Transfère à tous tes groupes ! » Junior trouve la photo drôle mais hésite. Que doit-il faire ? Justifie ta réponse en quatre étapes.
\tcblower
\textbf{Solution détaillée.}
\begin{enumerate}
\item \textbf{Appliquer le test des 4 questions.} Est-ce vrai ? Non, c'est un montage. Est-ce utile ? Non, cela n'apporte rien de positif. Est-ce respectueux ? Non, cela ridiculise un camarade. A-t-il l'autorisation ? Non. Quatre réponses négatives : il faut \textbf{s'abstenir}.
\item \textbf{Ne pas transférer, ne pas commenter, ne pas « aimer »} : toute réaction, même un rire, encourage l'auteur et fait circuler la photo.
\item \textbf{Agir positivement} : écrire en privé à l'auteur pour lui demander de supprimer la photo, et soutenir le camarade visé.
\item \textbf{Alerter un adulte} si la photo continue de circuler : professeur principal ou parents, en montrant les captures d'écran comme preuves.
\end{enumerate}
\textbf{Conclusion.} Junior protège son camarade, protège sa propre e-réputation et applique la netiquette : c'est cela, un usage responsable des réseaux sociaux.
\end{exoresolu}

\begin{exercice}[À toi de jouer]
\begin{enumerate}
\item Ta petite sœur veut publier sur Facebook une photo de famille où l'on voit ta maison et le nom de ton école sur ton uniforme. Explique-lui deux risques et deux conseils.
\item Un ami t'envoie une « information » surprenante sur une star camerounaise. Décris les étapes à suivre avant de la partager.
\item Rédige, avec ta famille, trois règles d'usage du téléphone à la maison, et justifie chacune d'elles.
\end{enumerate}
\end{exercice}

\begin{corrige}[Exercice 1]
\begin{enumerate}
\item \textbf{Risques} : (a) des inconnus peuvent identifier l'adresse de la maison et l'école, ce qui met la famille en danger ; (b) la photo peut être copiée et réutilisée hors de son contexte. \textbf{Conseils} : (a) flouter ou éviter les éléments identifiants (uniforme, portail, adresse) ; (b) régler la confidentialité sur « amis uniquement » et demander l'accord de toutes les personnes visibles.
\item (1) Regarder la source du message ; (2) vérifier sur au moins deux sources fiables (sites d'information reconnus) ; (3) vérifier la date de l'information ; (4) en cas de doute, ne pas partager.
\item Exemples de règles justifiées : « pas de téléphone pendant les repas » (pour préserver le dialogue familial) ; « téléphone éteint à 21 h » (pour protéger le sommeil) ; « les devoirs avant les réseaux sociaux » (pour protéger les résultats scolaires).
\end{enumerate}
\end{corrige}

\begin{aretenir}[L'essentiel de la section]
\begin{itemize}
\item La \cle{netiquette} : politesse et respect en ligne (pas d'insultes, pas de majuscules, respect des opinions).
\item \textbf{Vie privée} : paramètres de confidentialité, aucune information personnelle en public, prudence avec les inconnus.
\item L'\cle{e-réputation} : tout ce que tu publies construit ton image, durablement. Avant de publier : \textbf{vrai ? utile ? respectueux ? autorisé ?}
\item \textbf{Respect d'autrui} : jamais de photo sans autorisation, jamais de message privé transféré.
\item \textbf{Fake news} : vérifier la source, croiser l'information, ne jamais relayer une rumeur.
\item \textbf{Équilibre} : l'écran est une partie du temps libre, pas tout le temps libre ; famille et école sont des alliés, pas des adversaires.
\end{itemize}
\end{aretenir}

\coursec{Élaborer et justifier une démarche d'hygiène numérique}

Dans la section précédente, nous avons appris à utiliser les réseaux sociaux de manière responsable afin de préserver notre bien-être social. Mais savoir \emph{ce qu'il faut faire} ne suffit pas : encore faut-il savoir \emph{comment le faire}, de façon organisée, et être capable d'expliquer \emph{pourquoi} on le fait. C'est exactement l'objet de cette dernière section : transformer toutes les notions de l'unité (risques physiques, risques psychologiques et sociaux, bonnes pratiques) en une \cle{démarche} claire, ordonnée et justifiée, que tu pourras appliquer à ta propre vie ou à celle d'un camarade, et rédiger proprement lors d'une évaluation.

\courssub{Pourquoi une démarche plutôt que des conseils en vrac ?}

Imagine un infirmier scolaire à qui un élève dit : « Je suis fatigué. » Il ne va pas donner un médicament au hasard. Il va d'abord \emph{observer} (prendre la température, poser des questions), puis \emph{diagnostiquer} (identifier la cause), puis \emph{prescrire} (proposer un traitement), puis \emph{suivre} le patient (vérifier que ça marche). L'hygiène du numérique fonctionne exactement pareil : dire « utilise moins ton téléphone » sans savoir quelle habitude pose problème, c'est comme donner un médicament sans diagnostic. La démarche en cinq étapes que nous allons construire est l'équivalent, pour ta vie numérique, de la consultation médicale.

\begin{definition}[Démarche d'hygiène numérique]
Une \cle{démarche d'hygiène numérique} est une méthode organisée en étapes qui permet, face à une situation de vie impliquant les outils numériques, de : recueillir les faits, identifier les risques, proposer des solutions adaptées, justifier ses choix à l'aide des notions du cours, puis évaluer les résultats et ajuster si nécessaire.
\end{definition}

\courssub{Étape 1 — Observer et diagnostiquer : recueillir les faits}

Toute démarche commence par l'\cle{observation}. Il s'agit de décrire la situation \emph{telle qu'elle est}, sans juger et sans encore proposer quoi que ce soit. Les faits à recueillir sont de quatre types :

\begin{itemize}
  \item le \cle{temps d'écran} : combien d'heures par jour, sur quels appareils (téléphone, télévision, ordinateur), à quels moments de la journée ;
  \item les \cle{habitudes} : écran au lit, téléphone pendant les repas, notifications consultées en plein devoir, jeux en ligne tard le soir ;
  \item les \cle{symptômes} : maux de tête, yeux qui piquent, douleurs au cou ou au dos, fatigue, difficultés d'endormissement, baisse des notes ;
  \item les \cle{conflits et effets sociaux} : disputes avec les parents à cause du téléphone, isolement, comparaison aux autres sur les réseaux sociaux, cyberharcèlement éventuel.
\end{itemize}

\begin{methode}[Comment observer une situation numérique]
\begin{enumerate}
  \item Pendant quelques jours, noter chaque soir le temps passé sur chaque écran et à quelle heure on s'est couché.
  \item Noter les symptômes ressentis (fatigue, maux de tête, humeur) et les conflits éventuels.
  \item Regrouper les faits dans un petit tableau ou une liste : ce sont les \emph{données} de la démarche.
  \item Ne retenir que des faits observables (« il joue jusqu'à 23 h »), pas des jugements (« il est accro »).
\end{enumerate}
\end{methode}

\courssub{Étape 2 — Identifier les risques : relier les faits aux notions du cours}

Une fois les faits recueillis, on relie \emph{chaque habitude observée} à un \cle{risque} étudié dans les sections précédentes. C'est le diagnostic. On distingue trois familles :

\begin{itemize}
  \item les \cle{risques physiques} : fatigue visuelle (yeux secs, maux de tête), troubles musculo-squelettiques (cou, dos, poignets), troubles du sommeil causés par la lumière bleue qui retarde la production de mélatonine, sédentarité ;
  \item les \cle{risques psychologiques} : stress, anxiété, dépendance aux jeux ou aux réseaux sociaux, baisse de l'estime de soi par comparaison permanente ;
  \item les \cle{risques sociaux} : isolement, conflits familiaux, cyberharcèlement, baisse des résultats scolaires.
\end{itemize}

L'important est de faire un \emph{lien précis} : « écran au lit jusqu'à minuit » $\rightarrow$ « la lumière bleue retarde l'endormissement » $\rightarrow$ « sommeil insuffisant » $\rightarrow$ « fatigue et difficultés de concentration en classe ». Une chaîne de causes claire est la marque d'un bon diagnostic.

\courssub{Étape 3 — Proposer des solutions adaptées}

À chaque risque identifié doit correspondre une \cle{solution concrète} choisie parmi les bonnes pratiques étudiées : règle des pauses (toutes les 30 à 45 minutes, regarder au loin), distance yeux--écran d'au moins 50 cm, posture correcte (dos droit, écran à hauteur des yeux), écran éteint au moins une heure avant le coucher, téléphone hors de la chambre la nuit, plages sans écran (repas, devoirs), limitation du temps de jeu, paramétrage des notifications, signalement en cas de cyberharcèlement.

\begin{attention}[Erreur fréquente : la solution vague]
Proposer « utiliser moins le téléphone » ou « faire attention » n'est \textbf{pas} une solution acceptable : ce n'est ni concret, ni mesurable, ni justifié. Une bonne solution précise \emph{l'action exacte} (« éteindre l'écran à 21 h et laisser le téléphone dans le salon »), \emph{le moment} et \emph{la durée}. Demande-toi toujours : « Si je donne ce conseil à un camarade, saura-t-il exactement quoi faire ce soir ? »
\end{attention}

\courssub{Étape 4 — Rédiger et justifier : le schéma S-R-S-J}

C'est l'étape décisive, celle qui est évaluée au BEPC : rédiger sa démarche de façon claire et ordonnée. On suit le schéma \cle{S-R-S-J}.

\begin{methode}[Le schéma S-R-S-J pour rédiger une réponse justifiée]
\begin{enumerate}
  \item \textbf{S — Situation} : décrire les faits observés en une ou deux phrases (qui, quoi, combien de temps, quels symptômes).
  \item \textbf{R — Risques} : identifier le ou les risques en les nommant avec le vocabulaire de la leçon (fatigue visuelle, troubles du sommeil, dépendance, isolement...).
  \item \textbf{S — Solutions} : proposer des actions concrètes et précises, une par risque.
  \item \textbf{J — Justification} : expliquer \emph{pourquoi} chaque solution agit sur le risque, en citant la règle ou la notion du cours (« car la lumière bleue retarde la production de mélatonine », « car la règle des pauses repose les muscles des yeux »).
\end{enumerate}
\end{methode}

\begin{exemplebox}[Une réponse rédigée selon le schéma S-R-S-J]
« Aïcha dort mal et se réveille fatiguée (\textbf{S}) car elle utilise son téléphone au lit jusqu'à minuit : la lumière bleue de l'écran retarde son endormissement et réduit la qualité de son sommeil (\textbf{R}). Elle doit éteindre son écran au moins une heure avant le coucher et laisser le téléphone hors de la chambre, par exemple en le chargeant dans le salon (\textbf{S}), car l'absence de lumière bleue le soir restaure la production naturelle de mélatonine, l'hormone du sommeil, et permet un endormissement plus rapide (\textbf{J}). »
\end{exemplebox}

\begin{aretenir}[Critères d'une bonne justification]
Une justification est réussie si elle remplit trois conditions :
\begin{itemize}
  \item elle utilise le \textbf{vocabulaire de la leçon} (mélatonine, fatigue visuelle, posture, dépendance, cyberharcèlement...) ;
  \item elle \textbf{cite la règle appliquée} (écran éteint 1 h avant le coucher, pause toutes les 30 à 45 min, écran à 50 cm des yeux...) ;
  \item elle \textbf{explique le lien} entre la solution et le problème (le mot « car » ou « parce que » doit pouvoir relier les deux).
\end{itemize}
\end{aretenir}

\courssub{Étape 5 — Évaluer et ajuster}

Une démarche ne s'arrête pas à la rédaction : il faut vérifier que les solutions fonctionnent réellement. On prévoit donc un \cle{suivi} : chaque semaine, on compare la situation actuelle à la situation de départ (temps d'écran, heure du coucher, ressenti, résultats scolaires). Si une solution ne donne rien après deux ou trois semaines, on l'\emph{ajuste} : on revient à l'étape 3 pour choisir une autre action, ou à l'étape 1 pour vérifier que le diagnostic était bon. La démarche est donc un \emph{cycle}, pas une ligne droite.

\begin{propriete}[Principe d'organisation personnelle (admis)]
On admet qu'un changement d'habitude devient durable lorsqu'il est suivi et évalué régulièrement pendant plusieurs semaines (en général au moins trois à quatre semaines). C'est pourquoi toute démarche d'hygiène numérique comporte obligatoirement une étape d'évaluation et d'ajustement.
\end{propriete}

\begin{popfigure}
\begin{tikzpicture}[
  node distance=0.55cm and 1.1cm,
  etape/.style={rectangle, rounded corners=3pt, draw=popblue, fill=popblueL, thick, text width=3.1cm, align=center, minimum height=1.05cm, font=\small},
  fleche/.style={-{Stealth[length=2.5mm]}, thick, popdark},
  retour/.style={-{Stealth[length=2.5mm]}, thick, dashed, poporange}
]
\node[etape, align=center] (obs){\textbf{1. Observer}\\ recueillir les faits (temps d'écran, symptômes, conflits)};
\node[etape, right=of obs, align=center] (id){\textbf{2. Identifier}\\ relier chaque habitude à un risque (physique, psycho, social)};
\node[etape, right=of id, align=center] (prop){\textbf{3. Proposer}\\ choisir des solutions concrètes parmi les bonnes pratiques};
\node[etape, below=1.1cm of prop, align=center] (red){\textbf{4. Rédiger / Justifier}\\ schéma S-R-S-J avec le vocabulaire du cours};
\node[etape, left=of red, align=center] (eval){\textbf{5. Évaluer}\\ suivi hebdomadaire : ressenti, sommeil, notes};
\draw[fleche] (obs) -- (id);
\draw[fleche] (id) -- (prop);
\draw[fleche] (prop) -- (red);
\draw[fleche] (red) -- (eval);
\draw[retour] (eval.north) |- node[pos=0.75, above, font=\scriptsize, text=poporange]{si ça ne marche pas : ajuster} (obs.south);
\draw[retour] (eval.west) .. controls +(-1.2,0.8) and +(-1.2,-0.8) .. node[left, font=\scriptsize, text=poporange, align=center]{nouvelle\\solution} (prop.west);
\end{tikzpicture}
\legende{Organigramme des cinq étapes de la démarche d'hygiène numérique : les flèches en pointillés montrent les retours possibles pour ajuster la démarche.}
\end{popfigure}

\courssub{Construire un plan personnel d'hygiène numérique}

Pour passer de la théorie à l'action, on rassemble tout dans un \cle{plan personnel d'hygiène numérique} : un tableau à deux colonnes. Dans la première colonne, on note l'\emph{habitude à changer} (issue de l'observation) ; dans la seconde, l'\emph{action concrète} et sa \emph{justification} par les notions du cours. Ce tableau est un excellent support de rédaction pour le BEPC, car il oblige à relier systématiquement un problème à une solution justifiée.

\begin{popfigure}
\begin{tikzpicture}
\draw[fill=popblue, draw=popblue] (0,0) rectangle (15.2,0.9);
\node[white, font=\bfseries\small, anchor=west] at (0.25,0.45) {Habitude à changer (observée)};
\node[white, font=\bfseries\small, anchor=west] at (6.1,0.45) {Action concrète et justification};
\draw[fill=popblueL, draw=popline] (0,-1.7) rectangle (15.2,0);
\node[font=\small, anchor=north west, text width=5.4cm, align=flush left] at (0.2,-0.15) {Téléphone au lit jusqu'à minuit, réveil difficile le matin.};
\node[font=\small, anchor=north west, text width=8.7cm, align=flush left] at (6.1,-0.15) {Éteindre l'écran à 21 h et charger le téléphone dans le salon, \emph{car} la lumière bleue retarde la production de mélatonine et donc l'endormissement.};
\draw[fill=popcream, draw=popline] (0,-3.4) rectangle (15.2,-1.7);
\node[font=\small, anchor=north west, text width=5.4cm, align=flush left] at (0.2,-1.85) {Trois heures de jeux en ligne d'affilée le samedi, maux de tête et yeux qui piquent.};
\node[font=\small, anchor=north west, text width=8.7cm, align=flush left] at (6.1,-1.85) {Faire une pause de 10 min toutes les 45 min et regarder au loin, \emph{car} la règle des pauses repose les muscles des yeux et prévient la fatigue visuelle.};
\draw[popline] (6.0,0.9) -- (6.0,-3.4);
\end{tikzpicture}
\legende{Modèle de plan personnel d'hygiène numérique, avec deux lignes d'exemple remplies.}
\end{popfigure}

\begin{exoresolu}[La démarche complète appliquée à un cas concret]
Énoncé. Junior, élève de 3ème, passe environ 6 heures par jour sur son smartphone, dont 2 heures de réseaux sociaux le soir au lit. Depuis deux mois, il s'endort tard, ses notes ont baissé, il se dispute avec sa mère qui veut confisquer le téléphone, et il se sent « nul » en comparant sa vie à celle des influenceurs. Élabore et rédige une démarche d'hygiène numérique complète pour aider Junior.
\tcblower
Solution détaillée. On applique les cinq étapes et le schéma S-R-S-J.

\textbf{1. Observer (S).} Junior passe 6 h/jour sur son smartphone ; il consulte les réseaux sociaux au lit le soir ; symptômes : endormissement tardif, baisse des notes, baisse de l'estime de soi ; conflit familial avec sa mère.

\textbf{2. Identifier les risques (R).} Écran au lit le soir $\rightarrow$ troubles du sommeil (la lumière bleue retarde la mélatonine). Comparaison permanente aux influenceurs $\rightarrow{ }$risque psychologique : baisse de l'estime de soi. Six heures quotidiennes et conflit avec sa mère $\rightarrow$ risque social : tensions familiales et risque de dépendance. Sommeil réduit $\rightarrow$ fatigue $\rightarrow$ baisse des résultats scolaires.

\textbf{3. Proposer des solutions (S).} (a) Éteindre le téléphone une heure avant le coucher et le laisser hors de la chambre. (b) Limiter les réseaux sociaux à 30 min par jour, en journée, en désactivant les notifications. (c) Fixer des plages sans écran pendant les devoirs et les repas. (d) Dialoguer avec sa mère pour fixer ensemble des règles plutôt que subir une confiscation.

\textbf{4. Justifier (J).} (a) L'absence de lumière bleue le soir restaure la production de mélatonine : Junior s'endormira plus tôt et sera plus attentif en classe. (b) Réduire le temps passé à se comparer aux autres protège l'estime de soi, et couper les notifications diminue l'envie compulsive de consulter le téléphone. (c) Les plages sans écran améliorent la concentration, donc les notes. (d) Le dialogue transforme le conflit en règle acceptée, ce qui rend le changement durable.

\textbf{5. Évaluer et ajuster.} Junior tiendra chaque dimanche un petit bilan : heure de coucher, temps d'écran, ressenti, notes des contrôles. Si après trois semaines l'endormissement reste difficile, il avancera l'extinction de l'écran à 20 h 30 (retour à l'étape 3). Conformément au principe admis, ce suivi régulier sur plusieurs semaines rendra ses nouvelles habitudes durables.
\end{exoresolu}

\begin{exercice}[À toi de jouer]
Mariam garde les yeux rivés sur l'ordinateur du cybercafé familial 4 heures d'affilée chaque soir pour ses exposés et ses jeux ; elle se plaint de maux de tête et de douleurs au cou, et néglige ses révisions. Rédige pour elle une démarche complète en suivant le schéma S-R-S-J, puis présente ton plan sous forme de tableau à deux colonnes (habitude à changer / action concrète et justification) avec au moins trois lignes.
\end{exercice}

\begin{corrige}[Exercice 1]
\textbf{S} : 4 h d'écran d'affilée chaque soir, maux de tête, douleurs au cou, révisions négligées. \textbf{R} : fatigue visuelle et troubles musculo-squelettiques (posture prolongée), baisse des résultats scolaires. \textbf{S + J} (tableau) : (1) 4 h sans pause $\rightarrow$ pause de 10 min toutes les 45 min en regardant au loin, \emph{car} cela repose les yeux et prévient la fatigue visuelle ; (2) douleurs au cou $\rightarrow$ placer l'écran à hauteur des yeux, dos droit, à au moins 50 cm, \emph{car} une bonne posture évite les tensions du cou et du dos ; (3) révisions négligées $\rightarrow$ réserver 1 h de devoirs sans écran avant les jeux, \emph{car} les plages sans écran améliorent la concentration. \textbf{Évaluation} : bilan chaque semaine (douleurs, notes) et ajustement si besoin.
\end{corrige}

\begin{savaistu}[Le bien-être numérique, une affaire publique]
Dans plusieurs pays, des campagnes officielles encouragent les familles à fixer des « règles d'écran » à la maison, et certains établissements scolaires interdisent le téléphone portable en classe pour protéger la concentration des élèves. Au Cameroun, la sensibilisation à l'usage responsable du numérique fait partie des objectifs de l'éducation aux médias : savoir élaborer et justifier une démarche d'hygiène numérique, comme tu viens de l'apprendre, est déjà une compétence de citoyen numérique.
\end{savaistu}

\begin{aretenir}[L'essentiel de la section]
\begin{itemize}
  \item Une démarche d'hygiène numérique comporte cinq étapes : \cle{Observer} $\rightarrow$ \cle{Identifier} les risques $\rightarrow$ \cle{Proposer} des solutions $\rightarrow$ \cle{Rédiger et justifier} $\rightarrow$ \cle{Évaluer} et ajuster.
  \item La rédaction suit le schéma \cle{S-R-S-J} : Situation, Risques, Solutions, Justification.
  \item Une solution acceptable est \emph{concrète}, \emph{précise} et \emph{justifiée} par une notion du cours ; une solution vague est refusée.
  \item Le \cle{plan personnel d'hygiène numérique} (tableau habitude / action justifiée) est l'outil qui transforme le cours en changement durable, grâce au suivi régulier sur plusieurs semaines.
\end{itemize}
\end{aretenir}

\newpage

\coursec{Activité d'intégration}

\coursec{Gestion de la santé et du bien-être social}

\courssub{Objectifs de l'activité}

À l'issue de cette activité d'intégration, tu seras capable de :
\begin{itemize}
\item mobiliser les notions d'\cle{hygiène du numérique} vues dans la leçon (règle 20-20-20, règle 60/60, coupure des écrans avant le coucher, test des 4 questions avant publication, réflexes anti-cyberharcèlement) pour analyser une situation de vie courante ;
\item effectuer des calculs simples (temps d'écran hebdomadaire et annuel) et les comparer à des recommandations de santé ;
\item identifier et classer des risques en trois catégories : physiques, psychologiques et sociaux ;
\item construire un plan d'action concret et justifié, puis rédiger une conclusion argumentée, selon la démarche \cle{S-R-S-J} (Situation -- Risques -- Solutions -- Justification).
\end{itemize}

\courssub{Situation}

Aïcha est une élève de 3ème au collège bilingue de Ngaoundéré. Elle prépare son BEPC. Ses parents, commerçants au marché central, lui ont offert un smartphone pour qu'elle puisse faire ses recherches scolaires et communiquer avec ses camarades de classe. Voici la fiche « Une semaine d'Aïcha », établie avec l'aide de son professeur d'informatique.

\textbf{Relevé du temps d'écran d'Aïcha :}
\begin{itemize}
\item temps d'écran total : \textbf{5 heures par jour} (réseaux sociaux : 2 h ; vidéos : 1 h 30 ; jeux : 1 h ; travaux scolaires : 30 min) ;
\item dont \textbf{2 heures après 22 h}, allongée dans son lit, dans le noir ;
\item Aïcha ne fait \textbf{aucune pause} pendant ses sessions d'écran.
\end{itemize}

\textbf{Symptômes observés depuis deux mois :}
\begin{itemize}
\item yeux fatigués, secs et qui piquent en fin de journée ;
\item réveils nocturnes fréquents et difficultés à s'endormir ;
\item douleurs au dos et à la nuque le matin ;
\item maux de tête en classe.
\end{itemize}

\textbf{Incident survenu la semaine dernière :} une camarade a publié dans le groupe WhatsApp de la classe une photo d'Aïcha prise à la cantine, sans son autorisation. Plusieurs élèves ont laissé des commentaires moqueurs et des surnoms blessants. Aïcha n'ose plus aller au collège, elle pleure le soir et n'a prévenu ni ses parents ni un adulte de l'établissement.

\textbf{Conséquences scolaires :} ses notes ont chuté (moyenne passée de 14/20 à 10,5/20), elle s'endort parfois en cours et n'a pas rendu ses deux derniers devoirs.

Le professeur d'informatique confie à ta classe une mission : \textbf{établir un diagnostic complet de la situation d'Aïcha et lui proposer un plan d'hygiène numérique justifié}, qu'elle pourra appliquer dès cette semaine.

\courssub{Consignes}

Par groupes de 4 élèves, traite les tâches suivantes dans l'ordre. Rédige tes réponses sur une feuille qui servira à la restitution orale.

\begin{enumerate}
\item \textbf{Calcule} le temps d'écran hebdomadaire d'Aïcha (7 jours), puis son temps d'écran annuel (365 jours). Convertis le résultat annuel en nombre de journées complètes de 24 h. Que représente cette durée par rapport à une année ?
\item \textbf{Compare} ce temps aux recommandations vues en cours (pas d'écran au moins 1 h avant le coucher, pauses régulières, temps de loisir sur écran raisonnable). Repère dans la fiche toutes les habitudes d'Aïcha qui ne respectent pas ces recommandations.
\item \textbf{Identifie et classe} dans un tableau à trois colonnes (risques physiques / risques psychologiques / risques sociaux) tous les problèmes que vit Aïcha. Pour chaque problème, cite l'habitude numérique qui en est la cause.
\item \textbf{Construis un plan d'hygiène numérique en 5 actions concrètes}. Chaque action doit être justifiée par une règle précise du cours : règle 20-20-20, règle 60/60, coupure des écrans avant le coucher, test des 4 questions avant de publier, réflexes face au cyberharcèlement.
\item \textbf{Rédige une conclusion argumentée} de 8 à 10 lignes, adressée à Aïcha, qui résume le diagnostic, présente le plan d'action et explique pourquoi ces changements amélioreront sa santé, son bien-être et ses résultats scolaires.
\end{enumerate}

\courssub{Ressources à mobiliser}

\begin{aretenir}[Rappels du cours utiles pour la mission]
\begin{itemize}
\item \textbf{Règle 20-20-20 :} toutes les 20 minutes d'écran, regarder un objet situé à environ 20 pieds (6 mètres) pendant au moins 20 secondes, pour reposer les yeux.
\item \textbf{Règle 60/60 :} régler le volume d'un casque ou d'écouteurs à 60 \% du maximum et ne pas dépasser 60 minutes d'écoute continue.
\item \textbf{Coupure avant le coucher :} éteindre les écrans au moins 1 heure avant de dormir, car la lumière bleue perturbe le sommeil.
\item \textbf{Test des 4 questions avant de publier :} est-ce vrai ? est-ce utile ? est-ce gentil ? ai-je l'autorisation des personnes concernées ?
\item \textbf{Réflexes face au cyberharcèlement :} ne pas répondre, conserver les preuves (captures d'écran), bloquer le harceleur, en parler à un adulte de confiance, signaler le contenu.
\item \textbf{Posture :} écran à hauteur des yeux, dos droit, pieds au sol, pauses régulières pour se lever et s'étirer.
\item \textbf{Calcul :} temps hebdomadaire $=$ temps journalier $\times 7$ ; temps annuel $=$ temps journalier $\times 365$.
\end{itemize}
\end{aretenir}

\courssub{Démarche et corrigé commenté}

\begin{exoresolu}[Diagnostic et plan d'hygiène numérique pour Aïcha]
Énoncé : à partir de la fiche « Une semaine d'Aïcha », calculer son temps d'écran, classer les risques qu'elle encourt, proposer un plan de 5 actions justifiées et rédiger une conclusion argumentée.
\tcblower
\textbf{Étape 1 : calcul du temps d'écran (Situation).}

Temps hebdomadaire :
\[ 5 \text{ h/jour} \times 7 \text{ jours} = 35 \text{ heures par semaine.} \]

Temps annuel :
\[ 5 \text{ h/jour} \times 365 \text{ jours} = 1825 \text{ heures par an.} \]

Conversion en journées complètes de 24 h :
\[ \frac{1825}{24} \approx 76 \text{ journées complètes.} \]

Aïcha passe donc l'équivalent d'environ \textbf{76 jours et nuits sans interruption} par an devant un écran, soit plus de deux mois et demi d'une année. De plus, 2 h de ses 5 h quotidiennes ont lieu après 22 h, dans le lit et dans le noir : c'est la pire configuration pour le sommeil et pour les yeux.

\textbf{Étape 2 : comparaison aux recommandations.}

Aïcha ne respecte aucune des recommandations du cours : aucune pause (règle 20-20-20 non appliquée), écrans jusqu'à minuit (coupure avant le coucher non respectée), écran très proche des yeux dans le noir, et publication d'une photo sans autorisation dans son entourage (test des 4 questions non appliqué).

\textbf{Étape 3 : classification des risques (Risques).}

\begin{center}
\begin{tabularx}{\linewidth}{|l|X|X|}
\hline
\rowcolor{popblueL} \textbf{Catégorie} & \textbf{Problème observé} & \textbf{Cause numérique} \\
\hline
Physique & Yeux fatigués, secs, piqûres & 5 h d'écran sans pause, écran dans le noir, distance trop courte \\
\hline
Physique & Réveils nocturnes, insomnie, endormissement en classe & Écrans après 22 h (lumière bleue), non-respect de la coupure 1 h avant le coucher \\
\hline
Physique & Mal de dos et de nuque, maux de tête & Mauvaise posture allongée dans le lit, absence de pauses pour s'étirer \\
\hline
Psychologique & Stress, tristesse, pleurs, peur d'aller au collège, perte de confiance & Moqueries et surnoms blessants dans le groupe WhatsApp (cyberharcèlement) \\
\hline
Social & Photo publiée sans autorisation, isolement, rupture de communication avec adultes & Non-respect du test des 4 questions par une camarade (absence d'autorisation), silence de la victime \\
\hline
\end{tabularx}
\end{center}

\textbf{Étape 4 : plan d'hygiène numérique en 5 actions (Solutions et Justification).}

\begin{enumerate}
\item \textbf{Éteindre tous les écrans à 21 h} et poser le téléphone hors de la chambre (mode avion ou éteint). \emph{Justification :} la coupure des écrans au moins 1 h avant le coucher supprime l'effet de la lumière bleue sur la mélatonine et permet de retrouver un sommeil réparateur.
\item \textbf{Appliquer la règle 20-20-20} pendant les devoirs sur écran (30 min/jour). \emph{Justification :} regarder au loin (6 m) pendant 20 s toutes les 20 min repose les muscles ciliaires et fait disparaître la fatigue visuelle (yeux secs, picotements).
\item \textbf{S'installer correctement pour travailler :} assise à une table, dos droit, écran à hauteur des yeux, pieds au sol, et se lever 5 minutes toutes les heures. \emph{Justification :} une bonne posture et des pauses régulières préviennent les douleurs musculo-squelettiques (dos, nuque) et les maux de tête de tension.
\item \textbf{Limiter le temps de loisir sur écran à 2 h par jour} (minuteur activé), en réservant l'écran en priorité aux travaux scolaires. \emph{Justification :} un usage raisonnable et maîtrisé laisse du temps au sommeil, aux devoirs et aux activités hors écran (sport, amis, famille), réduisant le temps total à un niveau compatible avec la santé.
\item \textbf{Face au cyberharcèlement :} ne pas répondre aux moqueries, faire des captures d'écran comme preuves, bloquer les comptes moqueurs, signaler le contenu sur la plateforme et en parler immédiatement à ses parents et au professeur principal ; exiger le retrait de la photo publiée sans autorisation. \emph{Justification :} ce sont les réflexes anti-cyberharcèlement du cours ; le droit à l'image protège Aïcha : nul ne peut publier sa photo sans son accord (test des 4 questions : « ai-je l'autorisation ? »).
\end{enumerate}

\textbf{Étape 5 : conclusion argumentée (exemple de rédaction).}

Aïcha, ton diagnostic montre que tes problèmes de santé ne sont pas une fatalité : ils viennent directement de tes habitudes numériques. Cinq heures d'écran par jour, dont deux après 22 h dans le noir, expliquent ta fatigue des yeux, tes nuits perturbées et tes douleurs. La photo publiée sans ton autorisation et les moqueries qui ont suivi ont ajouté une souffrance psychologique et sociale qui t'a empêchée de travailler sereinement. En appliquant les cinq actions de ce plan — coupure des écrans à 21 h, règle 20-20-20, bonne posture, temps de loisir limité à 2 h et réflexes anti-cyberharcèlement — tu retrouveras progressivement un bon sommeil, des yeux reposés et un dos sans douleur. Ton téléphone redeviendra un outil utile pour tes études au lieu d'être une source de fatigue. Parler de ce que tu vis à un adulte est un acte de courage, pas de faiblesse. Avec un usage responsable du numérique, ta santé, ton moral et tes résultats au BEPC s'amélioreront ensemble.
\end{exoresolu}

\courssub{Bilan de l'activité}

Cette activité t'a permis de mettre en évidence plusieurs idées essentielles de la leçon :
\begin{itemize}
\item l'hygiène du numérique n'est pas un sujet abstrait : des habitudes précises et mesurables (5 h/jour, écrans après 22 h) produisent des effets concrets et observables sur le corps (yeux, sommeil, dos) ;
\item les risques du numérique sont de trois natures — physiques, psychologiques et sociales — et ils se renforcent mutuellement : la fatigue favorise la baisse des notes, le harcèlement aggrave l'isolement ;
\item chaque risque possède une parade simple et applicable immédiatement : règle 20-20-20, coupure avant le coucher, bonne posture, test des 4 questions, réflexes anti-cyberharcèlement ;
\item la démarche \cle{S-R-S-J} (Situation -- Risques -- Solutions -- Justification) est une méthode complète et transférable : tu peux l'utiliser pour analyser ta propre pratique numérique ou conseiller un proche ;
\item enfin, rédiger une conclusion argumentée t'a entraîné à \textbf{rédiger et justifier une démarche}, compétence exigée au BEPC : une réponse complète ne se contente pas de proposer des solutions, elle les relie explicitement aux règles du cours et aux données de la situation.
\end{itemize}

\begin{attention}[Erreurs fréquentes à éviter dans ce type d'activité]
\begin{itemize}
\item Proposer des actions vagues (« utiliser moins son téléphone ») sans règle précise ni justification : chaque action doit être reliée à une règle du cours.
\item Oublier de classer les risques : un même problème (ex. la baisse des notes) est souvent une \emph{conséquence} et non un risque numérique direct.
\item Conseiller à la victime de harcèlement de « répondre aux moqueries » : le premier réflexe est de ne pas répondre, de garder les preuves et d'alerter un adulte.
\end{itemize}
\end{attention}

\newpage

\coursec{Méthodes et exercices résolus}

\coursec{Gestion de la santé et du bien-être social}

\courssub{Notion d'hygiène du numérique}

\begin{definition}[Hygiène du numérique]
L'\cle{hygiène du numérique} (ou hygiène informatique) désigne l'ensemble des règles, habitudes et bonnes pratiques qui permettent de préserver sa \textbf{santé physique}, son \textbf{équilibre psychologique} et ses \textbf{relations sociales} lors de l'utilisation des outils numériques (ordinateur, smartphone, tablette, console de jeux). Elle concerne aussi bien l'ergonomie du poste de travail que la gestion du temps d'écran et le comportement en ligne.
\end{definition}

\begin{aretenir}[Pourquoi une hygiène du numérique ?]
L'omniprésence des écrans dans la vie quotidienne (école, travail, loisirs) expose à des risques réels : fatigue oculaire, douleurs musculaires, troubles du sommeil, stress, addiction, isolement, cyberharcèlement. Adopter une hygiène numérique, ce n'est pas refuser la technologie, c'est \textbf{l'utiliser de manière raisonnée, sécurisée et équilibrée} pour qu'elle reste un outil au service de notre bien-être.
\end{aretenir}

\courssub{Risques physiques liés aux écrans}

\begin{propriete}[Fatigue visuelle (asthénopie)]
L'exposition prolongée à un écran, surtout dans une pièce mal éclairée ou avec un contraste trop fort, fatigue les muscles oculaires.
\begin{itemize}
\item \textbf{Symptômes :} yeux rouges, secs ou larmoyants, picotements, vision floue temporaire, maux de tête (souvent frontaux).
\item \textbf{Causes :} clignement des paupières réduit (3 à 4 fois/minute au lieu de 12-15), lumière bleue, scintillement imperceptible, mauvaise distance de lecture.
\end{itemize}
\end{propriete}

\begin{propriete}[Troubles musculo-squelettiques (TMS)]
Une posture statique et inadaptée pendant des heures sollicite excessivement muscles, tendons et nerfs.
\begin{itemize}
\item \textbf{Zones touchées :} cou (torticolis, cervicalgies), dos (lombalgies, dorsalgie), épaules, poignets (syndrome du canal carpien), pouces (\emph{texting thumb} ou tendinite de De Quervain).
\item \textbf{Facteurs aggravants :} chaise sans dossier, écran trop bas/haut, clavier/souris mal positionnés, utilisation intensive du smartphone pouces seulement.
\end{itemize}
\end{propriete}

\begin{propriete}[Troubles du sommeil et rythme circadien]
La \cle{lumière bleue} (longueur d'onde ~460-480 nm) émise massivement par les écrans LED inhibe la sécrétion de \textbf{mélatonine} (hormone du sommeil).
\begin{itemize}
\item \textbf{Conséquences :} difficulté d'endormissement, sommeil moins réparateur, réveils nocturnes, somnolence diurne, dérèglement de l'horloge biologique.
\item \textbf{Aggravant :} usage de l'écran dans le lit, notifications sonores/lumineuses la nuit, contenu excitant (jeux, réseaux sociaux) avant le coucher.
\end{itemize}
\end{propriete}

\begin{methode}[La règle des 20-20-20 pour préserver la vue]
Pour limiter la fatigue visuelle lors d'un travail prolongé sur écran :
\begin{enumerate}
\item Toutes les \textbf{20 minutes},
\item Regarder un objet situé à au moins \textbf{20 pieds (environ 6 mètres)},
\item Pendant \textbf{20 secondes}.
\end{enumerate}
\emph{Astuce :} profite-en pour cligner des yeux délibérément 10 fois de suite pour réhydrater la cornée.
\end{methode}

\courssub{Risques psychologiques et sociaux}

\begin{definition}[Addiction aux écrans (cyberdépendance)]
Usage excessif, compulsif et incontrôlé des outils numériques (jeux vidéo, réseaux sociaux, streaming, navigation) qui entraîne une \textbf{perte de contrôle} (impossibilité de s'arrêter), un \textbf{sevrage} (irritabilité, anxiété sans écran) et des \textbf{conséquences négatives} sur la scolarité, le sommeil, la vie familiale et sociale.
\end{definition}

\begin{definition}[Cyberharcèlement]
Forme de harcèlement qui utilise les moyens numériques (réseaux sociaux, messageries, SMS, forums, jeux en ligne) pour \textbf{intimider, humilier, menacer ou exclure} une personne de manière \textbf{répétée} et \textbf{intentionnelle}. Exemples : rumeurs, photos détournées, insultes, usurpation d'identité, exclusion de groupes.
\end{definition}

\begin{propriete}[Autres risques psychologiques et sociaux]
\begin{itemize}
\item \textbf{Isolement social :} substitution des relations réelles par des interactions virtuelles, repli sur soi.
\item \textbf{Image de soi et estime :} comparaison permanente avec des vies « parfaites » sur Instagram/TikTok, dysmorphie corporelle (filtres), dépendance aux \emph{likes}.
\item \textbf{Empreinte numérique et réputation :} publications impulsives (photos, commentaires) qui restent accessibles des années plus tard (futur employeur, administration).
\item \textbf{Exposition à des contenus inappropriés :} violence, pornographie, fake news, radicalisation, arnaques (phishing, sextorsion).
\end{itemize}
\end{propriete}

\courssub{Les bonnes pratiques d'hygiène numérique}

\begin{propriete}[Ergonomie du poste de travail fixe (ordinateur de bureau)]
\begin{center}
\begin{tabularx}{\linewidth}{|l|X|}
\hline
\rowcolor{popblueL} \textbf{Élément} & \textbf{Règle d'or} \\
\hline
\textbf{Distance écran-yeux} & 50 à 70 cm (longueur du bras tendu) \\
\hline
\textbf{Hauteur écran} & Bord supérieur à \textbf{hauteur des yeux} ou légèrement en dessous (regard baissé de 10-20°) \\
\hline
\textbf{Chaise} & Dossier réglable, soutien lombaire, hauteur : pieds à plat au sol, cuisses horizontales \\
\hline
\textbf{Clavier / Souris} & Avant-bras horizontaux, poignets \textbf{neutres} (non cassés vers le haut), souris près du clavier \\
\hline
\textbf{Éclairage} & Lumière ambiante suffisante, \textbf{pas de reflet} sur l'écran (écran perpendiculaire à la fenêtre), éviter le travail dans le noir \\
\hline
\textbf{Pauses} & 5 à 10 min toutes les heures (se lever, marcher, étirements) + règle 20-20-20 \\
\hline
\end{tabularx}
\end{center}
\end{propriete}

\begin{propriete}[Bonnes pratiques sur smartphone / tablette (mobilité)]
\begin{itemize}
\item \textbf{Posture :} tenir l'appareil à hauteur du visage (bras relevés) plutôt que de baisser la tête (\emph{text neck}) ; alterner les mains et les doigts (pouce + index).
\item \textbf{Durée :} limiter les sessions continues à 20-30 min ; utiliser les outils de « Bien-être numérique » / « Temps d'écran » (Android/iOS) pour fixer des limites quotidiennes par application.
\item \textbf{Nuit :} activer le \textbf{mode « Confort visuel » / « Night Shift » / « Lumière nocturne »} (réduction lumière bleue) dès le coucher du soleil ; \textbf{mode avion} ou hors de la chambre pour dormir.
\item \textbf{Éclairage :} ne jamais utiliser l'écran dans le noir total ; activer la luminosité adaptative.
\end{itemize}
\end{propriete}

\begin{aretenir}[Règle des « 3 Pas » pour l'équilibre numérique]
\begin{enumerate}
\item \textbf{Pas d'écran} au réveil (premiers 30 min) ni au coucher (dernière heure).
\item \textbf{Pas d'écran} pendant les repas (moment de partage familial).
\item \textbf{Pas d'écran} dans la chambre la nuit (chargeur dans le salon).
\end{enumerate}
\end{aretenir}

\courssub{Usage responsable des réseaux sociaux et bien-être social}

\begin{methode}[Netiquette : les 5 règles d'or de la communication en ligne]
\begin{enumerate}
\item \textbf{Respect :} on n'écrit pas en ligne ce qu'on n'oserait pas dire en face. Pas d'insultes, de moqueries, de propos haineux (racisme, sexisme, homophobie).
\item \textbf{Vie privée :} on ne publie \textbf{jamais} de photo, vidéo ou information d'autrui sans son accord explicite (droit à l'image). On vérifie les paramètres de confidentialité (profil privé, liste d'amis restreinte).
\item \textbf{Réflexion :} avant de publier, on se pose 3 questions : \emph{Est-ce vrai ? Est-ce utile ? Est-ce bienveillant ?} (Règle des 3 sieves de Socrate adaptée).
\item \textbf{Empreinte numérique :} Internet n'oublie pas. Une blague à 14 ans peut nuire à une candidature à 22 ans. On gère sa \cle{e-réputation}.
\item \textbf{Esprit critique :} on vérifie l'info avant de partager (source, date, recoupement) pour ne pas propager de \emph{fake news}.
\end{enumerate}
\end{methode}

\begin{methode}[Réagir face au cyberharcèlement : la méthode « N.B.S.A. »]
\begin{enumerate}
\item \textbf{N} - \textbf{Ne pas répondre} : la réponse nourrit le harceleur et peut te retourner contre toi.
\item \textbf{B} - \textbf{Bloquer} l'auteur sur la plateforme (WhatsApp, Facebook, Instagram, TikTok, jeux en ligne).
\item \textbf{S} - \textbf{Signaler} le contenu via les outils de la plateforme (bouton « Signaler » / « Report ») et conserver les preuves (captures d'écran avec date, heure, pseudo).
\item \textbf{A} - \textbf{Alterter un adulte de confiance} : parents, professeur principal, CPE, conseiller d'orientation, ou appeler le \textbf{116 111} (numéro vert « Enfance en danger » au Cameroun) / contacter l'ANTIC (Agence Nationale des Technologies de l'Information et de la Communication). Le cyberharcèlement est un délit puni par la loi camerounaise (Loi n°2010/012 du 21 décembre 2010 relative à la cybercriminalité).
\end{enumerate}
\end{methode}

\begin{savaistu}[Le Cameroun et la protection des mineurs en ligne]
L'\cle{ANTIC} (Agence Nationale des Technologies de l'Information et de la Communication) est l'organe de régulation et de sécurité des systèmes d'information au Cameroun. Elle gère le \textbf{CERT-CM} (Computer Emergency Response Team) qui traite les incidents de sécurité, dont le cyberharcèlement et la sextorsion. La loi de 2010 prévoit des peines de prison et des amendes pour le harcèlement par voie électronique (Art. 75-76). En 2023, l'ANTIC a lancé la campagne « \emph{Stop Cyberharcèlement } » dans les établissements scolaires.
\end{savaistu}

\courssub{Appliquer les notions de l'unité à des situations concrètes de la vie courante}

\begin{methode}[Analyser une situation de vie et identifier les risques numériques]
Face à une situation concrète (un élève qui joue toute la nuit, un commerçant qui travaille 10 heures devant son écran, un adolescent harcelé sur WhatsApp), voici la démarche à suivre :
\begin{enumerate}
\item \textbf{Lire attentivement la situation} et souligner les éléments importants : durée d'utilisation, posture, moment de la journée, type d'activité (jeu, réseaux sociaux, travail), conséquences décrites (fatigue, douleurs, tristesse, baisse des notes).
\item \textbf{Classer les risques} en trois catégories : risques \cle{physiques} (yeux, dos, poignets, sommeil), risques \cle{psychologiques} (stress, addiction, isolement, anxiété) et risques \cle{sociaux} (cyberharcèlement, atteinte à la réputation, conflits familiaux, isolement relationnel).
\item \textbf{Nommer précisément chaque risque} avec le vocabulaire de la leçon : \cle{fatigue visuelle}, \cle{troubles musculo-squelettiques (TMS)}, \cle{troubles du sommeil}, \cle{addiction}, \cle{cyberharcèlement}, \cle{e-réputation}.
\item \textbf{Proposer une bonne pratique adaptée} à chaque risque identifié (règle des 20-20-20, posture correcte, limitation du temps d'écran via contrôle parental/bien-être numérique, arrêt écrans 1h avant coucher, blocage et signalement en cas de harcèlement, dialogue familial).
\item \textbf{Conclure} par une phrase qui relie la situation à la notion d'\cle{hygiène du numérique}.
\end{enumerate}
\textbf{Réflexe :} un risque identifié doit toujours être accompagné d'une solution concrète. Une réponse qui ne fait que décrire le problème sans proposer de solution est incomplète.
\end{methode}

\begin{exoresolu}[La nuit blanche de Tchoupo]
Tchoupo, élève en classe de 3ème à Yaoundé, joue chaque soir à des jeux vidéo sur son smartphone jusqu'à 2 heures du matin. Depuis quelques semaines, il s'endort en classe, a les yeux rouges et se plaint de maux de tête. Ses résultats scolaires baissent.
\begin{enumerate}
\item Identifie deux risques physiques et un risque psychologique auxquels Tchoupo s'expose.
\item Propose trois conseils d'hygiène numérique pour l'aider.
\end{enumerate}
\tcblower
\textbf{Solution détaillée}
\begin{enumerate}
\item \textbf{Analyse de la situation.} Les éléments importants sont : l'utilisation du smartphone \emph{la nuit} (jusqu'à 2 h du matin), la \emph{durée prolongée} devant l'écran, et les conséquences observées (sommeil en classe, yeux rouges, maux de tête, baisse des résultats).

\textbf{Risques physiques :}
\begin{itemize}
\item \cle{Troubles du sommeil} : la lumière bleue de l'écran, la nuit, perturbe l'endormissement et réduit la durée du sommeil. Cela explique qu'il s'endorme en classe.
\item \cle{Fatigue visuelle} : les yeux rouges et les maux de tête sont les signes classiques d'une exposition prolongée à l'écran, surtout dans le noir où le contraste est violent.
\end{itemize}
\textbf{Risque psychologique :}
\begin{itemize}
\item \cle{Addiction aux jeux vidéo} : le fait de jouer chaque nuit jusqu'à 2 h du matin, au détriment du sommeil et des études, traduit un usage excessif et difficilement contrôlable (perte de contrôle, conséquences négatives).
\end{itemize}

\item \textbf{Conseils d'hygiène numérique :}
\begin{itemize}
\item \textbf{Conseil 1 :} arrêter les écrans au moins \textbf{une heure avant de se coucher} et ne plus jouer la nuit, afin de retrouver un sommeil réparateur (règle du « Pas d'écran au coucher »).
\item \textbf{Conseil 2 :} limiter le temps de jeu à une durée raisonnable (par exemple 1 heure par jour) en utilisant la fonction « Bien-être numérique » du téléphone, et faire des \textbf{pauses régulières} : appliquer la règle des 20-20-20 (toutes les 20 minutes, regarder un objet à environ 6 mètres pendant 20 secondes).
\item \textbf{Conseil 3 :} ne pas jouer dans le noir complet : utiliser l'écran dans une pièce éclairée, activer le mode « Confort visuel » (filtre lumière bleue) et réduire la luminosité pour ménager les yeux.
\end{itemize}
\end{enumerate}
\textbf{Conclusion :} Tchoupo souffre de troubles du sommeil, de fatigue visuelle et d'un début d'addiction aux jeux vidéo ; en appliquant les règles d'hygiène du numérique (limitation du temps d'écran, pauses, arrêt des écrans avant le coucher, éclairage adapté), il retrouvera un meilleur équilibre entre sa santé et ses activités numériques.
\end{exoresolu}

\begin{methode}[Diagnostiquer et corriger un mauvais poste de travail]
Pour vérifier qu'un poste de travail respecte l'hygiène du numérique, procède par points de contrôle :
\begin{enumerate}
\item \textbf{L'écran} : il doit être à une distance de 50 à 70 cm des yeux, le haut de l'écran à hauteur des yeux ou légèrement en dessous.
\item \textbf{Le dos} : le dos doit être droit, appuyé contre le dossier de la chaise ; les pieds posent à plat sur le sol (ou sur un repose-pieds).
\item \textbf{Les bras et poignets} : les avant-bras reposent à l'horizontale, les poignets ne sont pas cassés vers le haut (utiliser un repose-poignet si besoin).
\item \textbf{L'éclairage} : la pièce doit être bien éclairée, sans reflet sur l'écran ; l'écran ne doit pas être face à une fenêtre (éblouissement) ni dos à la fenêtre (reflet).
\item \textbf{Les pauses} : prévoir une pause de 5 à 10 minutes toutes les heures (se lever, bouger) et appliquer la règle des 20-20-20 pour les yeux.
\end{enumerate}
Pour chaque défaut repéré, indique le risque (TMS, fatigue visuelle) et la correction à apporter.
\end{methode}

\begin{exoresolu}[Le cybercafé de Mme Etoundi]
Mme Etoundi gère un petit cybercafé à Douala. Ses clients s'assoient sur des tabourets sans dossier, les écrans sont posés très bas sur les tables, et la salle est mal éclairée. Plusieurs clients se plaignent de douleurs au cou et au dos après quelques heures.
\begin{enumerate}
\item Releve trois défauts d'ergonomie dans ce cybercafé.
\item Pour chaque défaut, indique le risque pour la santé et la correction à apporter.
\end{enumerate}
\tcblower
\textbf{Solution détaillée}
\begin{enumerate}
\item \textbf{Relevé des défauts :}
\begin{itemize}
\item Défaut 1 : les clients s'assoient sur des \textbf{tabourets sans dossier}.
\item Défaut 2 : les écrans sont posés \textbf{très bas} sur les tables.
\item Défaut 3 : la salle est \textbf{mal éclairée}.
\end{itemize}
\item \textbf{Risques et corrections :}
\begin{center}
\begin{tabularx}{\linewidth}{|X|X|X|}
\hline
\rowcolor{popblueL} \textbf{Défaut} & \textbf{Risque pour la santé} & \textbf{Correction} \\
\hline
Tabourets sans dossier & Douleurs au dos, \cle{troubles musculo-squelettiques (TMS)} car le dos n'est pas soutenu, position voûtée & Utiliser des chaises ergonomiques avec dossier et soutien lombaire, le dos bien appuyé \\
\hline
Écrans posés très bas & Douleurs au cou (\cle{cervicalgies}), torticolis car la tête est penchée en permanence vers le bas & Rehausser les écrans (support ou livres) pour que le haut de l'écran soit à hauteur des yeux \\
\hline
Salle mal éclairée & \cle{Fatigue visuelle}, maux de tête, car les yeux font un effort d'adaptation permanent (contraste écran/salle) & Installer un éclairage indirect suffisant (plafonniers, lampes), éviter les reflets sur les écrans (stores aux fenêtres) \\
\hline
\end{tabularx}
\end{center}
\end{enumerate}
\textbf{Conclusion :} en installant des chaises avec dossier, en rehaussant les écrans à hauteur des yeux et en améliorant l'éclairage de la salle, Mme Etoundi protégera la santé de ses clients et appliquera les règles d'ergonomie du poste de travail.
\end{exoresolu}

\begin{methode}[Réagir face au cyberharcèlement et aux dangers des réseaux sociaux]
Si toi ou un camarade êtes victime de cyberharcèlement (moqueries, insultes, menaces, rumeurs, photos détournées sur les réseaux sociaux), applique la démarche « N.B.S.A. » :
\begin{enumerate}
\item \textbf{N} - \textbf{Ne pas répondre} aux provocations : répondre aggrave souvent la situation et peut être utilisé contre toi.
\item \textbf{B} - \textbf{Bloquer} l'auteur des messages sur le réseau social concerné (WhatsApp, Facebook, Instagram, TikTok, Messenger).
\item \textbf{S} - \textbf{Signaler} les contenus à la plateforme (bouton « Signaler » / « Report » présent sur chaque publication/profil/message).
\item \textbf{Garder les preuves} : faire des captures d'écran des messages, photos ou commentaires injurieux (avec date, heure, pseudo visible) \textbf{AVANT} de bloquer/supprimer.
\item \textbf{A} - \textbf{Alterter un adulte de confiance} : parents, enseignant, censeur, conseiller d'orientation. Le cyberharcèlement est puni par la loi camerounaise (Loi 2010/012). Numéro vert : \textbf{116 111}.
\end{enumerate}
\textbf{Réflexe préventif :} avant de publier quoi que ce soit, se poser la question : « Est-ce que j'accepterais que mes parents, mes enseignants ou mon futur employeur voient cette publication dans 5 ans ? »
\end{methode}

\begin{exoresolu}[Les moqueries contre Aïcha]
Aïcha, élève de 3ème, a publié une photo d'elle sur un groupe WhatsApp de sa classe. Depuis, certains camarades détournent sa photo, la modifient et l'accompagnent de commentaires moqueurs. Aïcha ne veut plus aller au collège.
\begin{enumerate}
\item Comment appelle-t-on ce phénomène ?
\item Décris la démarche complète qu'Aïcha doit suivre pour s'en sortir.
\end{enumerate}
\tcblower
\textbf{Solution détaillée}
\begin{enumerate}
\item Ce phénomène s'appelle le \cle{cyberharcèlement} : il s'agit de moqueries, d'insultes ou de humiliations répétées, commises par le biais des outils numériques (réseaux sociaux, messageries). Le détournement de sa photo (\emph{deepfake} artisanal ou montage) et les commentaires moqueurs en sont des formes typiques. C'est aussi une atteinte au \cle{droit à l'image}.
\item \textbf{Démarche à suivre par Aïcha (méthode N.B.S.A.) :}
\begin{enumerate}
\item \textbf{Ne pas répondre} aux moqueries : répondre ou se justifier encourage souvent les harceleurs à continuer.
\item \textbf{Bloquer} les camarades auteurs des moqueries sur WhatsApp et quitter éventuellement le groupe concerné pour se protéger.
\item \textbf{Signaler} les contenus injurieux et les profils des harceleurs à la plateforme (fonction « Signaler » de WhatsApp : Maintenir le message $\to$ Signaler).
\item \textbf{Conserver les preuves} : faire des captures d'écran des photos détournées et des commentaires, avec les dates et les noms des auteurs. Ces preuves seront utiles pour la direction du collège et éventuellement la police/ANTIC.
\item \textbf{Alterter un adulte de confiance} : ses parents, son professeur principal ou le censeur du collège. L'établissement pourra intervenir auprès des auteurs (sanctions disciplinaires), et le cyberharcèlement étant puni par la loi, une plainte est possible si les faits continuent.
\end{enumerate}
\end{enumerate}
\textbf{Conclusion :} Aïcha est victime de cyberharcèlement ; en ne répondant pas, en gardant les preuves, en bloquant et signalant les auteurs, puis en alertant un adulte, elle se protège efficacement et permet que les responsables soient sanctionnés.
\end{exoresolu}

\courssub{Rédiger et justifier une démarche, une réponse ou une conclusion}

\begin{methode}[Rédiger une réponse justifiée en trois temps : Méthode A-J-C]
Au BEPC, une réponse sur l'hygiène du numérique doit être rédigée et justifiée. Suis la méthode « A-J-C » :
\begin{enumerate}
\item \textbf{Affirmer} : donne clairement ta réponse en une phrase complète (ex. : « Cette pratique est dangereuse pour la santé. »).
\item \textbf{Justifier} : appuie ton affirmation avec un savoir du cours (nom du risque, explication du mécanisme, règle d'hygiène numérique) : « Car... », « En effet... », « Cela s'explique par... ».
\item \textbf{Conclure} : termine par une phrase qui propose une solution ou qui relie la réponse à la situation posée.
\end{enumerate}
\textbf{Réflexes de rédaction :}
\begin{itemize}
\item Rédiger des \textbf{phrases complètes} (sujet, verbe, complément), jamais de simples mots isolés ou listes à puces sans verbes.
\item Employer le \textbf{vocabulaire exact} de la leçon : hygiène du numérique, fatigue visuelle, troubles musculo-squelettiques (TMS), addiction, cyberharcèlement, lumière bleue, mélatonine, e-réputation.
\item Soigner l'orthographe, la syntaxe et la présentation : une copie claire et structurée gagne des points.
\end{itemize}
\end{methode}

\begin{exoresolu}[Justifier une règle d'hygiène numérique]
Ton petit frère te demande : « Pourquoi maman m'interdit-elle de regarder la tablette juste avant de dormir ? » Rédige une réponse complète et justifiée (5 à 8 lignes) pour lui expliquer cette règle.
\tcblower
\textbf{Solution détaillée}

\textbf{Affirmation :} Ta maman a raison de t'interdire la tablette juste avant de dormir, car utiliser un écran le soir nuit gravement à la qualité de ton sommeil.

\textbf{Justification :} En effet, les écrans diffusent une \cle{lumière bleue} qui trompe ton cerveau : celui-ci croit qu'il fait encore jour et retarde la sécrétion de \cle{mélatonine}, l'hormone qui déclenche l'endormissement. De plus, les jeux et les vidéos excitent ton esprit au lieu de le préparer au repos. Un enfant qui regarde la tablette le soir s'endort plus tard, dort moins profondément, et le lendemain il est fatigué, inattentif en classe et souvent de mauvaise humeur. C'est ce qu'on appelle les \cle{troubles du sommeil} liés aux écrans.

\textbf{Conclusion :} Pour bien dormir et être en forme, il faut donc éteindre tous les écrans au moins une heure avant le coucher et préférer une activité calme comme la lecture d'un livre papier : c'est l'une des règles essentielles de l'\cle{hygiène du numérique}.
\end{exoresolu}

\begin{methode}[Élaborer et justifier une démarche d'hygiène numérique (Charte)]
Pour construire une démarche complète d'hygiène numérique (par exemple une charte d'utilisation des écrans pour ta famille, ta classe ou ton cybercafé), procède en 5 étapes :
\begin{enumerate}
\item \textbf{Faire le diagnostic} : recenser les usages (quels appareils ? combien d'heures par jour ? à quels moments ?) et les problèmes observés (fatigue, conflits, baisse des notes, douleurs, isolement).
\item \textbf{Fixer des règles de temps} : durée maximale quotidienne d'écran de loisir, interdiction des écrans la nuit et pendant les repas, pauses régulières (20-20-20, 10 min/h).
\item \textbf{Fixer des règles d'ergonomie} : posture correcte, écran à bonne distance et à hauteur des yeux, pièce éclairée, pas d'écran dans le lit.
\item \textbf{Fixer des règles de comportement en ligne} : respect des autres (netiquette), pas de publication de photos sans autorisation (droit à l'image), réflexion avant publication (empreinte numérique), signalement des abus (cyberharcèlement).
\item \textbf{Prévoir le suivi et l'évaluation} : vérifier régulièrement (ex. chaque mois) que les règles sont respectées, en discuter, et les ajuster si besoin (évolution de l'âge, des besoins).
\end{enumerate}
Chaque règle doit être \textbf{justifiée} par le risque qu'elle prévient : c'est cette justification qui est évaluée.
\end{methode}

\begin{exoresolu}[La charte numérique de la famille Ngo Bell]
La famille Ngo Bell possède deux smartphones, une tablette et un ordinateur portable. Les enfants passent en moyenne 6 heures par jour sur les écrans, y compris pendant les repas et tard le soir. Le père souhaite établir une charte familiale d'utilisation des écrans. Propose une charte de 5 règles, chacune justifiée par le risque qu'elle prévient.
\tcblower
\textbf{Solution détaillée}

\textbf{Diagnostic :} les usages sont excessifs (6 heures par jour, bien au-dessus des recommandations de 2h loisir), envahissent les moments de vie familiale (repas) et empiètent sur le sommeil (soirée tardive).

\textbf{Charte proposée :}
\begin{enumerate}
\item \textbf{Règle 1 :} limiter l'utilisation des écrans de loisir à 2 heures par jour maximum (hors devoirs). \emph{Justification :} un temps d'écran excessif favorise la \cle{fatigue visuelle}, la \cle{sédentarité} et le risque d'\cle{addiction} aux écrans.
\item \textbf{Règle 2 :} aucun écran pendant les repas (téléphones dans un panier à l'entrée de la salle à manger). \emph{Justification :} les repas sont des moments privilégiés d'échange familial ; les écrans isolent chacun, dégradent la communication et les relations sociales au sein de la famille.
\item \textbf{Règle 3 :} éteindre tous les écrans au moins une heure avant le coucher (chargeurs branchés dans le salon, pas dans les chambres). \emph{Justification :} la \cle{lumière bleue} des écrans perturbe la sécrétion de mélatonine et provoque des \cle{troubles du sommeil} (endormissement difficile, sommeil de mauvaise qualité).
\item \textbf{Règle 4 :} faire une pause de 10 minutes toutes les heures d'écran continu et s'installer correctement (dos droit, écran à hauteur des yeux, distance 50-70 cm, pièce éclairée). \emph{Justification :} ces mesures préviennent la \cle{fatigue visuelle} (règle 20-20-20) et les \cle{troubles musculo-squelettiques (TMS)} (douleurs au dos, au cou, aux poignets).
\item \textbf{Règle 5 :} sur les réseaux sociaux, respecter les autres (pas d'insultes, pas de moqueries), ne publier aucune photo d'un proche sans son accord explicite et signaler tout message blessant ou suspect à un parent. \emph{Justification :} ces règles préviennent le \cle{cyberharcèlement}, protègent le \cle{droit à l'image}, la \cle{réputation numérique} (e-réputation) et la vie privée de chacun.
\end{enumerate}
\textbf{Suivi :} chaque dimanche soir, la famille fait le point (15 min) sur le respect de la charte, discute des difficultés et l'ajuste si nécessaire (ex. : augmenter le temps le week-end, ajouter une application éducative autorisée).

\textbf{Conclusion :} cette charte, en limitant le temps d'écran, en protégeant le sommeil, la posture, les relations familiales et la vie en ligne, met en pratique l'ensemble des règles d'hygiène du numérique au service de la santé et du bien-être de toute la famille.
\end{exoresolu}

\begin{attention}[Les 5 erreurs qui coûtent des points à l'examen (BEPC)]
\begin{enumerate}
\item \textbf{Confondre les catégories de risques.} Beaucoup d'élèves classent le cyberharcèlement parmi les risques physiques. Retiens bien : yeux, dos, poignets, sommeil = risques \emph{physiques} ; stress, addiction, isolement, anxiété = risques \emph{psychologiques} ; cyberharcèlement, atteinte à la réputation, conflits, droit à l'image = risques \emph{sociaux}.
\item \textbf{Décrire le problème sans proposer de solution.} Une question du type « Que conseilles-tu ? » ou « Propose des mesures » exige des \emph{bonnes pratiques concrètes} (pauses, distance de l'écran, blocage, signalement, dialogue). Une réponse qui ne fait que répéter l'énoncé ou décrire les symptômes ne rapporte aucun point.
\item \textbf{Répondre par des mots isolés au lieu de phrases rédigées.} Écrire « sommeil, yeux » au lieu de « Il souffre de troubles du sommeil et de fatigue visuelle » fait perdre les points de rédaction et de justification. Chaque réponse doit être une phrase complète, affirmée puis justifiée (méthode A-J-C).
\item \textbf{Citer une règle sans la justifier.} Dire « il faut faire des pauses » ne suffit pas : il faut préciser \emph{pourquoi} (prévenir la fatigue visuelle / TMS) et éventuellement \emph{comment} (règle des 20-20-20, 10 minutes par heure). La justification est le cœur du savoir-faire évalué.
\item \textbf{Conseiller de répondre au harceleur ou de se venger.} Face au cyberharcèlement, la première règle est de \emph{ne pas répondre}, puis de garder les preuves, bloquer, signaler et en parler à un adulte. Proposer de « se venger », « insulter en retour », « pirater son compte » ou « régler ça entre nous » est une erreur grave qui contredit la leçon et la loi.
\end{enumerate}
\end{attention}

\newpage

\coursec{Exercices}

\courssub{Je vérifie mes connaissances}

\begin{exercice}[QCM : les bons réflexes]
Pour chaque question, entoure la (ou les) bonne(s) réponse(s).
\begin{enumerate}
\item La règle dite « 20-20-20 » consiste à :
\begin{enumerate}
\item[a)] regarder un écran 20 minutes puis dormir 20 minutes ;
\item[b)] toutes les 20 minutes, regarder pendant 20 secondes un objet situé à environ 20 pieds (6 mètres) ;
\item[c)] passer au maximum 20 minutes par jour devant un écran.
\end{enumerate}
\item La distance recommandée entre les yeux et l'écran d'un ordinateur est d'environ :
\begin{enumerate}
\item[a)] 10 cm ; \quad b) 50 à 70 cm ; \quad c) 2 m.
\end{enumerate}
\item Avant de publier la photo d'un camarade sur un réseau social, il faut :
\begin{enumerate}
\item[a)] flouter son visage ; \quad b) demander son autorisation ; \quad c) rien, c'est libre.
\end{enumerate}
\item Une bonne hygiène du numérique comprend :
\begin{enumerate}
\item[a)] éteindre les écrans au moins 30 minutes avant de dormir ;
\item[b)] consulter son téléphone toutes les nuits vers 2 h du matin ;
\item[c)] faire des pauses régulières pendant le travail sur écran.
\end{enumerate}
\end{enumerate}
\end{exercice}

\begin{exercice}[Vrai ou faux, en justifiant]
Réponds par VRAI ou FAUX, puis justifie ta réponse en une ou deux phrases.
\begin{enumerate}
\item La lumière bleue des écrans, le soir, peut perturber l'endormissement.
\item Le cyberharcèlement n'a pas de conséquences réelles car il se passe « seulement » sur Internet.
\item Régler la luminosité de l'écran en fonction de l'éclairage de la pièce fait partie des bonnes pratiques.
\item Partager son mot de passe avec son meilleur ami est sans danger.
\end{enumerate}
\end{exercice}

\begin{exercice}[Définitions]
Définis en deux ou trois phrases chacun des termes suivants, puis donne un exemple tiré de la vie courante :
\begin{enumerate}
\item hygiène du numérique ;
\item cyberharcèlement ;
\item e-réputation ;
\item netiquette.
\end{enumerate}
\end{exercice}

\begin{exercice}[Calcul de temps d'écran]
Amina passe en moyenne 4 h 30 par jour sur son téléphone portable.
\begin{enumerate}
\item Calcule son temps d'écran hebdomadaire (7 jours).
\item Calcule son temps d'écran sur une année de 365 jours (donne le résultat en heures, puis en journées de 24 h).
\item Elle décide de réduire ce temps à 2 h par jour. Combien d'heures par semaine libère-t-elle pour ses études ?
\end{enumerate}
\end{exercice}

\courssub{Je m'entraîne}

\begin{exercice}[Bonne ou mauvaise pratique ?]
Recopie le tableau ci-dessous et classe chacun des 10 comportements suivants dans la colonne qui convient, puis corrige chaque mauvaise pratique par une action concrète :
\begin{enumerate}
\item[a)] s'asseoir le dos droit, les pieds à plat devant l'ordinateur ;
\item[b)] manger en tapant sur le clavier sans se laver les mains ;
\item[c)] faire une pause de 10 minutes après une heure d'écran ;
\item[d)] régler le haut de l'écran à hauteur des yeux ;
\item[e)] publier une photo de classe sans demander l'avis des camarades ;
\item[f)] utiliser un mot de passe différent pour chaque compte ;
\item[g)] laisser son téléphone allumé sous l'oreiller la nuit ;
\item[h)] signaler à un adulte un message de harcèlement reçu en ligne ;
\item[i)] travailler sur l'ordinateur dans une pièce bien éclairée ;
\item[j)] passer 6 heures d'affilée sur un jeu vidéo sans pause.
\end{enumerate}
\begin{center}
\begin{tabularx}{\linewidth}{|X|l|X|}
\hline
\rowcolor{popblueL} \textbf{Comportement} & \textbf{Bonne / mauvaise pratique} & \textbf{Correction si mauvaise pratique} \\
\hline
a) & & \\
\hline
\end{tabularx}
\end{center}
\end{exercice}

\begin{exercice}[Étude de cas : Moussa et son téléphone]
Moussa, élève de 3ème, consulte son téléphone toutes les nuits vers 2 h du matin pour suivre des vidéos. Le jour, il s'endort en classe, ses notes baissent et il se plaint de maux de tête.
\begin{enumerate}
\item Identifie trois risques (physiques, scolaires ou psychologiques) auxquels Moussa s'expose.
\item Propose trois solutions concrètes pour l'aider.
\item Justifie chaque solution en expliquant le bienfait attendu (démarche : Situation -- Risque -- Solution -- Justification).
\end{enumerate}
\end{exercice}

\begin{exercice}[Mon plan personnel d'hygiène numérique]
Recopie et complète le tableau suivant avec quatre de tes propres habitudes numériques à améliorer.
\begin{center}
\begin{tabularx}{\linewidth}{|c|X|X|X|}
\hline
\rowcolor{popblueL} \textbf{N°} & \textbf{Habitude à changer} & \textbf{Action concrète} & \textbf{Justification} \\
\hline
1 & & & \\
\hline
2 & & & \\
\hline
3 & & & \\
\hline
4 & & & \\
\hline
\end{tabularx}
\end{center}
\end{exercice}

\begin{exercice}[Réagir face au cyberharcèlement]
Dans un groupe de discussion de la classe, un élève diffuse un montage photo moquant le physique de Nadège. Plusieurs camarades ajoutent des commentaires blessants.
\begin{enumerate}
\item Explique pourquoi ce comportement est qualifié de cyberharcèlement.
\item Cite deux conséquences possibles pour Nadège.
\item Décris trois attitudes responsables que tu peux adopter en tant que témoin.
\end{enumerate}
\end{exercice}

\courssub{Je me prépare au Probatoire}

\begin{exercice}[Type BEPC : le club informatique sensibilise]
Le club informatique du collège de Bafoussam organise une journée « Santé et numérique ». Le président du club relève les habitudes des élèves : en moyenne, un élève passe 5 h par jour devant un écran (téléphone, télévision, ordinateur), dont 2 h après 22 h.
\begin{enumerate}
\item Définis l'expression « hygiène du numérique ».
\item Calcule le temps d'écran hebdomadaire d'un élève, puis le temps passé devant un écran après 22 h sur une semaine.
\item Cite deux risques physiques et deux risques psychologiques liés à un usage excessif des écrans.
\item Propose quatre règles d'hygiène numérique que le club peut afficher dans la salle multimédia, en justifiant chacune.
\item Rédige un slogan de sensibilisation (une phrase) pour la journée « Santé et numérique ».
\end{enumerate}
\end{exercice}

\begin{exercice}[Type BEPC : les réseaux sociaux au collège]
Kamdem, élève de 6ème, vient de créer son premier compte sur un réseau social. Il te demande conseil.
\begin{enumerate}
\item Rédige un texte de 10 à 15 lignes expliquant à Kamdem comment utiliser les réseaux sociaux de façon responsable. Ton texte devra aborder obligatoirement : la netiquette, la protection de la vie privée, l'e-réputation et le cyberharcèlement.
\item Kamdem veut publier une photo de ses camarades prise à la cantine. Indique-lui la démarche à suivre avant la publication et justifie-la.
\item Un inconnu demande à Kamdem sa localisation et son numéro de téléphone en message privé. Que doit-il faire ? Donne deux actions concrètes.
\end{enumerate}
\end{exercice}

\begin{exercice}[Type BEPC : enquête dans une classe de 3ème]
Une enquête menée dans une classe de 3ème de 50 élèves donne les résultats suivants : 35 élèves dorment avec leur téléphone dans la chambre ; 20 élèves déclarent des douleurs au cou ou aux yeux après les écrans ; 10 élèves ont déjà reçu un message moqueur ou menaçant en ligne.
\begin{enumerate}
\item Calcule le pourcentage d'élèves concernés par chacune des trois situations.
\item Pour chaque situation, identifie le risque encouru et propose une bonne pratique permettant de le réduire.
\item Le conseiller d'orientation souhaite afficher un « plan d'hygiène numérique de la classe ». Rédige ce plan sous forme d'un tableau de quatre lignes : habitude à changer, action concrète, justification.
\item Conclus en rédigeant un court paragraphe (5 à 8 lignes) expliquant pourquoi le bien-être social d'un élève dépend aussi de son comportement numérique.
\end{enumerate}
\end{exercice}

\newpage

\coursec{Corrigés des exercices}

\begin{corrige}[Exercice 1 — QCM : les bons réflexes]
\begin{enumerate}
\item \textbf{Réponse b).} La règle « 20-20-20 » consiste, toutes les 20 minutes de travail sur écran, à regarder pendant 20 secondes un objet situé à environ 20 pieds (soit environ 6 mètres). Elle permet de reposer les muscles des yeux et de limiter la fatigue visuelle.
\item \textbf{Réponse b).} La distance recommandée entre les yeux et l'écran est d'environ 50 à 70 cm (environ la longueur d'un bras). À 10 cm, les yeux fatiguent très vite ; à 2 m, on est obligé de se pencher ou de plisser les yeux.
\item \textbf{Réponse b).} Avant de publier la photo d'une personne, il faut toujours demander son autorisation : c'est le respect du droit à l'image et de la vie privée. Publier sans accord peut nuire à l'e-réputation du camarade et constituer une faute.
\item \textbf{Réponses a) et c).} Éteindre les écrans au moins 30 minutes avant de dormir favorise l'endormissement (la lumière bleue perturbe le sommeil), et faire des pauses régulières protège les yeux, le dos et la concentration. Consulter son téléphone à 2 h du matin (réponse b) est une mauvaise pratique qui détruit la qualité du sommeil.
\end{enumerate}
\end{corrige}

\begin{corrige}[Exercice 2 — Vrai ou faux, en justifiant]
\begin{enumerate}
\item \textbf{VRAI.} La lumière bleue émise par les écrans trompe le cerveau qui « croit » qu'il fait encore jour : la production de mélatonine (l'hormone du sommeil) est retardée, ce qui rend l'endormissement plus difficile.
\item \textbf{FAUX.} Le cyberharcèlement a des conséquences bien réelles : tristesse, angoisse, isolement, baisse des résultats scolaires, parfois dépression. De plus, il est puni par la loi : « sur Internet » ne signifie pas « sans importance » ni « sans sanction ».
\item \textbf{VRAI.} Un écran trop lumineux dans une pièce sombre (ou trop sombre dans une pièce claire) fatigue les yeux. Adapter la luminosité de l'écran à l'éclairage de la pièce est une bonne pratique d'hygiène visuelle.
\item \textbf{FAUX.} Un mot de passe est personnel et secret, comme le code d'une carte bancaire. Même un meilleur ami peut le perdre, le révéler par erreur ou utiliser ton compte sans ton accord ; tu restes responsable de tout ce qui est fait avec ton compte.
\end{enumerate}
\end{corrige}

\begin{corrige}[Exercice 3 — Définitions]
\begin{enumerate}
\item \textbf{Hygiène du numérique :} ensemble des règles et des bonnes habitudes d'utilisation des outils numériques (ordinateur, téléphone, tablette) qui permettent de protéger sa santé physique, son équilibre psychologique et ses relations sociales. \emph{Exemple :} faire une pause de 10 minutes après une heure de travail sur l'ordinateur.
\item \textbf{Cyberharcèlement :} ensemble de moqueries, insultes, menaces, rumeurs ou humiliations répétées contre une personne à travers les outils numériques (réseaux sociaux, SMS, groupes de discussion). \emph{Exemple :} diffuser dans le groupe WhatsApp de la classe un montage photo ridiculisant un camarade.
\item \textbf{E-réputation :} image qu'une personne laisse d'elle-même sur Internet, formée par tout ce qu'elle publie (photos, commentaires, vidéos) et par ce que les autres publient sur elle. \emph{Exemple :} un employeur qui recherche ton nom sur Internet avant de t'embaucher.
\item \textbf{Netiquette :} ensemble des règles de politesse et de bon comportement à respecter sur Internet (saluer, ne pas insulter, ne pas écrire en majuscules qui équivalent à crier, respecter les autres). \emph{Exemple :} répondre calmement et respectueusement dans un groupe de discussion, même quand on n'est pas d'accord.
\end{enumerate}
\end{corrige}

\begin{corrige}[Exercice 4 — Calcul de temps d'écran]
\begin{enumerate}
\item Temps hebdomadaire : $4\text{ h }30 = 4,5$ h par jour.
\[ 4,5 \times 7 = 31,5 \text{ h} \]
Amina passe \textbf{31 h 30 par semaine} sur son téléphone (soit plus d'une journée complète !).
\item Sur 365 jours :
\[ 4,5 \times 365 = 1642,5 \text{ h} \]
Conversion en journées de 24 h :
\[ \frac{1642,5}{24} \approx 68,4 \text{ journées} \]
Amina passe environ \textbf{1642,5 heures}, c'est-à-dire environ \textbf{68 journées entières} de 24 h par an devant son téléphone (plus de deux mois !).
\item Nouveau temps quotidien : 2 h, donc un gain de $4,5 - 2 = 2,5$ h par jour.
\[ 2,5 \times 7 = 17,5 \text{ h} \]
Elle libère \textbf{17 h 30 par semaine} pour ses études.
\end{enumerate}
\end{corrige}

\begin{corrige}[Exercice 5 — Bonne ou mauvaise pratique ?]
\begin{center}
\begin{tabularx}{\linewidth}{|X|l|X|}
\hline
\rowcolor{popblueL} \textbf{Comportement} & \textbf{Bonne / mauvaise} & \textbf{Correction si mauvaise pratique} \\
\hline
a) Dos droit, pieds à plat & Bonne & --- \\
\hline
b) Manger en tapant sans se laver les mains & Mauvaise & Se laver les mains avant et ne pas manger au clavier (miettes, microbes, dégâts sur le matériel) \\
\hline
c) Pause de 10 min après 1 h d'écran & Bonne & --- \\
\hline
d) Haut de l'écran à hauteur des yeux & Bonne & --- \\
\hline
e) Publier une photo de classe sans avis & Mauvaise & Demander l'autorisation de chaque camarade visible avant de publier \\
\hline
f) Un mot de passe différent par compte & Bonne & --- \\
\hline
g) Téléphone allumé sous l'oreiller la nuit & Mauvaise & Éteindre le téléphone ou le laisser hors de la chambre la nuit \\
\hline
h) Signaler un message de harcèlement à un adulte & Bonne & --- \\
\hline
i) Travailler dans une pièce bien éclairée & Bonne & --- \\
\hline
j) 6 h de jeu vidéo sans pause & Mauvaise & Faire une pause de 10 minutes toutes les heures et limiter la durée de jeu \\
\hline
\end{tabularx}
\end{center}
\end{corrige}

\begin{corrige}[Exercice 6 — Étude de cas : Moussa et son téléphone]
\begin{enumerate}
\item \textbf{Trois risques :}
\begin{itemize}
\item \emph{Risque physique :} manque de sommeil, maux de tête, fatigue des yeux (lumière bleue à 2 h du matin).
\item \emph{Risque scolaire :} somnolence en classe, baisse de la concentration et des notes.
\item \emph{Risque psychologique :} dépendance au téléphone, difficulté à se passer des vidéos.
\end{itemize}
\item et 3. \textbf{Démarche Situation -- Risque -- Solution -- Justification :}
\begin{center}
\begin{tabularx}{\linewidth}{|X|X|X|X|}
\hline
\rowcolor{popblueL} \textbf{Situation} & \textbf{Risque} & \textbf{Solution} & \textbf{Justification} \\
\hline
Téléphone consulté à 2 h du matin & Sommeil perturbé, fatigue & Éteindre le téléphone et le laisser hors de la chambre la nuit & Supprime la tentation et la lumière bleue ; le sommeil redevient réparateur \\
\hline
Somnolence en classe, notes en baisse & Échec scolaire & Fixer une heure de coucher régulière (avant 22 h) & Un sommeil suffisant améliore l'attention et la mémorisation en classe \\
\hline
Maux de tête, yeux fatigués & Fatigue visuelle & Limiter le temps d'écran quotidien et faire des pauses & Moins d'écran le soir = moins de fatigue des yeux et moins de maux de tête \\
\hline
\end{tabularx}
\end{center}
\end{enumerate}
\end{corrige}

\begin{corrige}[Exercice 7 — Mon plan personnel d'hygiène numérique]
Exemple de correction (les réponses personnelles des élèves peuvent varier ; vérifier la cohérence habitude -- action -- justification) :
\begin{center}
\begin{tabularx}{\linewidth}{|c|X|X|X|}
\hline
\rowcolor{popblueL} \textbf{N°} & \textbf{Habitude à changer} & \textbf{Action concrète} & \textbf{Justification} \\
\hline
1 & Garder le téléphone dans la chambre la nuit & Le laisser charger dans le salon & Dormir sans être dérangé par les notifications \\
\hline
2 & Rester assis 3 h d'affilée devant l'ordinateur & Faire une pause de 10 min chaque heure & Reposer les yeux, le dos et la nuque \\
\hline
3 & Utiliser le même mot de passe partout & Créer un mot de passe différent pour chaque compte & Si un compte est piraté, les autres restent protégés \\
\hline
4 & Regarder des vidéos jusqu'à minuit & Éteindre les écrans 30 min avant de dormir & La lumière bleue retarde l'endormissement \\
\hline
\end{tabularx}
\end{center}
\end{corrige}

\begin{corrige}[Exercice 8 — Réagir face au cyberharcèlement]
\begin{enumerate}
\item Il s'agit de cyberharcèlement car Nadège subit, à travers un outil numérique (groupe de discussion), des moqueries et des humiliations répétées : le montage photo est diffusé publiquement et plusieurs camarades ajoutent des commentaires blessants. L'intention de nuire, la répétition et l'effet de groupe caractérisent le harcèlement.
\item \textbf{Conséquences possibles pour Nadège :} tristesse, angoisse et perte de confiance en elle ; isolement et refus d'aller au collège ; baisse des résultats scolaires ; atteinte durable à son e-réputation (le montage peut circuler longtemps).
\item \textbf{Trois attitudes responsables de témoin :}
\begin{itemize}
\item ne pas rire, ne pas commenter et ne pas partager le montage (ne pas devenir complice) ;
\item soutenir Nadège, lui parler et l'encourager à ne pas rester seule ;
\item signaler les faits à un adulte (parent, enseignant, CPE) et conserver les captures d'écran comme preuves.
\end{itemize}
\end{enumerate}
\end{corrige}

\begin{corrige}[Exercice 9 — Type BEPC : le club informatique sensibilise]
\begin{enumerate}
\item L'\textbf{hygiène du numérique} est l'ensemble des règles et bonnes habitudes d'utilisation des outils numériques qui protègent la santé physique, l'équilibre psychologique et la vie sociale de l'utilisateur.
\item Temps hebdomadaire : $5 \times 7 = 35$ h par semaine.\\
Temps après 22 h sur une semaine : $2 \times 7 = 14$ h.\\
Un élève passe donc \textbf{35 h par semaine} devant un écran, dont \textbf{14 h après 22 h} (soit 40 \% du temps total, très tard le soir).
\item \textbf{Risques physiques :} fatigue et douleur des yeux ; douleurs du cou et du dos (mauvaise posture) ; troubles du sommeil (on en cite deux). \textbf{Risques psychologiques :} dépendance aux écrans ; stress, anxiété, isolement ; baisse de la concentration.
\item Quatre règles à afficher (avec justification) :
\begin{itemize}
\item \emph{Fais une pause de 10 minutes toutes les heures} : pour reposer les yeux et le dos.
\item \emph{Règle 20-20-20} : toutes les 20 minutes, regarder 20 secondes à 6 mètres pour éviter la fatigue visuelle.
\item \emph{Éteins les écrans au moins 30 minutes avant de dormir} : la lumière bleue perturbe l'endormissement.
\item \emph{Assieds-toi le dos droit, écran à hauteur des yeux, à 50--70 cm} : pour éviter les douleurs du cou et du dos.
\end{itemize}
\item Exemple de slogan : « \emph{Un écran maîtrisé, une santé préservée !} » ou « \emph{Déconnecte le soir, réussis le matin !} ».
\end{enumerate}
\textbf{Barème indicatif (sur 10) :} définition 1 pt ; calculs 2 pts (1 pt chacun) ; risques 2 pts (0,5 pt par risque) ; règles justifiées 4 pts (1 pt par règle justifiée) ; slogan 1 pt.\\
\textbf{Points de vigilance :} les règles doivent être \emph{justifiées} (une règle sans justification ne vaut que la moitié des points) ; les calculs doivent montrer l'opération posée ; le slogan doit être court et lié au thème santé/numérique.
\end{corrige}

\begin{corrige}[Exercice 10 — Type BEPC : les réseaux sociaux au collège]
\begin{enumerate}
\item \textbf{Texte conseil à Kamdem (10 à 15 lignes) :}\\
Cher Kamdem, un réseau social est un outil formidable si tu l'utilises de façon responsable. D'abord, respecte la \emph{netiquette}, c'est-à-dire les règles de politesse en ligne : salue, exprime-toi calmement, n'insulte jamais et n'écris pas en majuscules (cela équivaut à crier). Ensuite, protège ta \emph{vie privée} : ne publie jamais ton adresse, ton numéro de téléphone ou le nom de ton collège, choisis un mot de passe solide et secret, et règle ton compte en mode privé. Pense aussi à ton \emph{e-réputation} : tout ce que tu publies laisse une trace durable ; avant de poster une photo ou un commentaire, demande-toi si tu serais fier de le montrer à tes parents ou à un futur employeur. Enfin, si tu es témoin ou victime de \emph{cyberharcèlement} (moqueries, menaces, montages humiliants), ne réponds pas aux provocations, garde les preuves (captures d'écran) et signale les faits à un adulte. En suivant ces règles, tu profiteras d'Internet en toute sécurité.
\item \textbf{Démarche avant de publier la photo de la cantine :} Kamdem doit d'abord \textbf{demander l'autorisation} de chaque camarade visible sur la photo (et idéalement l'avis d'un adulte responsable). \emph{Justification :} chaque personne a un droit à l'image ; publier sans accord viole la vie privée des camarades, peut les mettre mal à l'aise et nuire à leur e-réputation si la photo circule.
\item \textbf{Face à l'inconnu qui demande sa localisation et son numéro :} Kamdem doit (1) \textbf{ne jamais répondre ni donner ces informations personnelles} ; (2) \textbf{bloquer l'inconnu et en parler immédiatement à un adulte} (parent ou enseignant), en gardant les messages comme preuves.
\end{enumerate}
\textbf{Barème indicatif (sur 10) :} texte 6 pts (1,5 pt par notion obligatoire correctement expliquée : netiquette, vie privée, e-réputation, cyberharcèlement) ; démarche photo 2 pts ; réaction face à l'inconnu 2 pts.\\
\textbf{Points de vigilance :} les quatre notions doivent être \emph{définies ou clairement expliquées}, pas seulement citées ; la justification de la question 2 doit mentionner le droit à l'image ; à la question 3, « ne pas répondre » seul ne suffit pas : il faut aussi signaler à un adulte.
\end{corrige}

\begin{corrige}[Exercice 11 — Type BEPC : enquête dans une classe de 3ème]
\begin{enumerate}
\item Pourcentages (classe de 50 élèves) :
\begin{align*}
\text{Téléphone dans la chambre : } & \frac{35}{50} \times 100 = 70\ \% \\
\text{Douleurs cou/yeux : } & \frac{20}{50} \times 100 = 40\ \% \\
\text{Messages moqueurs ou menaçants : } & \frac{10}{50} \times 100 = 20\ \%
\end{align*}
\item Risques et bonnes pratiques :
\begin{center}
\begin{tabularx}{\linewidth}{|X|X|X|}
\hline
\rowcolor{popblueL} \textbf{Situation} & \textbf{Risque encouru} & \textbf{Bonne pratique} \\
\hline
Téléphone dans la chambre (70 \%) & Sommeil perturbé par les notifications et la lumière bleue, fatigue & Éteindre le téléphone ou le laisser hors de la chambre la nuit \\
\hline
Douleurs au cou ou aux yeux (40 \%) & Fatigue visuelle, troubles de la posture & Pauses régulières, règle 20-20-20, écran à hauteur des yeux à 50--70 cm \\
\hline
Messages moqueurs/menaçants (20 \%) & Cyberharcèlement : angoisse, isolement, baisse des notes & Ne pas répondre, garder les preuves, signaler à un adulte \\
\hline
\end{tabularx}
\end{center}
\item \textbf{Plan d'hygiène numérique de la classe :}
\begin{center}
\begin{tabularx}{\linewidth}{|X|X|X|}
\hline
\rowcolor{popblueL} \textbf{Habitude à changer} & \textbf{Action concrète} & \textbf{Justification} \\
\hline
Dormir avec le téléphone & Le laisser dans le salon pour la nuit & Retrouver un sommeil complet et réparateur \\
\hline
Longues séances d'écran sans pause & Pause de 10 min toutes les heures & Préserver les yeux, le cou et le dos \\
\hline
Écrans tard le soir & Éteindre les écrans 30 min avant de dormir & La lumière bleue retarde l'endormissement \\
\hline
Silence face aux moqueries en ligne & Signaler tout message blessant à un adulte & Protéger les victimes et stopper le cyberharcèlement \\
\hline
\end{tabularx}
\end{center}
\item \textbf{Conclusion (5 à 8 lignes) :} Le bien-être social d'un élève dépend aussi de son comportement numérique, car une grande partie de sa vie relationnelle se déroule désormais en ligne. Un élève qui respecte la netiquette entretient de bonnes relations avec ses camarades, tandis que celui qui publie sans réfléchir peut blesser les autres ou abîmer sa propre e-réputation. Le cyberharcèlement, subi ou commis, détruit la confiance et l'amitié au sein de la classe. À l'inverse, un usage maîtrisé des écrans (sommeil suffisant, pauses, respect d'autrui) rend l'élève plus disponible, plus serein et mieux intégré dans son groupe. Bien vivre avec les autres, aujourd'hui, c'est donc aussi bien vivre avec le numérique.
\end{enumerate}
\textbf{Barème indicatif (sur 10) :} pourcentages 1,5 pt (0,5 pt chacun) ; risques et bonnes pratiques 3 pts (1 pt par ligne) ; plan en tableau 3 pts (0,75 pt par ligne cohérente) ; conclusion 2,5 pts.\\
\textbf{Points de vigilance :} les pourcentages doivent être calculés à partir de l'effectif total (50) avec la formule visible ; chaque bonne pratique doit correspondre au risque identifié ; la conclusion doit faire le lien entre \emph{comportement numérique} et \emph{relations sociales}, pas seulement parler de santé physique.
\end{corrige}

\newpage

\coursec{Fiche bilan}

\begin{popfigure}
\begin{center}\tcbox[colback=popblue,colframe=popblue,coltext=white,arc=9pt,boxsep=4pt,left=6pt,right=6pt]{\bfseries\small Hygiène du numérique}\end{center}
\vspace{-2pt}
\begin{tcbraster}[raster columns=3,raster equal height=rows,raster column skip=6pt,raster row skip=6pt,enhanced,blankest]
\begin{tcolorbox}[enhanced,colback=poporange,colframe=poporange,coltext=white,boxrule=0.9pt,arc=3pt,left=4pt,right=4pt,top=3pt,bottom=3pt,boxsep=1pt,halign=left]\bfseries\scriptsize Notion : bien utiliser les outils numériques\end{tcolorbox}
\begin{tcolorbox}[enhanced,colback=poppurple,colframe=poppurple,coltext=white,boxrule=0.9pt,arc=3pt,left=4pt,right=4pt,top=3pt,bottom=3pt,boxsep=1pt,halign=left]\bfseries\scriptsize Risques physiques : yeux, dos, sommeil\end{tcolorbox}
\begin{tcolorbox}[enhanced,colback=poppink,colframe=poppink,coltext=white,boxrule=0.9pt,arc=3pt,left=4pt,right=4pt,top=3pt,bottom=3pt,boxsep=1pt,halign=left]\bfseries\scriptsize Risques psycho-sociaux : addiction, cyberharcèlement\end{tcolorbox}
\begin{tcolorbox}[enhanced,colback=popgreen,colframe=popgreen,coltext=white,boxrule=0.9pt,arc=3pt,left=4pt,right=4pt,top=3pt,bottom=3pt,boxsep=1pt,halign=left]\bfseries\scriptsize Bonnes pratiques : pauses, posture, luminosité\end{tcolorbox}
\begin{tcolorbox}[enhanced,colback=popgold,colframe=popgold,coltext=white,boxrule=0.9pt,arc=3pt,left=4pt,right=4pt,top=3pt,bottom=3pt,boxsep=1pt,halign=left]\bfseries\scriptsize Réseaux sociaux : usage responsable\end{tcolorbox}
\begin{tcolorbox}[enhanced,colback=popblue!70,colframe=popblue!70,coltext=white,boxrule=0.9pt,arc=3pt,left=4pt,right=4pt,top=3pt,bottom=3pt,boxsep=1pt,halign=left]\bfseries\scriptsize Démarche : observer, agir, justifier\end{tcolorbox}
\end{tcbraster}

\legende{Carte mentale de la leçon : l'hygiène du numérique et le bien-être social}
\end{popfigure}

\begin{bilan}
\textbf{Définitions clés}
\begin{itemize}
  \item \cle{Hygiène du numérique} : ensemble des bonnes habitudes et règles d'usage des outils numériques (ordinateur, téléphone, tablette) pour préserver sa \textbf{santé physique}, son \textbf{équilibre psychologique} et sa \textbf{vie sociale}.
  \item \cle{Fatigue visuelle} : trouble des yeux (sécheresse, maux de tête, vision floue) provoqué par une exposition prolongée aux écrans.
  \item \cle{Addiction numérique} : dépendance excessive aux écrans, aux jeux vidéo ou aux réseaux sociaux, qui nuit au sommeil, aux études et aux relations.
  \item \cle{Cyberharcèlement} : agressions répétées (moqueries, menaces, rumeurs) commises par Internet ou téléphone.
\end{itemize}

\textbf{Les risques à connaître par c\oe ur}
\begin{center}
\begin{tabularx}{\linewidth}{|l|X|}
\hline
\rowcolor{popblueL} \textbf{Type de risque} & \textbf{Exemples} \\
\hline
Physique & fatigue visuelle, maux de dos et de nuque, troubles du sommeil (lumière bleue), sédentarité \\
\hline
Psychologique & stress, anxiété, addiction, isolement, baisse de concentration \\
\hline
Social & cyberharcèlement, repli sur soi, conflits, mauvaise image de soi \\
\hline
\end{tabularx}
\end{center}

\textbf{Méthodes express}
\begin{itemize}
  \item \textbf{Règle 20-20-20} : toutes les 20 minutes d'écran, regarder au loin (environ 6 m) pendant 20 secondes.
  \item \textbf{Bonne posture} : écran à hauteur des yeux, dos droit, pieds au sol, pièce éclairée.
  \item \textbf{Limiter le temps d'écran} : pas d'écran au moins 1 heure avant de dormir ; fixer des horaires.
  \item \textbf{Réseaux sociaux} : respecter les autres, ne pas tout partager, bloquer et signaler les harceleurs, parler à un adulte en cas de problème.
  \item \textbf{Rédiger une démarche} : observer la situation $\rightarrow$ identifier le risque $\rightarrow$ proposer une solution $\rightarrow$ justifier le choix.
\end{itemize}
\end{bilan}

\begin{aretenir}[Checklist avant le BEPC]
\begin{itemize}
  \item[$\square$] Je sais définir l'hygiène du numérique avec mes propres mots.
  \item[$\square$] Je cite au moins trois risques physiques liés aux écrans.
  \item[$\square$] Je cite au moins deux risques psychologiques et deux risques sociaux.
  \item[$\square$] Je connais la règle 20-20-20 pour protéger mes yeux.
  \item[$\square$] Je sais décrire la bonne posture devant un ordinateur.
  \item[$\square$] Je sais pourquoi il faut éviter les écrans avant de dormir.
  \item[$\square$] Je connais les bons comportements sur les réseaux sociaux.
  \item[$\square$] Je sais quoi faire en cas de cyberharcèlement (bloquer, signaler, en parler).
  \item[$\square$] Je sais rédiger et justifier une démarche d'hygiène numérique en quatre étapes.
  \item[$\square$] Je sais appliquer ces notions à une situation concrète de ma vie quotidienne.
\end{itemize}
\end{aretenir}

\findecours

\end{document}
